%~Mouliné par MaN_auto v.0.36.1 (7beb38e9) 2025-09-22 10:15:29
\documentclass[AHL,Unicode,longabstracts]{cedram}

\usepackage{zref-clever}
\usepackage{mathtools}
%\usepackage{stmaryrd}
\usepackage{wasysym}
\usepackage{mathrsfs}
\usepackage{tikz}
%\usepackage{enumitem}
%\usepackage{subfloat}
%\usepackage{mleftright}
%\usepackage{caption}
\usepackage[labelformat=simple]{subcaption}
\usepackage{aligned-overset}
%\usepackage[caption=false]{subfig}

%%%%%%%%%%%%%%%%%%%%%%%%%%%%%%%%%%%%%%%%%%%%%%%%%%%%%%%%%%%%%%%%%%%%%
\renewcommand\thesubfigure{(\alph{subfigure})}
%\renewcommand*\thesubfloatfigure{\themainfigure(\alph{subfloatfigure})}
\zcsetup{ abbrev=false, cap=true, nameinlink=false, rangetopair=false, endrange=stripprefix }

\zcRefTypeSetup{equation}{ Name-sg =, name-sg =, Name-pl =, name-pl =, refbounds = {(,,,)}, +refbounds-rb = {(,,,}, +refbounds-re = {,,,)}, rangesep = {--} }
\zcsetup{countertype={cdrthm=theorem}}
\AddToHook{env/defi/begin}{\zcsetup{countertype={cdrthm=definition}}}
\AddToHook{env/prop/begin}{\zcsetup{countertype={cdrthm=proposition}}}
\AddToHook{env/coro/begin}{\zcsetup{countertype={cdrthm=corollary}}}
\AddToHook{env/lemm/begin}{\zcsetup{countertype={cdrthm=lemma}}}
\AddToHook{env/rema/begin}{\zcsetup{countertype={cdrthm=remark}}}
%%%%%%%%%%%%%%%%%%%%%%%%%%%%%%%%%%%%%%%%%%%%%%%%%%%%%%%%%%%%%%%%%%%%%
%\makeatletter
%\let\listoffiguresfontsize\footnotesize
%\makeatother




\let\oversetx\overset
%\mleftright

\newcounter{setofcase}
\newcounter{case}[setofcase]
%\zcrefname{case}{Case}{Cases}
\newenvironment{case}[1][]{\refstepcounter{case}\begin{proof}[{Case~\thecase\if#1\empty\else: #1\fi}]}{\let\qed\relax\end{proof}}
\newcommand{\newcaseset}{\stepcounter{setofcase}}


\newcounter{setofstep}
\newcounter{step}[setofstep]
%\crefname{step}{Step}{Steps}
\newenvironment{step}[1][]{\refstepcounter{step}\begin{proof}[{Step~\thestep\if#1\empty\else: #1\fi}]}{\let\qed\relax\end{proof}}
\newcommand{\newstepset}{\stepcounter{setofstep}}



\newcommand*\ssS{\mathsf{S}}
\newcommand*\GS{\mathsf{GS}}
\newcommand*\ssN{\mathsf{N}}
\newcommand*\sfs{\mathsf{s}}
\newcommand*\sfr{\mathsf{r}}
\newcommand*\sff{\mathsf{f}}
\newcommand*\sfT{\mathsf{T}}


\newcommand*\scM{\mathscr{M}}
\newcommand*\scN{\mathscr{N}}

\newcommand*\TT{\mathbb{T}}
\newcommand*\RR{\mathbb{R}}
\newcommand*\ZZ{\mathbb{Z}}
\newcommand*\NN{\mathbb{N}}
\newcommand{\bbQ}{\mathbb{Q}}
\newcommand{\bbP}{\mathbb{P}}
\newcommand{\bbU}{\mathbb{U}}
\newcommand{\bbE}{\mathbb{E}}

\newcommand{\bfun}{\mathbf{1}}
\newcommand{\clA}{\mathcal{A}}
\newcommand{\clJ}{\mathcal{J}}
\newcommand{\clR}{\mathcal{R}}
\newcommand{\clE}{\mathcal{E}}
\newcommand{\clF}{\mathcal{F}}
\newcommand{\clT}{\mathcal{T}}



\newcommand*\st{\mathrel{:}}

\DeclareMathOperator{\supp}{supp}
\DeclareMathOperator{\locc}{loc}
\newcommand*\e{\mathrm{e}}
\newcommand{\rms}{\mathrm{s}}
\newcommand*\dif{\mathrm{d}}
\newcommand*\Left{\mathrm{L}}
\newcommand*\Right{\mathrm{R}}
\newcommand*\LR{\mathrm{LR}}

\newcommand*\eps{\varepsilon}
%%%%%%%%%%%%%%%%%%%%%%%%%%%%%%%%%%%%%%%%%%%%%%%%%%%%%%%%%%%%%%%

\renewcommand*{\to}{\mathchoice{\longrightarrow}{\rightarrow}{\rightarrow}{\rightarrow}}

\let\oldmapsto\mapsto
\renewcommand*{\mapsto}{\mathchoice{\longmapsto}{\oldmapsto}{\oldmapsto}{\oldmapsto}}

\let\oldin\in
\renewcommand{\in}{\mathchoice{\oldin}{\oldin}{\,\oldin\,}{\,\oldin\,}}
\let\oldcap\cap
\renewcommand{\cap}{\mathchoice{\oldcap}{\oldcap}{\,\oldcap\,}{\,\oldcap\,}}

\let\originalleft\left 
\let\originalright\right
\renewcommand{\left}{\mathopen{}\mathclose\bgroup\originalleft}
\renewcommand{\right}{\aftergroup\egroup\originalright}
%%%%%%%%%%%%%%%%%%%%%%%%%%%%%%%%%%%%%%%%%%%%%%%%%%%%%%%%%%%%%%%%%%%%%%%%%%%%%%%%%%%%%%%%%%%
\graphicspath{{./figures/}}

\newcommand*{\mk}{\mkern -1mu}
\newcommand*{\Mk}{\mkern -2mu}
\newcommand*{\mK}{\mkern 1mu}
\newcommand*{\MK}{\mkern 2mu}

\hypersetup{urlcolor=purple, linkcolor=blue, citecolor=red}

\newcommand*{\relabel}{\renewcommand{\labelenumi}{(\theenumi)}}
\newcommand*{\romanenumi}{\renewcommand*{\theenumi}{\roman{enumi}}\relabel}
\newcommand*{\Romanenumi}{\renewcommand*{\theenumi}{\Roman{enumi}}\relabel}
\newcommand*{\alphenumi}{\renewcommand*{\theenumi}{\alph{enumi}}\relabel}
\newcommand*{\Alphenumi}{\renewcommand*{\theenumi}{\Alph{enumi}}\relabel}
\let\oldtilde\tilde
\renewcommand*{\tilde}[1]{\mathchoice{\widetilde{#1}}{\widetilde{#1}}{\oldtilde{#1}}{\oldtilde{#1}}}
\let\oldhat\hat
\renewcommand*{\hat}[1]{\mathchoice{\widehat{#1}}{\widehat{#1}}{\oldhat{#1}}{\oldhat{#1}}}
\let\oldexists\exists
\renewcommand*{\exists}{\mathrel{\oldexists}}

\newcommand{\llbracket}{\mathopen{[\mkern -2.7mu[}}
\newcommand{\rrbracket}{\mathclose{]\mkern -2.7mu]}}
\newcommand{\llb}{\llbracket}
\newcommand{\rrb}{\rrbracket}
%%%%%%%%%%%%%%%%%%%%%%%%%%%%%%%%%%%%%%%%%%%%%%%%%%%%%%%%%%%%%%%%%%%%%%%%%%%%%%%%%%%%%%%%%%%

\title[Simultaneous global inviscid periodic Burgers flows]{Simultaneous global inviscid Burgers flows with periodic Poisson forcing}

\alttitle{Flots de Burgers inviscides globaux simultanés avec forçage de Poisson périodique}

\subjclass{60K37, 37L55, 60H15, 60G55}
\keywords{inviscid Burgers equation, global solutions, invariant measures}


\author[\initial{A.} \lastname{Dunlap}]{\firstname{Alexander} \lastname{Dunlap}}
\address{Department of Mathematics,\\
Duke University,\\
120~Science~Dr,\\
Durham, NC, 27708 (USA)}
\email{alexander.dunlap@duke.edu}



\begin{abstract}
We study the inviscid Burgers equation on the circle $\TT:= \RR/\ZZ$ forced by the spatial derivative of a Poisson point process on $ \RR\times\TT$. We construct global solutions with mean $\theta$ simultaneously for all $\theta\in\RR$, and in addition construct their associated global shocks (which are unique except on a countable set of $\theta$). We then show that as $\theta$ changes, the solution only changes through the movement of the global shock, and give precise formulas for this movement. This can be seen as an analogue of previous results by the author and Yu Gu in the viscous case with white-in-time forcing, which related the derivative of the solution in $\theta$ to the density of a particle diffusing in the Burgers flow.
\end{abstract}

\begin{altabstract}
Nous étudions l'équation de Burgers inviscide sur le cercle $\TT := \RR/\ZZ$ forcée par la dérivée spatiale d'un processus ponctuel de Poisson sur $ \RR\times\TT$. Nous construisons des solutions globales avec une moyenne $\theta$ simultanément pour tout $\theta\in\RR$, et nous construisons en outre les chocs globaux associés (qui sont uniques sauf sur un ensemble dénombrable de $\theta$). Nous montrons ensuite que lorsque $\theta$ change, la solution ne change qu'à travers le mouvement du choc global, et nous donnons des formules précises pour ce mouvement. Ceci peut être considéré comme un analogue de résultats précédents par l'auteur et Yu Gu dans le cas visqueux avec un forçage blanc dans le temps, qui reliait la dérivée de la solution en $\theta$ à la densité d'une particule diffusant dans le flot de Burgers.
\end{altabstract}


\datereceived{2024-06-16}
\daterevised{2025-03-03}
\dateaccepted{2025-06-05}

\editors{S.~Gou\"ezel and H.~Lacoin}

\begin{DefTralics}
\newcommand*\TT{\mathbb{T}}
\newcommand*\RR{\mathbb{R}}
\newcommand*\ZZ{\mathbb{Z}}
\end{DefTralics}
%%%%%%%%%%%%%%%%%%%%%%%%%%%%%%%%%%%%%%%%%%%%%%%%%%%%%%%%%%%%%%%%%%%%%%%%%%%%%%%%%%%%%%%%%%%%%
\begin{document}
\maketitle



\tableofcontents{}
\listoffigures{\phantom{scriptsize}}

An animation is also included in the supplementary material\footnote{Available online at~\url{
https://annales.lebesgue.fr/public/assets/751/moving\_global\_shock.mp4}.} and described in Section~\ref{subsec:illustrations}.

\section{Introduction}

\subsection{Background}

Let $\TT\coloneqq\RR/\ZZ$ and let $\pi\colon\RR\to\TT$ be the projection map. Let $\mu$ be a purely atomic measure on $ \RR\times\TT$ such that $\ssN\coloneqq\supp\mu$ is a discrete set. We will think of $\mu$ as a random measure, and our main example will be when $\mu$ is a realization of a homogeneous Poisson point process. We are interested in global solutions to the forced inviscid Burgers equation given by
\begin{subequations}\label{eq:uproblem}
\begin{align}
\partial_{t}u_{\theta}(t,x)+\partial_{x}\left(\frac{1}{2}u_{\theta}^{2}+\mu\right)(t,x) & =0,\qquad\theta,t\in\RR,x\in\TT;\label{eq:ueqn}
\\
\int_{\TT}u_{\theta}(t,x)\,\dif x & =\theta,\qquad\theta,t\in\RR.\label{eq:uintegral}
\end{align}
\end{subequations}
This type of Burgers equation with discrete forcing has been considered in the whole-line setting, and briefly in the periodic setting, in~\cite{Bak13,BCK14}. As is usual for the inviscid Burgers equation, we make sense of the problem \eqref{eq:uproblem} via entropy solutions, which we define via Lagrangian minimizers. For $X\in H^{1}([s,t];\TT)$, and $\theta\in\RR$, we define the action
\begin{equation}\label{eq:Adef}
\clA_{\theta,s,t}[X]\coloneqq\frac{1}{2}\int_{s}^{t}(X'(r)-\theta)^{2}\,\dif r-\mu\bigl(\{(r,X(r))\st r\in[s,t)\}\bigr),
\end{equation}
where $X'$ denotes the derivative of $X$. Roughly speaking, entropy solutions to \eqref{eq:uproblem} with initial data $u_{\theta}(s,x)=\theta+\partial_{x}G(x)$ are given by $u_{\theta}(t,x)=X'(t)$, where $X$ minimizes $\clA_{\theta,s,t}[X]+G(X(s))$ over all paths satisfying $X(t)=x$. (See Definition~\ref{def:minimizers} below.) At some points, the slope of the minimizer is not unique, and at these ``shock'' points, the solution $u$ is discontinuous in space. This minimization problem resembles that defining the Hammersley process~\cite{AD95,Ham72}, but here we impose a quadratic penalty rather than a hard cutoff on the slope of the minimizers.

The integral on the left side of \zcref{eq:uintegral} is preserved by the Burgers evolution, as can be seen formally by integrating \zcref{eq:ueqn} in space. In terms of the minimization problem \zcref{eq:Adef}, the role of $\theta$ is to encourage potential minimizers to have average slope $\theta$. Previous works in the mathematics literature (see e.g.~\cite{Bak13,EKMS00,GIKP05,IK03}) on the forced Burgers equation in the periodic setting have mostly considered the case $\theta=0$. This is essentially equivalent to considering the problem for any fixed value of $\theta$, since if $u$ solves \zcref{eq:ueqn}, then $\tilde{u}(t,x)=u(t,x-\theta t)+\theta$ solves
\begin{equation}\label{eq:veqn}
\partial_{t}\tilde{u}(t,x)+\frac{1}{2}\partial_{x}\tilde{u}^{2}(t,x)+\partial_{x}\mu(t,x-\theta t)=0.
\end{equation}
If the forcing $\mu$ is taken to be random and shear-invariant in law, then the solution maps of the problems \zcref{eq:veqn} and \zcref{eq:ueqn} have the same probability distributions.

However, equivalence in law of the solutions with different values of $\theta$ does not tell us about the behavior of the equation solved for all values of $\theta$ \emph{simultaneously} (i.e.~with the same realization of the forcing $\mu$). This has been a topic of significant recent interest for various models in the KPZ universality class on the whole line; see e.g.~\cite{BFS23,BSS23,BSS24,FS20,GRASS23}. In the Burgers setting on the torus, a time-periodic version of the problem has been discussed at a physics level in~\cite{BK07}. Important questions in the multi-$\theta$ context include (1) how to construct jointly invariant measures for the randomly forced Burgers equation simultaneously for all values of $\theta$; and (2) how solutions sampled from these invariant measures change as the parameter $\theta$ changes. On the whole line, it has been found in all examples that have been considered that the spatial integral of the process (i.e.~the process of spatial increments of the KPZ equation) almost surely exhibits a countable dense set of discontinuities in $\theta$. However, the proofs of this property in each case are relying on some exact computation that can be performed for the model, and the physical phenomenology behind these discontinuities remains to be fully understood.

The author and Yu Gu have studied a related problem for the \emph{viscous} Burgers equation on the torus, forced by a white-in-time Gaussian noise, in~\cite{DG23}. The~equation we considered in that setting is
\begin{subequations}\label{eq:vproblem}
\begin{align}
\dif v_{\theta}(t,x) & =\frac{1}{2}\left[\partial_{x}^{2}v_{\theta}+\partial_{x}(v_{\theta}^{2})\right](t,x)\dif t+\dif V(t,x), && \theta,t\in\RR,x\in\TT;\label{eq:veqn-1}
\\
\int_{\TT}v_{\theta}(t,x)\,\dif x & =\theta, &&\theta,t\in\RR,\label{eq:vintegral}
\end{align}
\end{subequations}
where $\dif V$ is the noise. We showed that, at stationarity, we have
\begin{equation}\label{eq:viscouscase}
u_{\theta_{2}}(t,x)-u_{\theta_{1}}(t,x)=\int_{\theta_{1}}^{\theta_{2}}g_{\theta}(t,x)\,\dif\theta,
\end{equation}
where $g_{\theta}$ is a statistically stationary solution to the associated coupled PDE
\begin{subequations}\label{eq:gproblem}
\begin{align}
\partial_{t}g_{\theta}(t,x) & =\frac{1}{2}\Delta g_{\theta}(t,x)+\partial_{x}(u_{\theta}g_{\theta})(t,x), &&\theta,t\in\RR,x\in\TT;\label{eq:geqn}
\\
\int_{\TT}g_{\theta}(t,x)\,\dif x & =1, &&\theta,t\in\RR.\label{eq:gintegral}
\end{align}
\end{subequations}
The problem \eqref{eq:gproblem} can be obtained by differentiating
\eqref{eq:vproblem} in $\theta$ and setting $g_{\theta}=\partial_{\theta}u_{\theta}$. On the other hand, from the form of \eqref{eq:geqn}, we see that $g_{\theta}(t,\cdot)$ is the density of a particle diffusing in the flow given by $u_{\theta}$, with unit diffusivity. A particular consequence of \eqref{eq:viscouscase} (along with a moment bound on $g_{\theta}$ proved in~\cite{DG23}) is that $u_{\theta}(t,x)$ is continuous in $\theta$, which is in sharp contrast to the behavior that has been observed on the real line.

The purpose of the present paper is to study the jointly invariant measures of the \emph{inviscid} periodic Burgers equation \eqref{eq:uproblem} simultaneously for all $\theta$. The main results we will prove can be seen as analogues of those in~\cite{DG23}. In particular, we will show that jointly stationary solutions for \eqref{eq:uproblem} exist, that a one-force-one-solution principle holds, and that an inviscid analogue of the relation \eqref{eq:viscouscase} holds. It will no longer be the case that $u_{\theta}(t,x)$ is continuous in $\theta$ because the analogue of $g_{\theta}$ is less regular in our setting. However, the spatial integral of $g_{\theta}$ will be continuous in $x$, and hence $\int_{0}^{x}u_{\theta}(t,y)\,\dif y$ will still be continuous in $\theta$. Actually, this is a rather generic feature of the periodic case: if we define $h_{\theta}(t,x)=\int_{0}^{x}u_{\theta}(t,x)\,\dif x$, then by the comparison principle we also have for $\theta_{2}\ge\theta_{1}$ that $u_{\theta_{2}}\ge u_{\theta_{1}}$ (assuming a similar ordering for the initial condition), and so for $x\in[0,1]$ we have
\begin{equation}\label{eq:hsqueeze}
0\le h_{\theta_{2}}(t,x)-h_{\theta_{1}}(t,x)\le h_{\theta_{2}}(t,1)-h_{\theta_{1}}(t,1)\overset{\eqref{eq:uintegral}}{=}\theta_{2}-\theta_{1}.
\end{equation}

In the inviscid setting we consider here, the particle diffusing in the Burgers flow should be replaced by a particle simply moving in the flow, without diffusion. It is well-known that such a particle will eventually join a \emph{shock }of the Burgers flow. This suggests that, at least formally, the density $g_{\theta}$ should be replaced by a delta mass at a single ``global'' shock, and \zcref{eq:viscouscase} then suggests that the change in $u$ as $\theta$ is varied should occur only at the location of this global shock. We will prove a precise version of this statement in our main theorem Theorem~\ref{thm:mainthm} below. We also point the reader to the video included the supplementary material for a visualization of the shocks in the Burgers flow as $\theta$ is changed; see Section~\ref{subsec:illustrations} for a description.

The proof techniques in the present setting are entirely different from those of~\cite{DG23}. In particular, we study the minimizers of the functional \zcref{eq:Adef}, and the associated shocks as mentioned above, rather than using the stochastic analysis tools of~\cite{DG23}. This leads us to a fine study of the structure of one-sided minimizers and global shocks for the inviscid Burgers equation. One-sided minimizers and global shocks have been studied extensively in the literature on the stochastic Burgers equation. We refer to the survey~\cite{BK18} for an illuminating heuristic discussion. The use of minimizers rather than polymers makes many features of the problem more explicit than in the viscous case, and certain aspects of the phenomenology are clearer. In particular, we will see how the topology of the minimizers plays an important role in the analysis. The importance of the topology of the minimizers has been previously observed in the physics literature in~\cite{BIK02,BK07} for the Burgers equation with time-periodic forcing that is smoother in time than we consider here.

A consistent theme of work on stochastic Burgers with multiple means considered simultaneously is the presence of exceptional values of $\theta$ at which behavior is observed that happens with probability zero for any fixed $\theta$. This holds as well in our setting. Thus, our study of minimizers and shocks will go beyond that of previous work in that we will prove the behavior at these exceptional values of $\theta$ as well. In particular, an important feature of the study of minimizers for the Burgers equation is that distinct minimizers from the same point do not typically cross each other. At exceptional values of $\theta$, a certain amount of crossing is possible, and thus the picture that we describe exhibits significant additional topological complexity compared to the fixed-$\theta$ case.

\subsection{Main results}

We now state precisely the main results of our study. Although we are primarily interested in the setting in which $\mu$ is a Poisson point process, we can actually state simple deterministic conditions on $\mu$ under which our results hold. We will then check (see Theorem~\ref{thm:poisson-ok}) that these conditions are satisfied with probability $1$ by any homogeneous compound Poisson point process. To facilitate this, we introduce the following space of forcing measures.

\begin{defi}
Let $\Omega$ be the space of purely atomic measures $\mu$ on $\RR\times\TT$ such that $\ssN\coloneqq\supp\mu$ is a discrete set.
\end{defi}

Since much of our work concerns the behavior of Lagrangian minimizers, we first define them precisely.

\begin{defi}[Lagrangian minimizers]\label{def:minimizers}
Let $\theta\in\RR$ and $s<t$.
\begin{enumerate}
\item We define the set $\scM_{s,y\mid t,x}^{\theta}$ comprising all paths $X\in H^{1}([s,t];\TT)$ with $X(s)=y$ and $X(t)=x$ such that if $Y$ is another such path, then $\clA_{\theta,s,t}[X]\le\clA_{\theta,s,t}[Y]$.
\item We define the set of \emph{one-sided minimizers}
\[
\scM_{t}^{\theta}\coloneqq\left\{X\in H_{\locc}^{1}((-\infty,t];\TT)\st X|_{[s,t]}\in\scM_{s,X(s)\mid t,X(t)}^{\theta}\text{ for all }s\le t\right\}.
\]
\item For $x\in\TT$, we define
\[
\scM_{t,x}^{\theta}\coloneqq\left\{X\in\scM_{t}^{\theta}\st X(t)=x\right\}.
\]
\item We define the partial order $\preccurlyeq$ on $\scM_{t,x}^{\theta}$ by
\begin{equation}\label{eq:partialorderdef}
X_{1}\preccurlyeq X_{2}\quad\text{if }\int_{r}^{t}X_{1}'(s)\,\dif s\ge\int_{r}^{t}X_{2}'(s)\,\dif s\quad\text{for all }r\le t.
\end{equation}
\end{enumerate}
\end{defi}

The partial order $\preccurlyeq$ represents the ordering of minimizers when lifted to the universal cover $\RR$ of $\RR/\ZZ$. The reason for the ``$\ge$'' sign in \zcref{eq:partialorderdef} is that we say that $X_{1}\preccurlyeq X_{2}$ if the graph of $X_{1}$ lies (not necessarily strictly) to the left of that of $X_{2}$ when lifted to the universal cover, which corresponds to the integral of the derivative of $X_{1}$ being at least that of $X_{2}$. This ordering will play an important role in our topological arguments.

Our first theorem will state that one-sided minimizers exist. To state it, we first need to introduce another assumption on the noise field:
\begin{defi}\label{def:smallnoisezones}
We define $\widetilde{\Omega}_{1}$ to be the set of all $\mu\in\Omega$ such that there exists an $M=M(\mu)>0$ such that for all $t\in\RR$, there is an $s=s(t)\le t-M$ and a $y\in\TT$ such that
\begin{equation}\label{eq:intersectiononly}
\ssN\cap([s-M,s+M]\times\TT)=\{(s,y)\}
\end{equation}
and
\begin{equation}\label{eq:youshouldtake}
\mu(\{(s,y)\})>\frac{1}{4M},
\end{equation}
and moreover that we can choose $s(t)$ in such a way that
\begin{equation}\label{eq:stttoinfty}
\lim_{t\to+\infty}s(t)=+\infty.
\end{equation}
\end{defi}

Definition~\ref{def:smallnoisezones} encodes the notion of \emph{small-noise zones}: regions of space-time containing only a single, sufficiently strong forcing point. All minimizers started sufficiently far in the future and extending sufficiently far into the past must pass through such points. Indeed, since there are no other forcing points in the vicinity,\pagebreak{} there is nothing to be lost by passing through the forcing point; see Proposition~\ref{prop:gothroughy} below. These small-noise zones act as regeneration times for the dynamics, as the behavior of the minimizers before and after them is decoupled. They have already been mentioned for this problem in~\cite{Bak13}, and the simple new observation here is that the decoupling happens simultaneously over all $\theta$.

Given the existence of small-noise zones, we can prove the existence of global solutions to~\eqref{eq:uproblem}:

\begin{theo}\label{thm:minimizers-exist}
Suppose that $\mu\in\widetilde{\Omega}_{1}$. For each $\theta,t\in\RR$ and $x\in\TT$, the set~$\scM_{t,x}^{\theta}$ is nonempty, consists of piecewise-linear paths connecting $(t,x)$ and points of $\ssN$, and has unique minimal and maximal elements $X_{\theta,t,x,\Left}$ and $X_{\theta,t,x,\Right}$, respectively, under the partial order $\preccurlyeq$.
\end{theo}

Again, we note that, in the case when $\mu$ is a Poisson point process,
Theorem~\ref{thm:minimizers-exist} has been proved for a single $\theta$ at a time in~\cite{Bak13}. The novelty here (even in the case when $\mu$ is a Poisson point process) is that we prove it for all $\theta$ on a single set~$\widetilde{\Omega}_{1}$ (which will have probability $1$ when $\mu$ is a Poisson point process). Theorem~\ref{thm:minimizers-exist} is proved as part of Proposition~\ref{prop:globalminimizersexist} in Section~\ref{sec:Existence}.

Let $\LR=\{\Left,\Right\}$. We define, for $\Square\in\LR$,
\begin{equation}\label{eq:udef}
u_{\theta,\Square}(t,x)=X_{\theta,t,x,\Square}'(t-),
\end{equation}
the derivative from below of the process $X_{\theta,t,x,\Square}$ at $t$. Both $u_{\theta,\Left}$ and $u_{\theta,\Right}$ are global-in-time entropy solutions to~\eqref{eq:uproblem}. It is clear from the definitions that $u_{\theta,\Left}(t,x)\ge u_{\theta,\Right}(t,x)$ for all $(\theta,t,x)\in\RR\times\RR\times\TT$. They differ only on the set
\begin{equation}\label{eq:Sdef}
\ssS\coloneqq\left\{(\theta,t,x)\in\RR\times\RR\times\TT\st u_{\theta,\Left}(t,x)>u_{\theta,\Right}(t,x)\right\},
\end{equation}
which is the set of \emph{shocks}. We define
\begin{equation}\label{eq:Sthetadef}
\ssS_{\theta}\coloneqq\{(t,x)\in\RR\times\RR\times\TT\st(\theta,t,x)\in\ssS\}
\end{equation}
and
\begin{equation}\label{eq:Sthetatdef}
\ssS_{\theta,t}\coloneqq\{x\in\TT\st(t,x)\in\ssS_{\theta}\}.
\end{equation}

Theorem~\ref{thm:minimizers-exist} describes a picture (previously observed in~\cite{Bak13}) in which, for each fixed $\theta$, the time-space cylinder $\RR\times\TT$ is tessellated by regions of points $(t,x)$ for which the last forcing point on $X_{\theta,t,x,\Square}$ is a given forcing point. The shock set $\ssS_{\theta}$ is formed from the boundaries of these regions (except that points of $\ssN$ are generally on the boundaries of these regions but not in $\ssS_{\theta}$). The dynamics of shocks are also well-understood. Each forcing point creates a pair of shocks starting at that point. (See Proposition~\ref{prop:advanceint-N}.) A shock at position $(t,x)$ moves (as time advances) with velocity $\frac{1}{2}(u_{\theta,\Left}(t,x)+u_{\theta,\Right}(t,x))$ (the well-known \emph{Rankine--Hugoniot condition}; see Proposition~\ref{prop:advanceint} for a proof in our setting). When two shocks collide with one another, they merge to form a single shock. See Figure~\ref{fig:shocks_picture} for an illustration.

\begin{figure}[!hbp]
\begin{centering}
\input{figures/shocks_picture.pgf}
\end{centering}
\caption[Forcing points, minimizers, shocks, and the Burgers solution.]{The bottom plot shows the forcing points (yellow), a sample of minimizers (black dashed lines), and the shocks (blue solid lines), which are points at which minimizers extend in multiple directions. The top plot shows $u_{\theta}(t,\cdot)=u_{\theta,\Square}(t,\cdot)$, where $\Square\in\protect\LR$ is arbitrary from the point of view of plotting. Note that the discontinuities of $u_{\theta}(t,\cdot)$ correspond to the locations of the blue shock curves at time $t$.}
\label{fig:shocks_picture}
\end{figure}

It has been observed in~\cite{EKMS00} for a fixed value of $\theta$ (with a different type of forcing) that, with probability $1$, there is a unique \emph{global minimizer} and a unique \emph{global shock} (the latter also known as the \emph{main shock} or \emph{topological shock}). All minimizers merge with the global minimizer as time goes to $-\infty$, and all shocks merge with the global shock as time goes to $+\infty$. (In the forcing considered in~\cite{EKMS00}, the~shocks converge towards each other exponentially fast rather than literally merging.) The term \emph{topological shock }is particularly illuminating; it refers to the fact that the global shock is characterized by the presence of minimizers that, at the point that they merge back together, have between them completed a nontrivial winding about the torus. See~\cite[Theorem~5.2]{EKMS00}.

Let us now state this topological characterization of global shocks precisely. For $(\theta,t,x)\in\ssS$, we define
\begin{equation}\label{eq:Tveedef}
T_{\vee}(\theta,t,x)\coloneqq\sup\left\{s<t\st X_{\theta,t,x,\Left}(s)=X_{\theta,t,x,\Right}(s)\right\}.
\end{equation}

\begin{defi}\label{def:globalshocks}
We define the set $\GS$ of \emph{global shocks} comprising all shocks $(\theta,t,x)\in\ssS$ such that
\begin{equation}\label{eq:shockslopesdiffernt}
\int_{T_{\vee}(\theta,t,x)}^{t}X_{\theta,t,x,\Left}'(s)\,\dif s>\int_{T_{\vee}(\theta,t,x)}^{t}X_{\theta,t,x,\Right}'(s)\,\dif s.
\end{equation}
We also define
\[
\GS_{\theta}\coloneqq\{(t,x)\in\ssS_{\theta}\st(\theta,t,x)\in\GS\}\quad\text{and}\quad\GS_{\theta,t}\coloneqq\{x\in\ssS_{\theta,t}\st(\theta,t,x)\in\GS\}.
\]
\end{defi}

\begin{figure}[!htbp]
\begin{centering}
\input{figures/global_shock_picture.pgf}
\end{centering}
\caption[The difference between global and non-global shocks]{The global shock curve is shown as a thicker light blue line. The point $(\theta,t,y)$ is a global shock, since the left and right minimizer coming from $(t,y)$, considered up until their first meeting point, together complete a wrap around the torus before meeting again. The point $(\theta,t,x)$ is a shock, since minimizers come from $x$ in multiple directions, but not a global shock, since the minimizers do not accumulate a nontrivial winding before meeting again.
}
\label{fig:global_shock_picture}
\end{figure}

See Figure~\ref{fig:global_shock_picture} for an illustration of the definition of global shock. Additional pictures are available in~\cite{BIK02,BK07}, in particular in the higher-dimensional setting, which is also of significant interest but which we do not consider here.

When we consider all values of $\theta$ simultaneously, it is \emph{not }the case that there exists a unique global shock for each $\theta,t$. In fact, this is impossible, since the global shock must have asymptotic slope $\theta$, and so it cannot vary continuously as $\theta$ is varied. We can nonetheless make a strong statement about the structure of the global shock set $\GS$ if we make the following assumption on the noise, which holds with probability~$1$ for the Poisson point process.
\begin{defi}\label{def:Omegatilde2}
We define $\widetilde{\Omega}_{2}$ to be the set of all $\mu\in\Omega$ such that, for each $\theta\in\RR$, we have
\begin{equation}\label{eq:GSNlt1}
\#(\GS_{\theta}\cap\ssN)\le1
\end{equation}
and
\begin{equation}\label{eq:allshockNsglobal}
\ssS_{\theta}\cap\ssN\setminus\GS_{\theta}=\varnothing.
\end{equation}
We define
\begin{equation}\label{eq:Theta1def}
\Theta_{\otimes}\coloneqq\{\theta\in\RR\st\#(\GS_{\theta}\cap\ssN)=1\},
\end{equation}
and we define maps $s_{\otimes}\colon\Theta_{\otimes}\to\RR$ and $y_{\otimes}\colon\Theta_{\otimes}\to\TT$ by letting $(s_{\otimes}(\theta),y_{\otimes}(\theta))$ be the unique element of $\GS_{\theta}\cap\ssN$ for each $\theta\in\Theta_{\otimes}$.

We further define
\[
\widetilde{\Omega}=\widetilde{\Omega}_{1}\cap\widetilde{\Omega}_{2}.
\]
\end{defi}

\begin{theo}\label{thm:poisson-ok}
Let $\bbP$ be the probability measure associated to a homogeneous compound Poisson point process on $ \RR\times\TT$. Then $\bbP(\widetilde{\Omega})=1$.
\end{theo}

We prove Theorem~\ref{thm:poisson-ok} in Section~\ref{sec:poisson}. Now we can state our second result on the structure of the global shock set.
\begin{theo}\label{thm:shockstructure}
Suppose that $\mu\in\widetilde{\Omega}$. There are unique functions $\sfs_{\Left},\sfs_{\Right}\colon\RR\times\RR\to\TT$ (which we call the \emph{left} and \emph{right global shocks}, respectively) such that the following properties hold:
\begin{enumerate}
\item \label{enu:ssquareshocks}
For each $\theta,t\in\RR$ and $\Square\in\LR$, we have $\GS_{\theta,t}=\{\sfs_{\Left}(\theta,t),\sfs_{\Right}(\theta,t)\}$.
\item \label{enu:shocksleftright}
For each fixed $t\in\RR$ and $\Square\in\LR$, the function $\theta\mapsto\sfs_{\Square}(\theta,t)$ is piecewise continuous. In particular, for each fixed $t\in\RR$,
\begin{equation}\label{eq:sLnotsRdiscrete}
\text{the set }\left\{\theta\in\RR\st\sfs_{\Left}(\theta,t)\ne\sfs_{\Right}(\theta,t)\right\} \text{ is discrete},
\end{equation}
and
\begin{equation}\label{eq:sLsRareleftrightlimits}
\sfs_{\Left}(\theta,t)=\lim_{\theta'\uparrow\theta}\sfs_{\Diamond}(\theta,t)\quad\text{and}\quad\sfs_{\Right}(\theta,t)=\lim_{\theta'\downarrow\theta}\sfs_{\Diamond}(\theta,t)
\end{equation}
for each $\theta\in\RR$ and $\Diamond\in\LR$.
\end{enumerate}
Moreover, these functions have the following additional properties (which are not necessary for the uniqueness statement):
\begin{enumerate}
\setcounter{enumi}{2}
\item \label{enu:shockcontinuity}
For each fixed $\theta\in\RR$ and $\Square\in\LR$, the map $t\mapsto\sfs_{\Square}(\theta,t)$ is continuous.
\item \label{enu:splitatforcing}
If $\theta\in\RR\setminus\Theta_{\otimes}$, then $\sfs_{\Left}(\theta,t)=\sfs_{\Right}(\theta,t)$ for all $t\in\RR$. On the other hand, if $\theta\in\Theta_{\otimes}$, then there exists an $s_{\wedge}(\theta)\in(s_{\otimes}(\theta),\infty)$ such that
\[
\left\{t\in\RR\st\sfs_{\Left}(\theta,t)\ne\sfs_{\Right}(\theta,t)\right\}=\left(s_{\otimes}(\theta),s_{\wedge}(\theta)\right).
\]
\end{enumerate}
\end{theo}

The functions $\sfs_{\Left}$ and $\sfs_{\Right}$ are constructed in Definition~\ref{def:sLsRdefs}, which gives a characterization that is in many ways easier to work with than the characterization given in Theorem~\ref{thm:shockstructure}, but requires some additional definitions to state. Theorem~\ref{thm:shockstructure}$\MK$\eqref{enu:ssquareshocks} is implicit in that definition. The statement~\eqref{eq:sLnotsRdiscrete} is proved as Proposition~\ref{prop:shocksdiffdiscrete}, and~\eqref{eq:sLsRareleftrightlimits} is proved simultaneously with
Proposition~\ref{prop:relatetooneoverderiv} in Section~\ref{subsec:Movement-of-shocks}. Part~\eqref{enu:shockcontinuity} of the theorem statement is proved as Proposition~\ref{prop:smovement}$\MK$\eqref{enu:scts}, and part~\eqref{enu:splitatforcing} is proved in Section~\ref{subsec:movement-global}. The uniqueness statement in Theorem~\ref{thm:shockstructure} holds because part~\eqref{enu:ssquareshocks} and \eqref{eq:sLnotsRdiscrete} characterize, for each $t$, $\rms_{\Left}(\theta,t)$ and $\rms_{\Right}(\theta,t)$ at all except a discrete set of $\theta$, which means that the limits in~\eqref{eq:sLsRareleftrightlimits} are well-defined and characterize $\rms_{\Left}(\theta,t)$ and $\rms_{\Right}(\theta,t)$ for all $\theta$.

The discontinuity of $\sfs_{\Square}(\theta,t)$ is in accordance with the topological obstruction to the continuity of $\sfs_{\Square}(\theta,t)$ in $\theta$ mentioned above. The phenomenon we observe is that, generically, we have $\sfs_{\Left}(\theta,t)=\sfs_{\Right}(\theta,t)$: a single global shock for each $\theta$ and $t$. However, for $\theta\in\Theta_{\otimes}$, there is a time $s_{\otimes}(\theta)$ at which the global shock hits a forcing point and splits into two global shocks $\sfs_{\Left}(\theta,t)\ne\sfs_{\Right}(\theta,t)$. These global shocks then re-merge at a later time $s_{\wedge}(\theta)$, but they may have accumulated a nontrivial winding relative to one another by this merging time. By this last statement we mean that the union of the two branches is not contractible; see Figure~\ref{fig:global_shock_split_picture}.

\begin{figure}[!htbp]
\begin{centering}
\input{figures/global_shock_split_picture.pgf}
\end{centering}
\caption[The global shock splitting and re-merging]{The global shock splitting into two when it hits a forcing point. Later, the two global shocks merge back together. The two branches accumulate different a nontrivial winding relative to one another by the time they re-merge.}
\label{fig:global_shock_split_picture}
\end{figure} 

This motivates the conditions in Definition~\ref{def:Omegatilde2}, which state that for each $\theta$, there can be at most one forcing point that lies on a shock. For fixed $\theta$, this happens with probability $0$, but with probability $1$, it will happen for some values of $\theta$.

The global shock set is the analogue of $g_{\theta}$ defined in~\eqref{eq:gproblem} for the viscous problem. Indeed, $g_{\theta}$ represents the density of a passive particle that has been diffusing in the Burgers flow since time $-\infty$. A similar particle moving in the inviscid Burgers flow (without diffusivity, in accordance with the inviscidity) will end up in the global shock set, since all particles eventually merge with a shock and all shocks eventually merge with the global shock (up to the fact that the global shock may itself split at most once for each $\theta$).

We are now ready to state our main theorem, which describes how the global solutions to the Burgers equation change as $\theta$ is varied.

\begin{theo}\label{thm:mainthm}
Suppose that $\mu\in\widetilde{\Omega}$ and fix $t\in\RR$.
\begin{enumerate}
\item \label{enu:derivintheta}
The functions $\sfs_{\Left}(\cdot,t)$ and $\sfs_{\Right}(\cdot,t)$ are left- and right-differentiable, respectively. For any $\theta_{1}<\theta_{2}$ and any $\Square\in\LR$, we have
\begin{equation}\label{eq:utheta2theta1}
u_{\theta_{2},\Square}(t,x)-u_{\theta_{1},\Square}(t,x)=\sum_{\substack{\theta\in\llbracket\theta_{1},\theta_{2}\rrbracket_{\Square}\\
\sfs_{\Square}(\theta,t)=x } }\frac{1}{\partial_{\theta}\sfs_{\Square}\left(\theta\pm_{\Square},t\right)},
\end{equation}
where we have defined
\begin{equation}\label{eq:pmbookkeeping}
\llbracket\theta_{1},\theta_{2}\rrbracket_{\Square}=
\begin{cases}
(\theta_{1},\theta_{2}] & \text{if }\;\Square=\Left;\\{} [\theta_{1},\theta_{2}) & \text{if }\;\Square=\Right
\end{cases}
\quad\text{and}
\quad\pm_{\Square}=
\begin{cases}
- & \text{if }\;\Square=\Left;\\
+ & \text{if }\;\Square=\Right.
\end{cases}
\end{equation}
\item \label{enu:switching}
If $\sfs_{\Left}(\theta,t)\ne\sfs_{\Right}(\theta,t)$, then there is an $\eps=\eps(\theta,t)>0$ such that
\begin{equation}\label{eq:leftstillshock}
\left(\theta',t,\sfs_{\Left}(\theta,t)\right)\in\ssS\quad\text{for all }\theta'\in[\theta,\theta+\eps)
\end{equation}
and
\begin{equation}\label{eq:rightstillshock}
\left(\theta',t,\sfs_{\Right}(\theta,t)\right)\in\ssS\quad\text{for all }\theta'\in(\theta-\eps,\theta].
\end{equation}
\end{enumerate}
\end{theo}

\begin{figure}[p]
\centering{}
\foreach\mytheta in {-0.51,-0.49,-0.47,-0.45,-0.43,-0.41,-0.39,-0.37,-0.35,-0.33,-0.31,-0.29,-0.27,-0.25,-0.23,-0.21,-0.19,-0.17}{\input{figures/global_shock_moving_picture-\mytheta.pgf}\hfill{}}
\caption[Plots of the global shock set for a sequence of values of $\theta$]{\label{fig:global_shock_moving_picture}
Plots of the global shock set for a sequence of values of $\theta$, which increase as the figures are read like English text, top to bottom and left to right.}
\end{figure}

Let us now describe how \zcref{eq:utheta2theta1} is an inviscid analogue of \zcref{eq:viscouscase}. Suppose for sake of illustration that there is just a single global shock $\sfs(\theta,t)$ for each $\theta$ and $t$, and that it is differentiable as a function of $\theta$. Under this (incorrect) assumption, the analogy of \zcref{eq:viscouscase} would be
\begin{equation}\label{eq:viscouscase-delta}
u_{\theta_{2}}(t,x)-u_{\theta_{1}}(t,x)=\int_{\theta_{1}}^{\theta_{2}}\delta(x-\sfs(\theta,t))\,\dif\theta.
\end{equation}
Here we have used the fact that the inviscid analogue of $g_{\theta}$ is a delta mass at $\sfs(\theta,t)$, as discussed above. Formally performing the change of variables $y=\sfs(\theta,t)$, $\dif y=\partial_{\theta}\sfs(\theta,t)\dif\theta$ in \zcref{eq:viscouscase-delta}, we get
\begin{equation}\label{eq:viscouscase-delta-1}
u_{\theta_{2}}(t,x)-u_{\theta_{1}}(t,x)=\int_{\sfs(\theta_{1},t)}^{\sfs(\theta_{2},t)}\sum_{\substack{\theta\in[\theta_{1},\theta_{2}]\\
\sfs(\theta,t)=y } }\frac{\delta(x-y)}{\partial_{\theta}\sfs(\theta,t)}\,\dif y=\sum_{\substack{\theta\in[\theta_{1},\theta_{2}]\\
\sfs(\theta,t)=x } }\frac{1}{\partial_{\theta}\sfs(\theta,t)},
\end{equation}
which is almost but not precisely correct. Our main result \zcref{eq:utheta2theta1} is a corrected version of the statement: it takes into the account that in the inviscid case, neither $\sfs_{\Square}(\theta,t)$ nor $u_{\theta,\Square}(t,x)$ will be continuous in $\theta$, and selects the left or right versions of these functions as appropriate.

Theorem~\ref{thm:mainthm} implies that, if $t$ and $x$ are held fixed and $\theta$ varies, then the solutions $u_{\theta,\Square}(t,x)$ only change when $\sfs_{\Square}(\theta,t)=x$. Definition~\ref{def:globalshocks} makes transparent the reason for this phenomenon. Indeed, if we differentiate~\eqref{eq:Adef} in $\theta$, we get
\[
\frac{\dif}{\dif\theta}\clA_{\theta,s,t}[X]=-\int_{s}^{t}X'(r)\,\dif r+\theta(t-s).
\]
Therefore, roughly speaking, we see a change in the slope of the minimizer only when there are multiple minimizers with different values of the integral of $X'$ from $t$ until the point at which the minimizers meet, which is exactly what is encoded in~\eqref{eq:shockslopesdiffernt}.

Theorem~\ref{thm:mainthm}$\MK$\eqref{enu:switching} further develops the behavior of the jumps of $\sfs_{\Square}(\theta,t)$. Indeed, it shows that when $\theta\mapsto\sfs_{\Square}(\theta,t)$ has a jump, it jumps to a preexisting non-global shock, which then becomes a global shock and begins to move. See Figure~\ref{fig:global_shock_moving_picture} and also the animation in the supplementary material (described in
Section~\ref{subsec:illustrations}).

\subsection*{Organization of the paper}

In Section~\ref{sec:Existence}, we establish basic facts about the minimizers, and in particular prove Theorem~\ref{thm:minimizers-exist}. In Section~\ref{sec:Continuity-of-minimizers}, we set the stage for our perturbative arguments by using the discreteness of $\ssN$ to show that, as the parameters $\theta,t,x$ are changed, the points used in the minimizer can only change when there are multiple minimizers coming from the same point. In Section~\ref{sec:global-shocks}, we study the topological features of global shocks Definition~\ref{def:globalshocks}. In Section~\ref{sec:Movement-of-shocks}, we study how shocks, and in particular global shocks, move as $t$ is varied. In Section~\ref{sec:udependenceontheta}, we study how the flow changes as $\theta$ is changed. Finally, in
Section~\ref{sec:poisson}, we prove Theorem~\ref{thm:poisson-ok}, showing that our assumptions are satisfied almost surely for a compound Poisson process.

\subsection*{Key to illustrations}
\label{subsec:illustrations}


\typeout{(./anc/moving_global_shock.mp4)}%%mettre en sppp mat!

The paper features several figures and is also accompanied by one animation. In each of the figures, minimizers are drawn as dotted black lines, forcing points as yellow dots with black borders, non-global shocks by dark blue solid lines, and global shocks by thicker light blue lines. Time increases along the vertical axis and space is drawn along the horizontal axis.

In the animation, which is included in the supplementary material, the value of~$\theta$ starts negative and is increased as the animation progresses. The shocks, global shocks, and a sampling of minimizers are drawn, and the graph of $u_{\theta}(t,\cdot)$ for the last plotted time $t$ is also shown above as in Figure~\ref{fig:shocks_picture}. The reader will note that, as proved in Theorem~\ref{thm:mainthm}, the only movement in the picture is through the movement of the global shock. When the global shock hits a forcing point, it jumps to the shock extending from the other side of that forcing point.

\subsection*{Acknowledgments}

The author is very grateful to Yuri Bakhtin, Yu Gu, and Evan Sorensen for many inspiring discussions and important feedback throughout the project, and also to an anonymous referee for helpful remarks on the manuscript.

\section{Existence of one-sided minimizers}
\label{sec:Existence}

In this section we prove Theorem~\ref{thm:minimizers-exist} on the existence of one-sided minimizers. First we must establish some basic properties of minimizers. If $X\in H^{1}([s,r];\TT)$ and $Y\in H^{1}([r,t];\TT)$ are such that $X(r)=Y(r)$, we define the concatenation $X\odot_{r}Y\colon[s,t]\to\TT$ by
\begin{equation}\label{eq:odotdef}
(X\odot_{r}Y)(q)\coloneqq
\begin{cases}
X(q), & \text{if }q\le r;\\
Y(q), & \text{if }q\ge r.
\end{cases}
\end{equation}

\begin{prop}\label{prop:basicproperties}
Suppose that $\mu\in\Omega$. Let $\theta\in\RR$, $-\infty<s<t<+\infty$, and $x,y\in\TT$.
\begin{enumerate}
\item \label{enu:straightlines}
The set $\scM_{s,y\mid t,x}^{\theta}$ is nonempty. Every $X\in\scM_{s,y\mid t,x}^{\theta}$ consists of straight line segments connecting points of $\{(t,x),(s,y)\}\cup\ssN$.
\item \label{enu:subpathalsominimizes}
If $\varnothing\ne[s',t']\subseteq[s,t]$ and $X\in\scM_{s,y\mid t,x}^{\theta}$, then $X|_{[s',t']}\in\scM_{s',X(s')\mid t',X(t')}^{\theta}$ as well.
\item \label{enu:pathsurgery}
If $r\in(s,t)$, $X\in\scM_{s,y\mid t,x}^{\theta}$ and $Y\in\scM_{s,y\mid r,X(r)}^{\theta}$, then $Y\odot_{r}X\in\scM_{s,y\mid t,x}^{\theta}$ as well.
\end{enumerate}
\end{prop}

\begin{proof}
The first point is a standard property of the convexity of the Dirichlet energy in~\eqref{eq:Adef}. For the second point, we note that if not, then we could modify $X$ on $[s',t']$ to improve the value of $\clA_{\theta,s,t}$, contradicting the assumption that $X\in\scM(\theta,s,t)$. To see the third point, note that
\[
\clA_{\theta,s,t}[Y\odot_{r}X]=\clA_{\theta,s,r}[Y]+\clA_{\theta,r,t}[X]\le\clA_{\theta,s,r}[X]+\clA_{\theta,s,r}[X]=\clA_{\theta,s,t}[X]
\]
by the definitions and part~\eqref{enu:subpathalsominimizes}.
\end{proof}

To prove Theorem~\ref{thm:minimizers-exist}, we use the existence of small-noise zones described in Definition~\ref{def:smallnoisezones} to achieve decoupling. The point is that when a small-noise zone occurs, the behavior of polymers inside the small-noise zone is independent of what happens outside of the small zone. This then implies that the behaviors of the polymer before and after the small-noise zone are conditionally independent. The following proposition, whose statement appeared already in~\cite{Bak13} in the case when $\mu(\{(t,x)\})=1$ for all $(t,x)\in\ssN$, is the reason for the definition of $\widetilde{\Omega}_{1}$.


\begin{prop}\label{prop:gothroughy}
Suppose that $\mu\in\Omega$ and that $s\in\RR$, $M>0$, and $y\in\TT$ are such that~\eqref{eq:intersectiononly} and~\eqref{eq:youshouldtake} hold. Then, for any $\theta\in\RR$, $z_{1},z_{2}\in\TT$, and $X\in\scM_{s-M,z_{1}\mid s+M,z_{2}}^{\theta}$, we have $X(s)=y$.
\end{prop}

\begin{proof}
Assume towards a contradiction that we have some $\theta\in\RR$ and $X\in\scM_{s-M,z_{1}\mid s+M,z_{2}}^{\theta}$ such that $X(s)\ne y$. By the assumptions on $s$ and $y$, along with Proposition~\ref{prop:basicproperties}$\MK$\eqref{enu:straightlines}, we see that $X$ consists of a single straight line segment. Let
\begin{equation}\label{eq:xirange}
\xi\in(-1/2,1/2]
\end{equation}
be such that $X(s)+\xi=y$, and define
\[
Y(r)\coloneqq X(r)+\xi\cdot
\begin{cases}
\frac{s+M-r}{M} & \text{if }r\in[s,s+M];\\
\frac{r-s+M}{M} & \text{if }r\in[s-M,s];\\
0 & \text{otherwise}.
\end{cases}
\]
In particular, this means that $Y(s-M)=X(s-M)=z_{1}$, $Y(s+M)=X(s+M)=z_{2}$, and $Y(s)=X(s)+\xi=y$. Then we have
\begin{align*}
&\clA_{\theta,s-M,s+M}[Y]-\clA_{\theta,s-M,s+M}[X]\\
& \; =M\left[\frac{1}{2}\left(X'(s)-\theta+\xi/M\right)^{2}\!+\frac{1}{2}\left(X'(s)-\theta-\xi/M)^{2}-(X'(s)-\theta\right)^{2}\right]-\mu(\{(s,y)\})\\
& \; =\frac{\xi^{2}}{M}-\mu(\{(s,y)\})\overset{\zcref{eq:xirange}}{\le}\frac{1}{4M}-\mu(\{(s,y)\})\overset{\zcref{eq:youshouldtake}}{<}0.
\end{align*}
But since $Y(s\pm M)=X(s\pm M)$, this contradicts the assumption $X\in\scM_{s-M,z_{1}\mid s+M,z_{2}}^{\theta}$.
\end{proof}

Now we can make the following important definition.

\begin{defi}\label{def:Tstar}
Suppose that $\mu\in\widetilde{\Omega}_{1}$ and let $M(\mu)$ be as in Definition~\ref{def:smallnoisezones}. For $t\in\RR$, we define $T_{*}(t)$ to be the supremum of all $s<t$ such that there exists a $y\in\TT$ such that $(s,y)\in\ssN$ and $X(s)=y$ for all
\[
X\in\bigcup_{\substack{\theta\in\RR\\
x,z\in\TT } }\scM_{s-M(\mu),z\mid t,x}^{\theta}.
\]
\end{defi}

We note that it is an immediate consequence of Definition~\ref{def:Tstar} and the definition~\eqref{eq:Tveedef} of $T_{\vee}$ that
\begin{equation}\label{eq:TstarTvee}
T_{*}(t)\le T_{\vee}(\theta,t,x)\quad\text{for any }\theta,t\in\RR\;\text{and}\; x\in\TT.
\end{equation}
The point of Proposition~\ref{prop:gothroughy} is that $T_{*}(t)$ is finite:


\begin{prop}\label{prop:Tstartfinite}
Suppose that $\mu\in\widetilde{\Omega}_{1}$ and $t\in\RR$.
\begin{enumerate}
\item We have $T_{*}(t)>-\infty$.
\item There is a $y_{*}(t)\in\TT$ such that
\begin{equation}\label{eq:allpassthrough}
X(T_{*}(t))=y_{*}(t)\text{ for all }X\in\bigcup_{\substack{\theta\in\RR\\
x,z\in\TT } }\scM_{T_{*}(t)-M(\mu),z\mid t,x}^{\theta}.
\end{equation}
\item Finally, we have
\begin{equation}\label{eq:Tstarttoinfinity}
\lim_{t\to+\infty}T_{*}(t)=+\infty.
\end{equation}
\end{enumerate}
\end{prop}

\begin{proof}
Fix $t\in\RR$. Let $M>0$, $s\in(-\infty,t-M]$, and $y\in\TT$ be as in the definition of $\widetilde{\Omega}$. We claim that in fact $T_{*}(t)\ge s$. Proposition~\ref{prop:basicproperties}$\MK$\eqref{enu:subpathalsominimizes} implies that if $\theta\in\RR$, $x,z\in\TT$, and $X\in\scM_{s-M,z\mid t,x}^{\theta}$, then $X\in\scM_{s-M,z\mid s+M,X(s+M)}^{\theta}$ as well, and so by
Proposition~\ref{prop:Tstartfinite} we have $X(s)=y$. Since $(s,y)\in\ssN$ by definition, we have $T_{*}(t)\ge s>-\infty$, and the first assertion is proved. The second assertion is tantamount to asserting that the supremum in Definition~\ref{def:Tstar} is achieved; this follows from the discreteness of $\ssN$ and the compactness of $\TT$. Finally, \eqref{eq:Tstarttoinfinity} follows from~\eqref{eq:stttoinfty}.
\end{proof}

The following proposition contains the statement of Theorem~\ref{thm:minimizers-exist}.

\begin{prop}\label{prop:globalminimizersexist}
Suppose that $\mu\in\widetilde{\Omega}_{1}$. For each $\theta,t\in\RR$ and $x\in\TT$, the set $\scM_{t,x}^{\theta}$ is nonempty. In particular, we have
\begin{equation}\label{eq:Mstarisrestriction}
\left\{X|_{[T_{*}(t),t]}\st X\in\scM_{t,x}^{\theta}\right\}=\scM_{T_{*}(t),y_{*}(t)\mid t,x}^{\theta}.
\end{equation}
Moreover, $\scM_{t,x}^{\theta}$ has minimal and maximal elements under the partial order $\preccurlyeq$.
\end{prop}

\begin{proof}
Since $\scM_{T_{*}(t),y_{*}(t)\mid t,x}^{\theta}$ is nonempty by
Proposition~\ref{prop:basicproperties}$\MK$\eqref{enu:straightlines}, to prove that the set $\scM_{t,x}^{\theta}$ is nonempty it suffices to prove~\eqref{eq:Mstarisrestriction}. The ``$\subseteq$'' direction is an immediate consequence of Proposition~\ref{prop:basicproperties}$\MK$\eqref{enu:subpathalsominimizes} and the definition of $\scM_{t,x}^{\theta}$, so we turn our attention to the ``$\supseteq$'' direction. In other words, given $X_{1}\in\scM_{T_{*}(t),y_{*}(t)\mid t,x}^{\theta}$, we seek to extend $X_{1}$ to an element $X$ of $\scM_{t,x}^{\theta}$. Let $t_{0}=t$, and let $t_{k}=T_{*}(t_{k-1})$ and $y_{k}=y_{*}(t_{k-1})$ for $k\ge1$. We note that, since $t_{k}\le t_{k-1}-M(\mu)$, we have
\begin{equation}\label{eq:tksfillup}
\bigsqcup_{k=1}^{\infty}(t_{k},t_{k-1}]=(-\infty,t].
\end{equation}
For $k\ge2$, let $X_{k}\in\scM_{t_{k},y_{k}\mid t_{k-1},y_{k-1}}^{\theta}$. (This inclusion is satisfied for $k=1$ as well, but $X_{1}$ has already been chosen.) Now for $s\in[t_{k},t_{k-1})$, define $X(s)\coloneqq X_{k}(s)$, so $X$ is defined as an element of $H_{\locc}^{1}((-\infty,t];\TT)$ by~\eqref{eq:tksfillup} and the fact that $X_{k}(t_{k-1})=y_{k-1}=X_{k-1}(t_{k-1})$ by definition.

We claim that $X\in\scM_{t,x}^{\theta}$. Let $s<t$ and let $k$ be large enough that $t_{k}\le s$. Suppose that $z\in\TT$ and that $Y\in\scM_{t_{k}-M,z\mid t,x}^{\theta}$. Then, by~\eqref{eq:allpassthrough}, we have $Y(t_{j})=y_{j}=X(t_{j})$ whenever $j\le k$. This means that $Y|_{[t_{j},t_{j-1}]}\in\scM_{t_{k},y_{k}\mid t_{k-1},y_{k-1}}^{\theta}$. Since the same is true for $X|_{[t_{j},t_{j-1}]}$, we have $\clA_{\theta,t_{j},t_{j-1}}[X]=\clA_{\theta,t_{j},t_{j-1}}[Y]$. Summing this up over all $j$, we obtain
\begin{equation}\label{eq:YXsame}
\clA_{\theta,t_{k},t}[Y]=\clA_{\theta,t_{k},t}[X].
\end{equation}
Now since $Y\in\scM_{t_{k},y_{k}\mid t,x}^{\theta}$ by Proposition~\ref{prop:basicproperties}$\MK$\eqref{enu:subpathalsominimizes},
\eqref{eq:YXsame} means that $X\in\scM_{t_{k},y_{k}\mid t,x}^{\theta}$ as well. But since $t_{k}\le s$, we can apply Proposition~\ref{prop:basicproperties}$\MK$\eqref{enu:subpathalsominimizes} once again to see that $X\in\scM_{s,X(s)\mid t,x}^{\theta}$ Since this is true for any $s<t$, we conclude that $X\in\scM_{t,x}^{\theta}$.

To show that $\scM_{t,x}^{\theta}$ has minimal/leftmost and maximal/rightmost elements under~$\preccurlyeq$, we observe that each $\scM_{t_{k},y_{k}\mid t_{k-1},y_{k-1}}^{\theta}$ is finite and has leftmost and rightmost elements (as can be seen using Proposition~\ref{prop:basicproperties}$\MK$\eqref{enu:pathsurgery} to build a path that is weakly to the left/right of any other minimizer). Then we note that a concatenation of these leftmost and rightmost elements in a similar manner to the above argument will yield leftmost and rightmost elements of $\scM_{t,x}^{\theta}$.
\end{proof}

The following corollary emphasizes how times of the form $T_{*}(t)$ serve as regeneration times such that the behavior of minimizers before and after them is independent.

\begin{coro}\label{cor:justdependsonendstuff}
If $\mu\in\widetilde{\Omega}_{1}$, $\theta\in\RR$, $x\in\TT$, and $s\le t$, then
\begin{equation}\label{eq:Misconcat}
\begin{aligned}
\scM_{t,x}^{\theta}&=\scM_{T_{*}(s),y_{*}(s)}^{\theta}\odot_{T_{*}(s)}\scM_{T_{*}(s),y_{*}(s)\mid t,x}^{\theta}\\
&=\left\{X\odot_{T_{*}(s)}Y\st X\in\scM_{T_{*}(s),y_{*}(s)}^{\theta}\quad \text{and}\quad Y\in\scM_{T_{*}(s),y_{*}(s)\mid t,x}^{\theta}\right\}.
\end{aligned}
\end{equation}
In particular, if we make the shorthand definition
\begin{equation}\label{eq:Mstartx}
\scM_{*\mid t,x}^{\theta}\coloneqq\scM_{T_{*}(t),y_{*}(t)\mid t,x}^{\theta},
\end{equation}
then we have
\[
\scM_{t,x}^{\theta}=\scM_{T_{*}(s),y_{*}(s)}^{\theta}\odot_{T_{*}(s)}\scM_{*\mid t,x}^{\theta}.
\]
\end{coro}

\begin{proof}
This follows from Proposition~\ref{prop:globalminimizersexist}, Proposition~\ref{prop:basicproperties}$\MK$\eqref{enu:pathsurgery}, and induction.
\end{proof}

%? pourquoi ya ça ?\FloatBarrier{}

\subsection{Minimizers and shocks}

Now that we have established the existence of global minimizers, the definition \eqref{eq:Sdef} of the shock set $\ssS$ makes sense. Here we establish a few basic properties about minimizers and shocks that will be useful in the sequel.


\begin{lemm}\label{lem:whencanyousplit}
Suppose that $\mu\in\widetilde{\Omega}_{1}$. Let $\theta\in\RR$, $s<t$, and $X_{1},X_{2}\in\scM_{t}^{\theta}$. If $r\in[s,t]$ is such that $X_{1}(r)=X_{2}(r)$ but $X_{1}'(r-)\ne X_{2}'(r-)$, then $(r,X_{i}(r))\in\ssS_{\theta}$. If we moreover assume that $r<t$, then $(r,X_{i}(r))\in\ssN$ as well.
\end{lemm}

\begin{proof}
It follows from the definitions and Proposition~\ref{prop:basicproperties}$\MK$\eqref{enu:subpathalsominimizes} that
\begin{equation}\label{eq:minimizelater}
X_{i}|_{(-\infty,r]}\in\scM_{r,X_{i}(r)}^{\theta}\qquad\text{for each }i\in\{1,2\},
\end{equation}
so $(\theta,r,X_{i}(r))\in\ssS$ by the definition of $\ssS$.

Now we assume that $r<t$ and prove that $(r,X_{i}(r))\in\ssN$. We note that the restrictions of $X_{1}$ and $X_{2}$ are both elements of $\scM_{T_{*}(r),y_{*}(r)\mid r,X_{i}(r)}^{\theta}$ by~\eqref{eq:minimizelater} and~\eqref{eq:Mstarisrestriction}, so Proposition~\ref{prop:basicproperties}$\MK$\eqref{enu:pathsurgery} and the assumption that $X_{1}'(r-)\ne X_{2}'(r-)$ imply that there is an element of $\scM_{T_{*}(r),y_{*}(r)\mid t,X_{1}(t)}^{\theta}$ that changes direction at $(r,X_{1}(r))$. Hence Proposition~\ref{prop:basicproperties}$\MK$\eqref{enu:straightlines} implies that $(r,X_{i}(r))\in\ssN$.
\end{proof}

\begin{prop}\label{prop:allshocksmerge}
Suppose that $\mu\in\widetilde{\Omega}_{1}$. For any $\theta,t_{0}\in\RR$, there is a $t>t_{0}$ such that $\#\GS_{\theta,t}=\#\ssS_{\theta,t}=1$.
\end{prop}

\begin{proof}
Let $M$ be as in Definition~\ref{def:smallnoisezones}. Using~\eqref{eq:stttoinfty}, we can find an $s\ge t_{0}$ such that~\eqref{eq:intersectiononly} and~\eqref{eq:youshouldtake} hold. By Proposition~\ref{prop:gothroughy}, we have $X(s)=y$ for any $X\in\scM_{t}^{\theta}$. Thus all minimizers at time $s+M$ begin with straight line segments to $(s,y)$, and there is a single shock at the point where the line segments switch the direction they go around the torus, as shown in
Figure~\ref{fig:single_shock_picture}.\qedhere
\end{proof}

\begin{figure}[!htbp]
\begin{centering}
\input{figures/single_shock_picture.pgf}
\end{centering}
\caption[Sufficient criterion to have a single global shock at a given time]{When all minimizers at time $s+M$ start with straight line segments to $(s,y)$, then there is a single shock, which is in fact a global shock, at time $s+M$.
}
\label{fig:single_shock_picture}
\end{figure}



\section{Continuity of minimizers with respect to parameters}
\label{sec:Continuity-of-minimizers}
We now want to explore how the sets $\scM_{t,x}^{\theta}$ change as we vary $\theta$, $t$, and $x$. The main result of this section is that, if $\theta,t,x$ are perturbed only slightly, then each new minimizer uses the same forcing points as one of the original minimizers.

For a path $X$ in $H_{\locc}^{1}((-\infty,t])$, $t_{0}\le t$, $\tau\in(t_{0}-t,+\infty)$, and $\eta\in\RR$, if $X|_{[t_{0},t]}$ is linear, we define a new path $\clT_{t_{0},\tau,\eta}X\in H_{\locc}^{1}((-\infty,t+\tau])$ by
\[
\clT_{t_{0},\tau,\eta}X(s)\coloneqq
\begin{cases}
X(s), & s\le t_{0};\\
X(t_{0})+(s-t_{0})\cdot\frac{(t-t_{0})X'(t-)+\eta}{t+\tau-t_{0}}, & s\in[t_{0},t+\tau].
\end{cases}
\]
This means in particular that
\[
\clT_{t_{0},\tau,\eta}X(t+\tau)=X(t_{0})+(t-t_{0})X'(t-)+\eta=X(t)+\eta
\]
(here we use the assumption that $X|_{[t_{0},t]}$ is linear) and
\begin{equation}\label{eq:TXderiv}
(\clT_{t_{0},\tau,\eta}X)'(t+\tau)=\frac{(t-t_{0})X'(t-)+\eta}{t+\tau-t_{0}}.
\end{equation}
We also define
\begin{equation}\label{eq:tminus1}
t_{-;1}[X]\coloneqq\min\{r<t\st(r,X(r))\in\ssN\}
\end{equation}
and
\begin{equation}\label{eq:tminus2}
t_{-;2}[X]\coloneqq\min\{r\le t\st(r,X(r))\in\ssN\}
\end{equation}
and, for $i=1,2$,
\begin{equation}\label{eq:Tidef}
\clT_{\tau,\eta}^{i}X=\clT_{t_{-;i}[X],\tau,\eta}X.
\end{equation}

We note that if $(t,X(t))\not\in\ssN$, then $\clT_{\tau,\eta}^{i}X$ does not depend on $i$. If $(t,X(t))\in\ssN$, then both $\clT_{\tau,\eta}^{1}$ and $\clT_{\tau,\eta}^{2}$ move the endpoint of $X$ to $(t+\tau,X(t)+\eta)$, but $\clT_{\tau,\eta}^{2}X$ keeps the forcing point at $(t,X(t))$, while $\clT_{\tau,\eta}^{1}X$ skips over it. See Figure~\ref{fig:move_endpoint_picture}.


\begin{figure}[!htbp]
\begin{centering}
\input{figures/move_endpoint_picture.pgf}
\end{centering}
\caption[Distinction between different notions of perturbing a minimizer]{The distinction between $\clT_{\tau,\eta}^{1}X$ and $\clT_{\tau,\eta}^{2}X$ when $(t,X(t))\in\protect\ssN$.}
\label{fig:move_endpoint_picture}
\end{figure}

Now we can state our proposition.

\begin{prop}\label{prop:main-continuity}
Suppose that $\mu\in\widetilde{\Omega}_{1}$. Fix $\theta,t\in\RR$ and $x\in\TT$. There exists an $\eps=\eps(\mu,\theta,t,x)\in(0,t-t_{-;1}[X])$ such that whenever $\zeta,\tau,\eta\in(-\eps,\eps)$, we have
\begin{equation}\label{eq:main-continuity}
\scM_{t+\tau,x+\eta}^{\theta+\zeta}\subseteq
\begin{cases}
\clT_{\tau,\eta}^{1}\scM_{t,x}^{\theta} & \text{if }(t,x)\not\in\ssN\text{ or }\tau\le0;\\
\clT_{\tau,\eta}^{1}\scM_{t,x}^{\theta}\cup\clT_{\tau,\eta}^{2}\scM_{t,x}^{\theta} & \text{otherwise.}
\end{cases}
\end{equation}
\end{prop}

The reason for the two cases on the right side of~\eqref{eq:main-continuity} is that $\clT_{\tau,\eta}^{2}X$ is not defined for $\tau\le0$, and if $(t,x)\not\in\ssN$, then $\clT_{\tau,\eta}^{1}=\clT_{\tau,\eta}^{2}$ by definition.
\begin{proof}
First we choose $\eps_{0}>0$ small enough that
\begin{multline}\label{eq:eps0cond}
\left\{ (s,X(s))\st\zeta,\tau,\eta\in(-\eps_{0},\eps_{0}),X\in\scM_{t+\tau,x+\eta}^{\theta+\zeta},s\in[t-\eps_{0},t+\eps_{0}]\right\}\cap\ssN\\
\subseteq\{(t,x)\}.
\end{multline}
This is possible since $\ssN$ is discrete: first we choose $\eps_{0}$ small enough that any element $(s,y)\in\ssN\cap([t-\eps_{0},t+\eps_{0}]\times\TT)$ must have $s=t$, and then we can make $\eps_{0}$ even smaller if necessary to ensure that it is not advantageous for a minimizer in $\scM_{t+\tau,x+\eta}^{\theta+\zeta}$ to use any forcing point $(t,y)$ with $y\ne x$.

Now let $T_{0}=T_{*}(t-\eps_{0})$ and $y_{0}=y_{*}(t-\eps_{0})$. By~\eqref{eq:Mstarisrestriction}, whenever $\zeta,\tau,\eta\in(-\eps_{0},\eps_{0})$, we have $\scM_{t+\tau,x+\eta}^{\theta+\zeta}=\scM_{T_{0},y_{0}}^{\theta+\zeta}\odot_{T_{0}}\scM_{T_{0},y_{0}\mid t+\tau,x+\eta}^{\theta+\zeta}$, so to complete the proof of the proposition it suffices to show that, when $\zeta,\tau,\eta$ are sufficiently small, we have
\begin{equation}\label{eq:goal}
\scM_{T_{0},y_{0}\mid t+\tau,x+\eta}^{\theta+\zeta}\subseteq
\begin{cases}
\clT_{\tau,\eta}^{1}\scM & \text{if }(t,x)\not\in\ssN\text{ or }\tau\le0;\\
\clT_{\tau,\eta}^{1}\scM\cup\clT_{\tau,\eta}^{2}\scM & \text{otherwise,}
\end{cases}
\end{equation}
where we have defined
\begin{equation}\label{eq:Mdef}
\scM\coloneqq\scM_{T_{0},y_{0}\mid t,x}^{\theta}.
\end{equation}

Fix $X\in\scM$. Let $\scN$ be the set of \emph{all} paths $Y$ with $Y(T_{0})=y_{0}$, $Y(t)=x$, and $Y$ consisting of straight line segments connecting a subset of the points of $\{(t,x),(T_{0},y_{0})\}\cup\ssN\cap([T_{0},t+\eps_{0}]\times\TT)$. We see from the definitions and Proposition~\ref{prop:basicproperties}$\MK$\eqref{enu:straightlines} that
\begin{equation}\label{eq:cangetallfromshifting}
\scM_{T_{0},y_{0}\mid t+\tau,x+\eta}^{\theta+\zeta}\subseteq
\begin{cases}
\clT_{\tau,\eta}^{1}\scN & \text{if }(t,x)\not\in\ssN\text{ or }\tau\le0;\\
\clT_{\tau,\eta}^{1}\scN\cup\clT_{\tau,\eta}^{2}\scN & \text{otherwise.}
\end{cases}.
\end{equation}

The discreteness of $\ssN$ implies that
\begin{equation}\label{eq:deltadef}
\min_{Y\in\scN\setminus\scM}\clA_{\theta,T_{0},t}[Y]-\clA_{\theta,T_{0},t}[X]>0.
\end{equation}
Now~\eqref{eq:eps0cond} implies that, for each fixed $Y\in\scN$, the map
\[
(\zeta,\tau,\eta)\mapsto\clA_{\theta,T,t+\tau}[\clT_{\tau,\eta}^{1}Y]
\]
is continuous on $[-\eps_{0}/2,\eps_{0}/2]^{3}$. Combining this observation with~\eqref{eq:deltadef} and using the discreteness of $\ssN$ again, we can find an $\eps\in(0,\eps_{0}/2)$ such that, if $\zeta,\tau,\eta\in(-\eps,\eps)$, then
\[
\min_{Y\in\scN\setminus\scM}\clA_{\theta,T_{0},t+\tau}[\clT_{\tau,\eta}^{1}Y]-\clA_{\theta,T_{0},t+\tau}[\clT_{\tau,\eta}^{1}X]>0,
\]
and so
\begin{equation}\label{eq:nootherguysT1}
\left\{\clT_{\tau,\eta}^{1}Y\st Y\in\scN\setminus\scM\right\}\cap\scM_{T_{0},y_{0}\mid t+\tau,x+\eta}^{\theta+\zeta}=\varnothing.
\end{equation}
Combining this with~\eqref{eq:cangetallfromshifting} and recalling that $\clT_{\tau,\eta}^{1}Y=\clT_{\tau,\eta}^{2}Y$ whenever $(t,x)\not\in\ssN$, we see that~\eqref{eq:goal} is proved in the case when $(t,x)\not\in\ssN$, and also when $\tau\le t$.

Thus we now have to consider the case when $(t,x)\in\ssN$ and $\tau>t$. By~\eqref{eq:nootherguysT1}, we see that
\begin{equation}\label{eq:cangetallfromshifting-1}
\scM_{T_{0},y_{0}\mid t+\tau,x+\eta}^{\theta+\zeta}\subseteq\clT_{\tau,\eta}^{1}\scM\cup\left\{\clT_{\tau,\eta}^{2}Y\st Y\in\scN\right\}.
\end{equation}
If $Z\in\scM_{T_{0},y_{0}\mid t+\tau,x+\eta}^{\theta+\zeta}\cap\{\clT_{\tau,\eta}^{2}Y\st Y\in\scN\}$, then in particular $Z(t)=x$, so by Proposition~\ref{prop:basicproperties}$\MK$\eqref{enu:subpathalsominimizes}, we have $Z|_{[T_{0},t]}\in\scM_{T_{0},y_{0}\mid t,x}^{\theta+\zeta}$. On the other hand, by the case $\tau\le t$ considered above, we have $\scM_{T_{0},y_{0}\mid t,x}^{\theta+\zeta}\subseteq\scM$. Therefore, we have $Z=\clT_{\tau,\eta}^{2}(Z|_{[T_{0},t]})\in\clT_{\tau,\eta}^{2}\scM$. This completes the proof of~\eqref{eq:goal}.
\end{proof}

We will apply Proposition~\ref{prop:main-continuity} many times throughout the paper. The following application is quite simple.

\begin{coro}\label{cor:minimizers-movex}
Suppose that $\mu\in\widetilde{\Omega}_{1}$. Fix $\theta,t\in\RR$ and $x\in\TT$. There exists an $\eps=\eps(\mu,\theta,t,x)>0$ such that whenever $\eta\in[0,\eps)$,
\begin{equation}\label{eq:Mx+eta}
\scM_{t,x+\eta}^{\theta}=
\left\{\clT_{0,\eta}^{1}X\st X\in\scM_{t,x}^{\theta}\quad\text{and}\quad X'(t-)=X_{\theta,t,x,\Right}'(t-)\right\}
\end{equation}
and
\begin{equation}\label{eq:Mx-eta}
\scM_{t,x-\eta}^{\theta}=
\left\{\clT_{0,-\eta}^{1}X\st X\in\scM_{t,x}^{\theta}\quad\text{and}\quad X'(t-)=X_{\theta,t,x,\Left}'(t-)\right\}.
\end{equation}
In particular,
\begin{equation}\label{eq:Xx+eta}
X_{\theta,t,x+\eta,\Right}=\clT_{0,\eta}^{1}X_{\theta,t,x,\Right}
\end{equation}
and
\begin{equation}\label{eq:Xx-eta}
X_{\theta,t,x-\eta,\Left}=\clT_{0,-\eta}^{1}X_{\theta,t,x,\Left}.
\end{equation}
\end{coro}

\begin{proof}
We prove the first statement, as the proof of the second is symmetrical, and the third and fourth statements follow immediately from the first and second, respectively. Let $\eps$ be as in Proposition~\ref{prop:main-continuity} (although we will make it even smaller shortly). By Corollary~\ref{cor:justdependsonendstuff}, $c\coloneqq\clA_{\theta,T_{*}(t),x}[X]$ does not depend on $X\in\scM_{t,x}^{\theta}$. Thus, for any $\eta\in[0,\eps)$ and any $X\in\scM_{t,x}^{\theta}$, we can compute using the definitions that
\begin{multline*}
\clA_{\theta,T_{*}(t),x+\eta}[\clT_{0,\eta}^{1}X]-c \\
\begin{aligned}
& =\frac{1}{2}(t-t_{-;1}[X])\left[\left(X'(t-)-\theta+\frac{\eta}{t-t_{-;1}[X]}\right)^{2}-\left(X'(t-)-\theta\right)^{2}\right]\\
& =\eta(X'(t-)-\theta)+\frac{\eta^{2}}{2(t-t_{-;1}[X])}.
\end{aligned}
\end{multline*}
From this expression we see that, when $\eta$ is sufficiently small and nonnegative, the action $\clA_{\theta,T_{*}(t),x+\eta}[\clT_{0,\eta}^{1}X]$ is minimized over $X\in\scM_{t,x}^{\theta}$ exactly when $X'(t-)=X_{\theta,t,x,\Right}'(t-)$ (as all such $X$ will have the same value of $t_{-;1}[X]$). Using this along with Proposition~\ref{prop:main-continuity} yields the conclusion.
\end{proof}

\section{Topology of global shocks}
\label{sec:global-shocks}

In this section we explore the properties of the global shock set $\GS$ introduced in Definition~\ref{def:globalshocks}.

\subsection{Basic properties}

We begin with two refinements of the definition of global shock. We recall that a global shock is one for which the left and right minimizers, when they merge back together, do so with different winding numbers. We now show that the difference in winding number is always exactly $1$.

\begin{prop}\label{prop:exactlyoneoff}
Suppose that $\mu\in\Omega$. If $(\theta,t,x)\in\GS$, then
\begin{equation}\label{eq:exactlyoneoff}
\int_{T_{\vee}(\theta,t,x)}^{t}X'_{\theta,t,x,\Left}(s)\,\dif s=\int_{T_{\vee}(\theta,t,x)}^{t}X'_{\theta,t,x,\Right}(s)\,\dif s+1.
\end{equation}
\end{prop}

\begin{proof}
For $\Square\in\LR$, let $I_{\Square}\coloneqq\int_{T_{\vee}(\theta,t,x)}^{t}X'_{\theta,t,x,\Square}(s)\,\dif s$. Using the fact that
\[
X_{\theta,t,x,\Left}(T_{\vee}(\theta,t,x))=X_{\theta,t,x,\Right}(T_{\vee}(\theta,t,x))
\]
along with~\eqref{eq:shockslopesdiffernt}, we see that $I_{\Left}-I_{\Right}$ is a positive integer. On the other hand, if $I_{\Left}-I_{\Right}>1$, then the intermediate value theorem yields an $s\in(T_{\vee}(\theta,t,x),t)$ such that
\[
\int_{s}^{t}X'_{\theta,t,x,\Left}(r)\,\dif r-\int_{s}^{t}X'_{\theta,t,x,\Right}(r)\,\dif r=1,
\]
and then we would have $X_{\theta,t,x,\Left}(s)=X_{\theta,t,x,\Right}(s)$, contradicting the definition~\eqref{eq:Tveedef} of $T_{\vee}(\theta,t,x)$. Thus we must have $I_{\Left}-I_{\Right}=1$, which was the claim.
\end{proof}

The next proposition says that, in order to check that a shock is a global shock, it is sufficient to find two minimizers that have different slopes at time $t$ and a positive winding number upon their first reconnection. (They do not have to be the left and right minimizers.)

\begin{figure}[!htbp]
\begin{centering}
\scalebox{0.9}{\input{figures/XLXXR_picture.pgf}}
\end{centering}
\caption[Illustration of the situation in Proposition~\ref{prop:XLXXR}]{Illustration of the situation in Proposition~\ref{prop:XLXXR}.}
\label{fig:XLXXR_picture}
\end{figure}

\begin{prop}\label{prop:XLXXR}
Suppose that $\mu\in\widetilde{\Omega}_{1}$. If there are $X,Y\in\scM_{t,x}^{\theta}$ such that $X'(t)>Y'(t)$ and, with $T_{X,Y}\coloneqq\max\{s<t\st X(s)=Y(s)\}$, we assume that
\begin{equation}\label{eq:XYdiscrep}
\int_{T_{X,Y}}^{t}X'(s)\,\dif s\ne\int_{T_{X,Y}}^{t}Y'(s)\,\dif s,
\end{equation}
then $(\theta,t,x)\in\GS$.
\end{prop}

\begin{proof}
We first note that
\begin{equation}\label{eq:XYdiscrep-pos}
\int_{r}^{t}(X'-Y')(s)\,\dif s>0\quad\text{for all }r\in[T_{X,Y},t]
\end{equation}
by the intermediate value theorem, the assumption that $X'(t)>Y'(t)$, and~\eqref{eq:XYdiscrep}, since if there were a point $r\in(T_{X,Y},t]$ such that $\int_{r}^{t}(X'-Y')(s)\,\dif s=0$ then we would have $X(r)=Y(r)$ and hence $T_{X,Y}\ge r$. In particular, this means that~\eqref{eq:XYdiscrep} can be refined to
\begin{equation}\label{eq:XYdiscrep-sign}
\int_{T_{X,Y}}^{t}X'(s)\,\dif s>\int_{T_{X,Y}}^{t}Y'(s)\,\dif s
\end{equation}
We also note that
\[
X_{\theta,t,x,\Left}'(t-)\ge X'(t-)>Y'(t-)\ge X_{\theta,t,x,\Right}'(t),
\]
so $(\theta,t,x)\in\ssS$. Furthermore,
\[
\int_{T_{X,Y}}^{t}\left(X_{\theta,t,x,\Left}'-X_{\theta,t,x,\Right}'\right)(s)\,\dif s\ge\int_{T_{X,Y}}^{t}(X'-Y')(s)\,\dif s\ge1
\]
(with the last inequality by~\eqref{eq:XYdiscrep-sign} and the fact that the last integral must be an integer), and so $T_{\vee}(\theta,t,x)\ge T_{X,Y}$. (See Figure~\ref{fig:XLXXR_picture}.)


But then we must in fact have
\[
\int_{T_{\vee}(\theta,t,x)}^{t}\left(X_{\theta,t,x,\Left}'-X_{\theta,t,x,\Right}' \right)(s)\,\dif s\ge\int_{T_{\vee}(\theta,t,x)}^{t}(X'-Y')(s)\,\dif s\overset{\zcref{eq:XYdiscrep-pos}}{>}0,
\]
and the proof is complete.
\end{proof}



\subsection{Existence of global shocks}

Now we show that there must be at least one global shock for each $\theta$ and $t$. The reason for this is topological: since all minimizers for a given $\theta$ and $t$ must merge, there must be some point at which the minimizers switch the direction they travel around the torus.
\begin{prop}\label{prop:findglobalshock}
Suppose that $\mu\in\widetilde{\Omega}$. For each $\theta,t\in\RR$, we have $\#\GS_{\theta,t}\ge1$.
\end{prop}

\begin{proof}
Fix $\theta,t\in\RR$. For $x\in\TT$ and $\Square\in\LR$, let
\[
f_{\Square}(x)=\int_{T_{*}(t)}^{t}X_{\theta,t,x,\Square}'(s)\,\dif s.
\]
Now since $y\coloneqq X_{\theta,t,x,\Square}(T_{*}(t))$ depends neither on $x$ nor on $\Square$, we know that $f_{\Square}(x)\in x-y+\ZZ$ for each $\Square,x$.\footnote{Since $x,y\in\TT$, $x-y$ is not well-defined, but $x-y+\ZZ$ is.} This implies that there exists an $x\in\TT$ such that $f_{\Left}(x)-f_{\Right}(x)\in\ZZ\setminus\{0\}$. Let
\begin{equation}\label{eq:s0def}
s_{\otimes}=\min\left\{ s\in[T_{*}(t),t]\st\quad
\begin{multlined}
X_{\theta,t,x,\Left}(s)=X_{\theta,t,x,\Right}(s)\\
\text{and}\quad\int_{T_{*}(t)}^{s}\left(X_{\theta,t,x,\Left}'-X_{\theta,t,x,\Right}'\right)(s)\,\dif s\in\ZZ\setminus\{0\}
\end{multlined}
\right\}.
\end{equation}
Let $y_{\otimes}=X_{\theta,t,x,\Left}(s_{\otimes})=X_{\theta,t,x,\Right}(s_{\otimes})$. Then we can see from the definitions and Lemma~\ref{lem:whencanyousplit} that $(s_{\otimes},y_{\otimes})\in\GS_{\theta}$. If $s_{\otimes}=t$, then we are done, so assume that $s_{\otimes}<t$. In this case, Lemma~\ref{lem:whencanyousplit} further implies that $(s_{\otimes},y_{\otimes})\in\ssN$. (Thus we have $(s_{\otimes},y_{\otimes})=(s_{\otimes}(\theta),y_{\otimes}(\theta))$ as in Definition~\ref{def:Omegatilde2}, justifying the notation.) On the other hand, the definition~\eqref{eq:s0def} implies that $s_{\otimes}>T_{*}(t)$, so in fact
\[
\eta\coloneqq\min\left\{\xi\ge0\st\text{there is an }X\in\scM_{t,x+\xi}^{\theta}\text{ such that }X(s_{\otimes})\ne y_{\otimes}\right\}
\]
exists. Define $x'=x+\eta$. By Proposition~\ref{prop:main-continuity}, there are $X,Y\in\scM_{t,x'}^{\theta}$ such that
\begin{equation}\label{eq:Xs0Ys0}
X(s_{\otimes})=y_{\otimes}\ne Y(s_{\otimes}).
\end{equation}
Moreover, since $(s_{\otimes},y_{\otimes})\in\GS_{\theta}$, we can use Proposition~\ref{prop:basicproperties}$\MK$\eqref{enu:pathsurgery} to modify $X$ on $[T_{*}(t),s_{\otimes}]$ if necessary to ensure that
\begin{equation}\label{eq:winding}
\int_{T_{*}(t)}^{t}(X'-Y')(s)\,\dif s\in\ZZ\setminus\{0\}.
\end{equation}

Let
\[
s_{1}\coloneqq\sup\{s\in[T_{*}(t),t]\st X(s)\ne Y(s)\}
\]
and
\[
s_{2}\coloneqq\sup\{s\in[T_{*}(t),t)\st X(s)=Y(s)\}.
\]
We claim that
\begin{equation}\label{eq:claim}
s_{1}=t\qquad\text{and}\qquad\int_{s_{2}}^{t}(X'-Y')(s)\,\dif s\in\ZZ\setminus\{0\}.
\end{equation}
Indeed, if either of these conditions fail, then by Lemma~\ref{lem:whencanyousplit} and~\eqref{eq:winding}, we can find an $s\in[T_{*}(t),t]$ such that $(s,X(s))=(s,Y(s))\in\ssS_{\theta}\cap\ssN$. But then~\eqref{eq:Xs0Ys0} implies that $(s,X(s))\ne(s_{\otimes},y_{\otimes})$, and since $(s_{\otimes},y_{\otimes})\in\ssN$ this contradicts the assumption that $\mu\in\widetilde{\Omega}_{2}$ (recalling Definition~\ref{def:Omegatilde2}). Thus we can conclude that \eqref{eq:claim} holds. Using Proposition~\ref{prop:XLXXR}, we see that this implies that $(t,x')\in\GS_{\theta}$, and hence the proposition is proved.
\end{proof}

\subsection{Closedness of the global shock set}

It will be helpful to know that the set $\GS$ is closed.

\begin{prop}\label{prop:GSclosed}
Suppose that $\mu\in\widetilde{\Omega}_{1}$. The set $\GS$ is a closed subset of $\TT\times\RR\times\RR$.
\end{prop}

\begin{proof}
Suppose that we have $(\theta_{n},t_{n},x_{n})\to(\theta,t,x)$ and $(\theta_{n},t_{n},x_{n})\in\GS$ for each $n\in\NN$. We claim that $(\theta,t,x)\in\GS$. By Proposition~\ref{prop:main-continuity}, we can find $n\in\NN$, $\zeta,\tau,\eta_{\Square}\in\RR$, $i_{\Square}\in\{1,2\}$, and $Z_{\Square}\in\scM_{t,x}^{\theta}$ (for $\Square\in\LR$) such that $\theta_{n}=\theta+\zeta$, $t_{n}=t+\tau$, $x_{n}=x+\eta_{\Square}$, $\tau>t_{-;i_{\Square}}[Z_{\Square}]-t$, and
\[
X_{\theta_{n},t_{n},x_{n},\Square}=\clT_{\tau,\eta_{\Square}}^{i_{\Square}}Z_{\Square}
\]
for $\Square\in\LR$. In fact, by taking $n$ sufficiently large, it is further possible to ensure that $\eta_{\Left}=\eta_{\Right}\eqqcolon\eta$ (which will be very small), so we have
\[
X_{\theta_{n},t_{n},x_{n},\Square}=\clT_{\tau,\eta}^{i_{\Square}}Z_{\Square}\overset{\eqref{eq:Tidef}}{=}\clT_{t_{-;i_{\Square}}\left[Z_{\Square}\right],\tau,\eta}Z_{\Square}.
\]
Now if $Z_{\Left}'(t-)=Z_{\Right}'(t-)$, then $T_{\vee}(\theta_{n},t_{n},x_{n})\ge t_{-;1}[Z_{\Left}]=t_{-;1}[Z_{\Right}]$ and it is impossible that $(\theta_{n},t_{n},x_{n})\in\GS$. Therefore, $Z_{\Left}'(t-)\ne Z_{\Right}'(t-)$. From this we see that $T_{\vee}(\theta_{n},t_{n},x_{n})=\sup\{s<t\st Z_{\Left}(s)\ne Z_{\Right}(s)\}$ and that
\[
1=\int_{T_{\vee}(\theta_{n},t_{n},x_{n})}^{t_{n}}\left(X_{\theta_{n},t_{n},x_{n},\Left}'-X_{\theta_{n},t_{n},x_{n},\Right}'\right)(s)\,\dif s=\int_{T_{\vee}(\theta_{n},t_{n},x_{n})}^{t}\left(Z_{\Left}'-Z_{\Right}'\right)(s)\,\dif s
\]
which means by Proposition~\ref{prop:XLXXR} that $(\theta,t,x)\in\GS$.
\end{proof}


\begin{rema}
The set $\ssS$ of shocks is \emph{not} a closed subset of $\TT\times\RR\times\RR$. Indeed, each forcing point generates two shocks extending from it forward in time (see Proposition~\ref{prop:advanceint-N} below), but most forcing points are not shocks. But Proposition~\ref{prop:GSclosed} says that a forcing point with a \emph{global} shock coming from it must be a global shock. This is illustrated in Figure~\ref{fig:global_shock_split_picture}. (In fact, the closure of $\ssS$ is $\ssS\cup(\RR\times\ssN)$, but we do not prove this since we do not need this fact in the paper.)
\end{rema}


\subsection{Multiple global shocks}

We have shown in Proposition~\ref{prop:findglobalshock} that there is at least one global shock for every $\theta$ and $t$. In this section we show that there cannot be more than two such global shocks. The key ingredients for this are Lemma~\ref{lem:whencanyousplit}, which says that minimizers can only cross each other at an element of $\ssS_{\theta}\cap\ssN$, and Definition~\ref{def:Omegatilde2}, which says that if $\mu\in\widetilde{\Omega}_{2}$ then there can be at most one element of $\ssS_{\theta}\cap\ssN$. Since, in order to have multiple global shocks at the same time, there must be some crossing of minimizers in order to satisfy the topological conditions in Definition~\ref{def:globalshocks}, these conditions restrict the structure of multiple global shocks.

Define
\[
\GS^{\Left}=\bigl\{(\theta,t,x)\in\GS\st\text{there is an }s\in(T_{\vee}(\theta,t,x),t)\text{ s.t. }(s,X_{\theta,t,x,\Right}(s))\in\GS_{\theta}\cap\ssN\bigr\}
\]
and similarly
\[
\GS^{\Right}=\bigl\{(\theta,t,x)\in\GS\st\text{there is an }s\in(T_{\vee}(\theta,t,x),t)\text{ s.t. }(s,X_{\theta,t,x,\Left}(s))\in\GS_{\theta}\cap\ssN\bigr\}.
\]
Define $\GS_{\theta}^{\Square}=\{(t,x)\st(\theta,t,x)\in\GS^{\Square}\}$ and $\GS_{\theta,t}^{\Square}=
\{x\st(\theta,t,x)\in\GS^{\Square}\}$.
\begin{lemm}\label{lem:GSLnotintersect}
If $\mu\in\widetilde{\Omega}$, then $\GS^{\Left}\cap\GS^{\Right}=\varnothing$.
\end{lemm}

\begin{proof}
If there were a $(\theta,t,x)\in\GS^{\Left}\cap\GS^{\Right},$then since $X_{\theta,t,x,\Left}$ and $X_{\theta,t,x,\Right}$ do not intersect on the time interval $(T_{\vee}(\theta,t,x),t)$ by definition, there would have to be two distinct elements of $\GS_{\theta}\cap\ssN$, contradicting the assumption that $\mu\in\widetilde{\Omega}_{2}$.
\end{proof}


\begin{prop}\label{prop:findsplitpoint}
Suppose that $\mu\in\widetilde{\Omega}$. Let $\theta,t\in\RR$ and $x_{1},x_{2}\in\GS_{\theta,t}$ be such that $x_{1}\ne x_{2}$. Then there are $i_{\Left},i_{\Right}\in\{1,2\}$ such that $\{i_{\Left},i_{\Right}\}=\{1,2\}$ and $x_{i_{\Square}}\in\GS^{\Square}$ for $\Square\in\LR$.
\end{prop}

\begin{proof}
Define $T_{i}=T_{\vee}(\theta,t,x_{i})$, and assume without loss of generality that
\begin{equation}\label{eq:T1geT2}
T_{1}\ge T_{2}.
\end{equation}
Let $\hat{x}_{i}\in\RR$ be such that
\begin{equation}\label{eq:xhatordering}
\hat{x}_{1}\le\hat{x}_{2}\le\hat{x}_{1}+1
\end{equation}
and $\pi(\hat{x}_{i})=x_{i}$ for $i\in\{1,2\}$. (Recall that $\pi\colon\RR\to\TT$ is the projection map.) Now for $i\in\{1,2\}$ and $\Square\in\LR$, define
\[
\hat{X}_{i,\Square}(s)=\hat{x}_{i}+\int_{t}^{s}X_{\theta,t,x_{i},\Square}'(r)\,\dif r+\bfun\{(i,\Square)=(1,\Left)\},
\]
so $\pi\circ\hat{X}_{i,\Square}=X_{\theta,t,x_{i},\Square}$. In particular, we have
\begin{equation}\label{eq:Xhat1soffset}
\hat{X}_{1,\Left}(t)-1=\hat{X}_{1,\Right}(t)=\hat{x}_{1}.
\end{equation}
and
\begin{equation}\label{eq:Xhat2smatch}
\hat{X}_{2,\Left}(t)=\hat{X}_{2,\Right}(t)=\hat{x}_{2}.
\end{equation}
Using these together with~\eqref{eq:xhatordering}, we summarize that
\begin{equation}\label{eq:initialordering}
\underbrace{\hat{X}_{1,\Left}(t)-1=\hat{X}_{1,\Right}(t)}_{=\hat{x}_{1}}\le\underbrace{\hat{X}_{2,\Left}(t)=\hat{X}_{2,\Right}(t)}_{=\hat{x}_{2}}\le\underbrace{\hat{X}_{1,\Left}(t)}_{=\hat{x}_{1}+1}.
\end{equation}
Also, \eqref{eq:Xhat1soffset} along with~\eqref{eq:exactlyoneoff} imply that
\begin{equation}\label{eq:Xhat1satS2}
\hat{X}_{1,\Left}(T_{1})=\hat{X}_{1,\Right}(T_{1}),
\end{equation}
while~\eqref{eq:Xhat2smatch} along with~\eqref{eq:exactlyoneoff} imply that
\begin{equation}\label{eq:Xhat2satS2}
\hat{X}_{2,\Left}(T_{2})=\hat{X}_{2,\Right}(T_{2})-1.
\end{equation}
Finally, we have
\[
\left|\hat{X}_{i,\Left}(s)-\hat{X}_{i,\Right}(s)\right|<1\quad\text{for all }i\in\{1,2\}\text{ and }s\in(S_{i},t)
\]
by the definition of $T_{i}$.

We claim that
\begin{equation}\label{eq:matchcontrad}
\hat{X}_{2,\Left}(T_{1})<\hat{X}_{2,\Right}(T_{1}).
\end{equation}
Indeed, if not, then we must have
\begin{equation}\label{eq:whatifmatch}
\hat{X}_{2,\Left}(T_{1})=\hat{X}_{2,\Right}(T_{1})
\end{equation}
by the definitions, and since we assumed in~\eqref{eq:T1geT2} that $T_{1}\ge T_{2}$, this implies that $T_{1}=T_{2}$, and then~\eqref{eq:whatifmatch} contradicts~\eqref{eq:Xhat2satS2}. Therefore, we conclude that~\eqref{eq:matchcontrad} holds.

Define
\begin{align*}
s_{1}&\coloneqq\min\left\{s\in(T_{1},t)\st\hat{X}_{2,\Left}(s)\ge\hat{X}_{1,\Right}(s)\right\}
\\[-0.2em]
\intertext{and}
s_{2}&\coloneqq\min\left\{s\in(T_{1},t)\st\hat{X}_{2,\Right}(s)\le\hat{X}_{2,\Left}(s)\right\}.
\end{align*}
Now~\eqref{eq:matchcontrad} and~\eqref{eq:initialordering} along with the intermediate value theorem imply that\linebreak $s_{1}\vee s_{2}>T_{1}$. If $s_{1}>S_{1}$, then Lemma~\ref{lem:whencanyousplit} implies that
\[
\hat{X}_{2,\Left}(s_{1})=\hat{X}_{1,\Right}(s_{1})\eqqcolon\hat{y}_{1}
\]
and $(s_{1},\pi(\hat{y}_{1}))\in\ssS_{\theta}\cap\ssN$, while if $s_{2}>S_{1}$, then Lemma~\ref{lem:whencanyousplit} similarly implies that
\[
\hat{X}_{2,\Right}(s_{2})=\hat{X}_{1,\Left}(s_{2})\eqqcolon\hat{y}_{2}
\]
and $(s_{2},\pi(\hat{y}_{2}))\in\ssS_{\theta}\cap\ssN$. Using the fact that $\mu\in\widetilde{\Omega}_{2}$, we conclude that in fact there is an $i\in\{1,2\}$ such that $(s_{i},\pi(\hat{y}_{i}))\in\GS_{\theta}\cap\ssN$, and the proof of Proposition~\ref{prop:findsplitpoint} is complete.
\end{proof}

\begin{coro}\label{cor:atmosttwoshocks}
Suppose that $\mu\in\widetilde{\Omega}$. For any $\theta,t\in\RR$, we have $\#\GS_{\theta,t}\le2$.
\end{coro}

\begin{proof}
Suppose for the sake of contradiction that we have distinct $x_{1},x_{2},x_{3}\in\GS_{\theta,t}$. By Proposition~\ref{prop:findsplitpoint}, we can assume without loss of generality that $x_{1}\in\GS_{\theta,t}^{\Left}$ and $x_{2}\in\GS_{\theta,t}^{\Right}$. But then applying Proposition~\ref{prop:findsplitpoint} twice more we see that $x_{3}\in\GS_{\theta,t}^{\Left}\cap\GS_{\theta,t}^{\Right}$, contradicting
Lemma~\ref{lem:GSLnotintersect}.
\end{proof}

Proposition~\ref{prop:findsplitpoint} and Corollary~\ref{cor:atmosttwoshocks} allow us to make the following definition.

\begin{defi}\label{def:sLsRdefs}
Suppose that $\mu\in\widetilde{\Omega}$. We define functions $\sfs_{\Left},\sfs_{\Right}\colon\RR\times\RR\to\TT$ as follows. For $(\theta,t)\in\RR\times\RR$, if $\GS_{\theta,t}=\{x\}$, then we define $\sfs_{\Left}(\theta,t)=\sfs_{\Right}(\theta,t)=x$. If $\GS_{\theta,t}=\{x_{\Left},x_{\Right}\}$ with $x_{\Square}\in\GS_{\theta,t}^{\Square}$ for $\Square\in\LR$, then we define $\sfs_{\Square}(\theta,t)=x_{\Square}$.
\end{defi}

Then the following proposition follows immediately from the definitions and Proposition~\ref{prop:findsplitpoint}.

\begin{prop}\label{prop:ifnotequalthenwherecollide}
Suppose that $\mu\in\widetilde{\Omega}$. If $\sfs_{\Left}(\theta,t)\ne\sfs_{\Right}(\theta,t)$, then $\theta\in\Theta_{\otimes}$, $s_{\otimes}(t)>T_{\vee}(\theta,t,x)\ge T_{*}(t)$, and
\[
X_{\theta,t,\sfs_{\Left}(\theta,t),\Right}(\sfs_{\otimes}(\theta))=X_{\theta,t,\sfs_{\Right}(\theta,t),\Left}(\sfs_{\otimes}(\theta))=y_{\otimes}(\theta).
\]
\end{prop}


\subsection{Topological classification of minimizers}

For certain purposes, it will be helpful to have a finer-grained classification of minimizers according to the way they wind about the torus $\TT$. Recall the definition \eqref{eq:Mstartx} of $\scM_{*\mid t,x}^{\theta}$.

\begin{defi}\label{def:equiv-relation}
Suppose that $\mu\in\widetilde{\Omega}_{1}$. Let $\theta,t\in\RR$ and $x\in\TT$. We define an equivalence relation $\sim$ on $\scM_{*\mid t,x}^{\theta}$ by $X\sim Y$ whenever
\[
\int_{T_{*}(t)}^{t}(X'-Y')(r)\,\dif r=0.
\]
We define $[X]$ to be the equivalence class of $X$ under $\sim$. We also define, for $\Square\in\LR$,
\[
\scM_{*|t,x,\Square}^{\theta}\coloneqq[X_{\theta,t,x,\Square}].
\]
The partial order $\preccurlyeq$ is defined in~\eqref{eq:partialorderdef} as a partial order on $\scM_{t,x}^{\theta}$. We extend its definition to $\scM_{*\mid t,x}^{\theta}$ in the obvious way. For $\Square\in\LR$, we define $X_{\theta,t,x,\Square,\Left}$ and $X_{\theta,t,x,\Square,\Right}$ to be the minimal and maximal elements, respectively, of $\scM_{*\mid t,x,\Square}^{\theta}$ under $\preccurlyeq$. Finally, for $\Square,\Diamond\in\LR$, we define
\begin{equation}\label{eq:usquarediamonddef}
u_{\theta,\Square,\Diamond}(t,x)=X_{\theta,t,x,\Square,\Diamond}'(t-).
\end{equation}
\end{defi}

We note that
\begin{equation}\label{eq:leftleftrightright}
X_{\theta,t,x,\Square,\Square}=X_{\theta,t,x,\Square}\qquad\text{and}\qquad u_{\theta,\Square,\Square}=u_{\theta,\Square}\qquad\text{for }\Square\in\LR.
\end{equation}

First we show that global shocks lead to multiple equivalence classes.

\begin{prop}\label{prop:GS-equiv-class}
Suppose that $\mu\in\widetilde{\Omega}_{1}$. If $(\theta,t,x)\in\GS$, then $\scM_{*\mid t,x,\Left}^{\theta}\ne\scM_{*\mid t,x,\Right}^{\theta}$.
\end{prop}

\begin{proof}
If $(\theta,t,x)\in\GS$, then we have
\[
\int_{T_{*}(t)}^{t}\left(X_{\theta,t,x,\Left}'-X_{\theta,t,x,\Right}'\right)(s)\,\dif s\ge\int_{T_{\vee}(\theta,t,x)}^{t}\left(X_{\theta,t,x,\Left}'-X_{\theta,t,x,\Right}'\right)(s)\,\dif s>0,
\]
so $X_{\theta,t,x,\Left}\not\sim X_{\theta,t,x,\Right}$ and hence $\scM_{*\mid t,x,\Left}^{\theta}\ne\scM_{*\mid t,x,\Right}^{\theta}$.
\end{proof}

The converse of Proposition~\ref{prop:GS-equiv-class} is false, even if we additionally assume that $(\theta,t,x)\in\ssS$. Indeed, if $\sfs_{\Left}(\theta,t)\ne\sfs_{\Right}(\theta,t)$, then any minimizer that passes through $(s_{\otimes}(\theta),y_{\otimes}(\theta))$ will split into two minimizers with different winding numbers on $[T_{*}(t),t]$. There may be additional non-global shocks at time $t$ whose minimizers pass through $(s_{\otimes}(\theta),y_{\otimes}(\theta))$ after re-merging for the first time. However, we do have the following:

\begin{prop}\label{prop:globalshocknojump}
Suppose that $\mu\in\widetilde{\Omega}_{1}$. If $(\theta,t,x)\not\in\GS$, then $u_{\theta,\Left,\Square}(t,x)=u_{\theta,\Right,\Square}(t,x)$ for each $\Square\in\LR$.
\end{prop}

\begin{proof}
We prove this for $\Square=\Left$; the proof for $\Square=\Right$ is symmetrical. We abbreviate $X_{\Diamond}\coloneqq X_{\theta,t,x,\Diamond,\Left}$, so (recalling~\eqref{eq:usquarediamonddef}) we have $u_{\theta,\Diamond,\Left}(t,x)=X_{\Diamond}'(t-)$. Suppose for the sake of contradiction that $X_{\Left}'(t-)\ne X_{\Right}'(t-)$. Let $T\coloneqq\max\{s<t\st X_{\Left}(s)=X_{\Right}(s)\}$ and let $y\coloneqq X_{\Left}(T)=X_{\Right}(T)$. We note that
\begin{equation}\label{eq:samewindingwhentheyhit}
\int_{T}^{t}\left(X_{\Left}'-X_{\Right}'\right)(s)\,\dif s=0,
\end{equation}
since otherwise we would have $(\theta,t,x)\in\GS$ by Proposition~\ref{prop:XLXXR}. Now we consider two cases (which are in fact symmetrical, but we prove both for clarity).
\begin{case}
Suppose first that $X_{\Left}'(t-)<X_{\Right}'(t-)$.
Define $Y=X_{\Left}\odot_{T}X_{\Right}$. Then we have $Y\ne X_{\Left}$ and $Y\preccurlyeq X_{\Left}$. Also, we have $\int_{T_{*}(t)}^{t}(Y'-X_{\Left}')(s)\,\dif s=0$ by~\eqref{eq:samewindingwhentheyhit}, so $Y\in\scM_{*|t,x,\Left}^{\theta}$. This contradicts the definition of $X_{\Left}$ as the leftmost element of~$\scM_{*|t,x,\Left}^{\theta}$.
\end{case}

\begin{case}
Suppose now that $X_{\Left}'(t-)>X_{\Right}'(t-)$.
Define $Y=X_{\Right}\odot_{T}X_{\Left}$. Then we have $Y\ne X_{\Right}$ and $Y\preccurlyeq X_{\Right}$. Also, we have $\int_{T_{*}(t)}^{t}(Y'-X_{\Right}')(s)\,\dif s=0$ by~\eqref{eq:samewindingwhentheyhit}, so $Y\in\scM_{*|t,x,\Right}^{\theta}$. This contradicts the definition of $X_{\Right}$ as the leftmost element\linebreak  of $\scM_{*|t,x,\Right}^{\theta}$.\qedhere
\end{case}\let\qed\relax
\end{proof}

The following lemma will also be useful.

\begin{lemm}\label{lem:onlytwoguys}
Suppose that $\mu\in\widetilde{\Omega}_{1}$. Let $\theta,t\in\RR$ and $x\in\TT$. If
\begin{equation}\label{eq:twowraps}
\int_{T_{*}(t)}^{t}\left[X_{\theta,t,x,\Left}'(s)-X_{\theta,t,x,\Right}'(s)\right]\,\dif s\ge2,
\end{equation}
then there is some $(s_{\otimes},y_{\otimes})\in[(T_{*}(t),t)\times\TT]\cap\GS_{\theta}\cap\ssN$ such that $X_{\theta,t,x,\Left}(s_{\otimes})=X_{\theta,t,x,\Right}(s_{\otimes})=y_{\otimes}$.
\end{lemm}

\begin{proof}
For~\eqref{eq:twowraps} to hold, the minimizers $X_{\theta,t,x,\Left}$ and $X_{\theta,t,x,\Right}$ must cross in the interval $(T_{*}(t),t)$, so the conclusion is a consequence of Lemma~\ref{lem:whencanyousplit}.
\end{proof}

We now use the equivalence relation introduced in Definition~\ref{def:equiv-relation} to give a further characterization of the left and right global shocks, when they differ.

\begin{prop}\label{prop:findleftrightshocks}
Suppose that $\mu\in\widetilde{\Omega}$ and that $(\theta,t,x)\in\GS$. If
\begin{align}
u_{\theta,\Left,\Right}(t,x) &=u_{\theta,\Right,\Right}(t,x),\nonumber
\\[-0.2em]
\intertext{then}
\sfs_{\Left}(\theta,t) &=x\ne\sfs_{\Right}(\theta,t).\label{eq:shocksdifferent-xleft}
\end{align}
Similarly, if $u_{\theta,\Left,\Left}(t,x)=u_{\theta,\Right,\Left}(t,x)$, then
\begin{equation}\label{eq:shocksdifferent-xright}
\sfs_{\Left}(\theta,t)\ne x=\sfs_{\Right}(\theta,t).
\end{equation}
\end{prop}

\begin{proof}
We will only prove the first statement, as the second is symmetrical. So assume that $u_{\theta,\Left,\Right}(t,x)=u_{\theta,\Right,\Right}(t,x)$. Let $s_{\otimes}=\sup\{r\le t\st X_{\theta,t,x,\Left,\Right}(r)\ne X_{\theta,t,x,\Right,\Right}(r)\}$. Since $u_{\theta,\Left,\Right}(t,x)=u_{\theta,\Right,\Right}(t,x)$, we have $s_{\otimes}<t$. Also, we cannot have $X_{\theta,t,x,\Left,\Right}=X_{\theta,t,x,\Right,\Right}$ since these paths are in different equivalence classes of $\sim$, so we must have $s_{\otimes}>-\infty$. Define
\begin{equation}\label{eq:y0def}
y_{\otimes}=X_{\theta,t,x,\Right}(s_{\otimes}).
\end{equation}
By Lemma~\ref{lem:whencanyousplit}, we have
\begin{equation}\label{eq:s0y0inSN}
(s_{\otimes},y_{\otimes})\in\ssS_{\theta}\cap\ssN.
\end{equation}

We claim that
\begin{equation}\label{eq:s0gt}
s_{\otimes}>T_{\vee}(\theta,t,x).
\end{equation}
Suppose for the sake of contradiction that $s_{\otimes}\le T_{\vee}(\theta,t,x)$. Define the path
\[
Y(s)\coloneqq\left(X_{\theta,t,x,\Left,\Right}\odot_{T_{\vee}(\theta,t,x)}X_{\theta,t,x,\Left}\right)(s)=
\begin{cases}
X_{\theta,t,x,\Left}(s) & \text{if }s\ge T_{\vee}(\theta,t,x);\\
X_{\theta,t,x,\Left,\Right}(s), & \text{if }s\in[T_{*}(t),T_{\vee}(\theta,t,x)].
\end{cases}
\]
This path is continuous since
\[
X_{\theta,t,x,\Left}(T_{\vee}(\theta,t,x))=X_{\theta,t,x,\Right}(T_{\vee}(\theta,t,x))=X_{\theta,t,x,\Left,\Right}(T_{\vee}(\theta,t,x)).
\]
The first identity by the definition~\eqref{eq:Tveedef} of $T_{\vee}$ and the second identity is because, since $s_{\otimes}\le T_{\vee}(\theta,t,x)$, we see that, whenever $s\ge T_{\vee}(\theta,t,x)$, we have $X_{\theta,t,x,\Left,\Right}(s)=X_{\theta,t,x,\Right}(s)$. This last observation moreover allows us to compute
\begin{align*}
\int_{T_{*}(t)}^{t}\left(X_{\theta,t,x,\Left}'-X_{\theta,t,x,\Left,\Right}'\right)(s)\,\dif s & \ge\int_{T_{*}(t)}^{t}\left(Y'-X_{\theta,t,x,\Left,\Right}'\right)(s)\,\dif s\\
& =\int_{T_{\vee}(\theta,t,x)}^{t}\left(X'_{\theta,t,x,\Left}-X'_{\theta,t,x,\Right}\right)(s)\,\dif s\overset{\zcref{eq:shockslopesdiffernt}}{>}0,
\end{align*}
but this contradicts the definition of $X_{\theta,t,x,\Left,\Right}$. Thus we conclude that~\eqref{eq:s0gt} holds.

Now~\eqref{eq:y0def} and~\eqref{eq:s0gt} imply that $X_{\theta,t,x,\Left}(s_{\otimes})\ne y_{\otimes}$. In fact, using Lemma~\ref{lem:whencanyousplit}, \eqref{eq:s0y0inSN}, and the fact that $\mu\in\widetilde{\Omega}_{2}$, we see that
\begin{equation}\label{eq:noty0forabit}
X(s_{\otimes})\ne y_{\otimes}\quad\text{for all }X\in\scM_{t,x}^{\theta}\text{ with }X'(t-)=X_{\theta,t,x,\Left}'(t-).
\end{equation}
Therefore, if we define $\xi\coloneqq\inf\{\eta\ge0\st X_{\theta,t,x+\eta,\Right}(s_{\otimes})\ne y_{\otimes}\}$, then~\eqref{eq:y0def} and~\eqref{eq:Xx+eta} imply that $\xi>0$, while~\eqref{eq:noty0forabit} and~\eqref{eq:Mx-eta} imply that $\xi<1$. Therefore,
\begin{equation}\label{eq:xprimedef}
x'\coloneqq x+\xi\ne x.
\end{equation}
Similar logic implies that
\begin{equation}\label{eq:lefthitsy0}
X_{\theta,t,x',\Left}(s_{\otimes})=y_{\otimes}\ne X_{\theta,t,x',\Right}(s_{\otimes})
\end{equation}
which means (again using Lemma~\ref{lem:whencanyousplit} and the assumption that $\mu\in\widetilde{\Omega}_{2}$) that
\begin{equation}\label{eq:dontintersectbys0}
T_{\vee}(\theta,t,x')<s_{\otimes}.
\end{equation}
Now~\eqref{eq:lefthitsy0} implies that
\begin{equation}\label{eq:follows0y0}
X_{\theta,t,x',\Left}(s)=X_{\theta,s_{\otimes},y_{\otimes},\Left}(s)\quad\text{for all }s\in[T_{*}(t),s_{\otimes}].
\end{equation}
Also, again using Lemma~\ref{lem:whencanyousplit}, in light of the fact that the path $X_{\theta,t,x',\Right}$ does not cross any element of $\ssS_{\theta}\cap\ssN$ by~\eqref{eq:lefthitsy0} and the assumption that $\mu\in\widetilde{\Omega}_{2}$, we see that
\[
X_{\theta,t,x',\Right}(T_{\vee}(\theta,t,x))=X_{\theta,s_{\otimes},y_{\otimes},\Right}(T_{\vee}(\theta,t,x))
\]
and thus that
\[
X_{\theta,t,x',\Right}(s)=X_{\theta,s_{\otimes},y_{\otimes},\Right}(s)\quad\text{for all }s\in[T_{*}(t),T_{\vee}(\theta,t,x)].
\]
Combined with~\eqref{eq:dontintersectbys0}, this implies that
\[
T_{\vee}(\theta,t,x')=T_{\vee}(\theta,s_{\otimes},y_{\otimes}).
\]
See Figure~\ref{fig:global_shock_minimizers_picture} for an illustration.



Now if we define
\begin{equation}\label{eq:Zdef}
Z\coloneqq X_{\theta,s_{\otimes},y_{\otimes},\Right}\odot_{s_{\otimes}}X_{\theta,t,x,\Left},
\end{equation}
then $Z\preccurlyeq X_{\theta,t,x,\Right}$ and so
\begin{align*}
\int_{T_{\vee}(\theta,t,x')}^{t}\left(X_{\theta,t,x',\Left}'-X_{\theta,t,x',\Right}'\right)(r)\,\dif r & \ge\int_{T_{\vee}(\theta,t,x')}^{t}\left(X_{\theta,t,x',\Left}'-Z'\right)(r)\,\dif r\\
& =\int_{T_{\vee}(\theta,t,x')}^{s_{\otimes}}\left(X_{\theta,s_{\otimes},y_{\otimes},\Left}'-X_{\theta,s_{\otimes},y_{\otimes},\Right}'\right)(r)\,\dif r=1,
\end{align*}
with the first identity by~\eqref{eq:follows0y0} and~\eqref{eq:Zdef} and the second by Proposition~\ref{prop:exactlyoneoff}. This implies that $x'\in\GS_{\theta,t}$. But since~\eqref{eq:follows0y0} implies in particular that $X_{\theta,t,x',\Left}(s_{\otimes})=y_{\otimes}$, and $(s_{\otimes},y_{\otimes})\in\ssS_{\theta}\cap\ssN$ and hence in $\GS_{\theta}\cap\ssN$ by Definition~\ref{def:Omegatilde2}, we actually have $x'\in\GS_{\theta,t}^{\Right}$, and hence that $x\in\GS_{\theta,t}^{\Left}$ by Proposition~\ref{prop:findsplitpoint} and Lemma~\ref{lem:GSLnotintersect}. Therefore, we have $\sfs_{\Left}(\theta,t)=x\ne x'=\sfs_{\Right}(\theta,t)$ by Definition~\ref{def:sLsRdefs} and~\eqref{eq:xprimedef}. This completes the proof of Proposition~\ref{prop:findleftrightshocks}.
\end{proof}


\begin{figure}[!hbp]
\begin{centering}
\input{figures/global_shock_minimizers_picture.pgf}
\end{centering}
\caption[Illustration of the situation in the proof of Proposition~\ref{prop:findleftrightshocks}]{Illustration of the situation in the proof of Proposition~\ref{prop:findleftrightshocks}.}
\label{fig:global_shock_minimizers_picture}
\end{figure}


\begin{prop}\label{prop:converse}
Suppose that $\mu\in\widetilde{\Omega}$. Suppose that $\theta,t\in\RR$ are such that $\sfs_{\Left}(\theta,t)\ne\sfs_{\Right}(\theta,t)$. Then we have
\begin{align*}
u_{\theta,\Left,\Right}(t,\sfs_{\Left}(\theta,t)) &=u_{\theta,\Right,\Right}(t,\sfs_{\Left}(\theta,t))
\\[-0.2em]
\intertext{and similarly}
u_{\theta,\Left,\Left}(t,\sfs_{\Right}(\theta,t)) &=u_{\theta,\Right,\Left}(t,\sfs_{\Right}(\theta,t)).
\end{align*}
\end{prop}

\begin{proof}
Again, we just prove the first statement, relying on symmetry to prove the second. It is trivial that $u_{\theta,\Left,\Right}(t,\sfs_{\Left}(\theta,t))\ge u_{\theta,\Right,\Right}(t,\sfs_{\Left}(\theta,t))$, so it suffices to prove that
\begin{equation}\label{eq:goaluLRuRR}
u_{\theta,\Left,\Right}(t,\sfs_{\Left}(\theta,t))\le u_{\theta,\Right,\Right}(t,\sfs_{\Left}(\theta,t)).
\end{equation}
Because of the assumption $\sfs_{\Left}(\theta,t)\ne\sfs_{\Right}(\theta,t)$, we know from Proposition~\ref{prop:findsplitpoint} that there is some $(s_{\otimes},y_{\otimes})\in\GS_{\theta}\cap\ssN$ such that
\[
y_{\otimes}=X_{\theta,t,\sfs_{\Left}(\theta,t),\Right}(s_{\otimes})=X_{\theta,t,\sfs_{\Right}(\theta,t),\Left}(s_{\otimes}).
\]
Now define, for $\Square\in\LR$,
\[
Y_{\Square}\coloneqq X_{\theta,s_{\otimes},y_{\otimes},\Square}\odot_{s_{\otimes}}X_{\theta,t,\sfs_{\Left}(\theta,t),\Right}.
\]
Then we see that $Y_{\Square}\in\scM_{t,\sfs_{\Left}(\theta,t)}^{\theta}$ by Proposition~\ref{prop:basicproperties}$\MK$\eqref{enu:pathsurgery}. We also note that
\begin{equation}\label{eq:YRisXR}
Y_{\Right}=X_{\theta,t,\sfs_{\Left}(\theta,t),\Right},
\end{equation}
while
\begin{equation}\label{eq:diffgtr1}
\begin{split}
\int_{T_{*}(t)}^{t}[Y_{\Left}'(s)-Y_{\Right}'(s)]\,\dif s &=\int_{T_{*}(t)}^{s_{\otimes}}
\left[X'_{\theta,s_{\otimes},y_{\otimes},\Left}(s)-X'_{\theta,s_{\otimes},y_{\otimes},\Left}(s)\right]\,\dif s\\
&\ge\int_{T_{\vee}(\theta,s_{\otimes},y_{\otimes})}^{s_{\otimes}}\left[X'_{\theta,s_{\otimes},y_{\otimes},\Left}(s)-X'_{\theta,s_{\otimes},y_{\otimes},\Left}(s)\right]\,\dif s=1
\end{split}
\end{equation}
since $(s_{\otimes},y_{\otimes})\in\GS_{\theta}$. On the other hand, since $X_{\theta,t,\sfs_{\Left}(\theta,t),\Left}(s_{\otimes})\ne y_{\otimes}$ by Lemma~\ref{lem:GSLnotintersect}, Lemma~\ref{lem:onlytwoguys} tells us that
\[
\int_{T_{*}(t)}^{t}\left[X_{\theta,t,\sfs_{\Left}(\theta,t),\Left}'-X_{\theta,t,\sfs_{\Left}(\theta,t),\Right}'\right](s)\,\dif s=1,
\]
so~\eqref{eq:YRisXR} and~\eqref{eq:diffgtr1} tell us that $Y_{\Left}\in\scM_{*\mid t,\sfs_{\Left}(\theta,t),\Left}^{\theta}$. Therefore, $X_{\theta,t,\sfs_{\Left}(\theta,t),\Left,\Right}$ must lie to the right of $Y_{\Left}$, so
\[
u_{\theta,\Left,\Right}(t,\sfs_{\Left}(\theta,t))=X_{\theta,t,\sfs_{\Left}(\theta,t),\Left,\Right}'(t-)\le Y_{\Left}'(t-)=X_{\theta,t,\sfs_{\Left}(\theta,t),\Right}'(t)=u_{\theta,\Right,\Right}(\theta,\sfs_{\Left}(\theta,t)),
\]
and so~\eqref{eq:goaluLRuRR} is proved.
\end{proof}

\begin{prop}\label{prop:findshocktostay}
Let $\mu\in\widetilde{\Omega}$ and $\theta,t\in\RR$. If $\sfs_{\Left}(\theta,t)\ne\sfs_{\Right}(\theta,t)$, then we have
\begin{equation}\label{eq:rightdiff}
u_{\theta,\Right,\Left}(t,\sfs_{\Right}(\theta,t))>u_{\theta,\Right,\Right}(t,\sfs_{\Right}(\theta,t))
\end{equation}
and
\begin{equation}\label{eq:leftdiff}
u_{\theta,\Left,\Left}(t,\sfs_{\Left}(\theta,t))>u_{\theta,\Left,\Right}(t,\sfs_{\Left}(\theta,t)).
\end{equation}
\end{prop}

\begin{proof}
We limit ourselves to proving~\eqref{eq:rightdiff}, as the proof of~\eqref{eq:leftdiff} is symmetrical. By Proposition~\ref{prop:ifnotequalthenwherecollide}, we have $\theta\in\Theta_{\otimes}$,
\begin{equation}\label{eq:mergeabove}
s_{\otimes}(t)>T_{\vee}(\theta,t,\sfs_{\Right}(\theta,t)),
\end{equation}
and
\[
X_{\theta,t,\sfs_{\Right}(\theta,t),\Left}(s_{\otimes}(\theta))=y_{\otimes}(\theta).
\]
Define
\begin{equation}\label{eq:Xtildedef}
\tilde{X}(s)=X_{\theta,s_{\otimes}(\theta),y_{\otimes}(\theta),\Right}\odot_{s_{\otimes}(\theta)}X_{\theta,t,\sfs_{\Right}(\theta,t),\Left},
\end{equation}
so $\tilde{X}\in\scM_{t,x}^{\theta}$ by Proposition~\ref{prop:basicproperties}$\MK$\eqref{enu:pathsurgery}. Suppose for the sake of contradiction that
\[
\int_{T_{*}(t)}^{t}\left(\tilde{X}'-X_{\theta,t,\sfs_{\Right}(\theta,t),\Right}'\right)(s)\,\dif s>0.
\]
Then we would have
\[
\int_{T_{*}(t)}^{t}\left(X'_{\theta,t,\sfs_{\Right}(\theta,t),\Left}-X_{\theta,t,\sfs_{\Right}(\theta,t),\Right}'\right)(s)\,\dif s>1,
\]
which implies by Lemma~\ref{lem:onlytwoguys} that $X_{\theta,t,\sfs_{\Right}(\theta,t),\Right}(s_{\otimes}(\theta))=y_{\otimes}(\theta)$, contradicting~\eqref{eq:mergeabove}. Therefore, we must have $\int_{T_{*}(t)}^{t}(\tilde{X}'-X_{\theta,t,\sfs_{\Right}(\theta,t),\Right}')(s)\,\dif s=0$ (as it is clearly nonnegative), and hence $\tilde{X}\in\scM_{*|t,x,\Right}^{\theta}$. But this means that
\begin{align*}
u_{\theta,\Right,\Left}(t,\sfs_{\Right}(\theta,t))=X_{\theta,t,\sfs_{\Right}(\theta,t),\Right,\Left}'(t-)&\ge\tilde{X}'(t-)\overset{\eqref{eq:Xtildedef}}{=}X_{\theta,t,\sfs_{\Right}(\theta,t),\Left}'(t-)
\\
& >X_{\theta,t,\sfs_{\Right}(\theta,t),\Right}'(t-) =u_{\theta,\Right,\Right}(t,\sfs_{\Right}(\theta,t)),
\end{align*}
and the proof of Proposition~\ref{prop:findshocktostay} is complete.
\end{proof}

\section{Movement of shocks as time is varied}
\label{sec:Movement-of-shocks}
In this section we study how shocks, and in particular in Section~\ref{subsec:movement-global} global shocks, move as $t$ changes.

\subsection{General case}

\begin{figure}[!htbp]
\begin{centering}
\subfloat[Three minimizers from a single shock at time $t$.\label{fig:shockmoveintime_picture0}
]{\input{figures/shockmoveintime_picture0.pgf}

}\hfill{}\subfloat[Proposition~\ref{prop:advanceint}: just after $t$, there is a single shock. The minimizers follow the left and right time-$t$ minimizers. \label{fig:global_shock_minimizers_picture-1-2}
]{\input{figures/shockmoveintime_picture+1.pgf}

}\hfill{}\subfloat[Proposition~\ref{prop:backwards}: just before $t$, there is one shock for each pair of time-$t$ minimizers with ``adjacent'' slopes.\label{fig:shockmoveintime_picture1}
]{\input{figures/shockmoveintime_picture-1.pgf}}
\end{centering}
\caption{Dynamics of a shock as time changes when there are more than two minimizers coming from a single shock.}
\end{figure}

Proposition~\ref{prop:main-continuity} tells us that, if $(t,x)\in\ssS_{\theta}$ and we perturb $t$ and $x$ by $\tau$ and $\eta$, respectively, then the minimizers will be perturbations of a subset of the original minimizers starting at $t$ and $x$. Thus, we can seek another shock at $(t+\tau,x+\eta)$ by trying to solve for $\eta$ such that there are multiple distinct slopes of minimizers starting at $(t+\tau,x+\eta)$. If there are only two minimizers starting at $(t,x)$, then this procedure is relatively straightforward since in that case those two minimizers must be the ones that we perturb to find minimizers at $(t+\tau,x+\eta)$. In this case we simply recover the usual Rankine--Hugoniot condition (\eqref{eq:RHcondition} below). However, it could also be the case that there are more than two minimizers meeting at the same shock. (See Figure~\ref{fig:shockmoveintime_picture0} for an example.) In this case, when time is moved slightly forward, only the greatest and least slopes of minimizers at time $t$ persist, as illustrated in Figure~\ref{fig:global_shock_minimizers_picture-1-2} and proved in Proposition~\ref{prop:advanceint}. On the other hand, when time is moved slightly backward, if the slopes of minimizers are ordered, then every pair of adjacent slopes corresponds to a separate shock, as illustrated in Figure~\ref{fig:shockmoveintime_picture1} and proved in Proposition~\ref{prop:backwards}.

We begin by considering the perturbative theory for single pairs of minimizers. Recall the definition~\eqref{eq:Tidef} of the path perturbation operator $\clT_{\tau,\eta}^{1}$.

\begin{lemm}\label{lem:differentiate-A}
Suppose that $\mu\in\widetilde{\Omega}_{1}$. Suppose that $\theta\in\RR$ and $(t,x)\in\ssS_{\theta}\setminus\ssN$. If $X_{1},X_{2}\in\scM_{t,x}^{\theta}$ are such that
\begin{equation}\label{eq:derivsnotthesame}
X_{1}'(t-)\ne X_{2}'(t-),
\end{equation}
then there is an $\eps>0$ and a unique function $\sfr_{X_{1},X_{2}}\colon(-\eps,\eps)\to\RR$ such that $\sfr_{X_{1},X_{2}}(0)$ and, for all $\tau\in(-\eps,\eps)$, we have, defining $T_{*}=T_{*}(t-\eps)$, that
\begin{equation}\label{eq:Aequal}
\clA_{\theta,T_{*},t+\tau}\left[\clT_{\tau,\sfr_{X_{1},X_{2}}(\tau)}^{1}X_{i}\right]\text{ is equal for }i\in\{1,2\}.
\end{equation}
Moreover, we have
\begin{equation}\label{eq:RHcondition}
\sfr_{X_{1},X_{2}}'(0)=\frac{1}{2}\left(X_{1}'(t-)+X_{2}'(t-)\right).
\end{equation}
\end{lemm}

\begin{proof}
For any piecewise-linear path $X\colon[T_{*},t]\to\TT$, we first compute from~\eqref{eq:TXderiv} that
\begin{equation}\label{eq:derivdiff}
\left(\clT_{\tau,\eta}^{1}X\right)'(t+\tau-)-X'(t-) =\frac{\eta-\tau X'(t-)}{t+\tau-t_{-;1}[X]}.
\end{equation}
Thus we can compute, using the definition~\eqref{eq:Adef} as well as~\eqref{eq:derivdiff}, that
\begin{multline}\label{eq:AX}
%\begin{aligned}
\clA_{\theta,T_{*},t+\tau}\left[\clT_{\tau,\eta}^{1}X\right]\\
=\frac{1}{2}(t+\tau-t_{-;1}[X])\left(X'(t-)+\frac{\eta-\tau X'(t-)}{t+\tau-t_{-;1}[X]}-\theta\right)^{2}
%\\&\qquad
+\clA_{\theta,T_{*},t_{-;1}[X]}[X].
\end{multline}
Using this with $X=X_{1},X_{2}$ and differentiating with respect to $\eta$, we see that
\begin{equation}\label{eq:etaderiv}
\frac{\dif}{\dif\eta}\left[\clA_{\theta,T_{*},t+\tau}\left[\clT_{\tau,\eta}^{1}X_{2}\right]-\clA_{\theta,T_{*},t+\tau}\left[\clT_{\tau,\eta}^{1}X_{1}\right]\right]_{\eta=\tau=0}=X_{2}'(t-)-X_{1}'(t-).
\end{equation}
Now the assumption~\eqref{eq:derivsnotthesame} let us use the implicit function theorem to find an $\eps>0$ and a unique function $\sfr_{X_{1},X_{2}}\colon(-\eps,\eps)\to\RR$ such that $\sfr_{X_{1},X_{2}}(0)=0$ and~\eqref{eq:Aequal} holds for $\tau\in(-\eps,\eps)$. We can also compute from~\eqref{eq:AX} that
\begin{equation}\label{eq:tauderiv}
\frac{\dif}{\dif\tau}\left[\clA_{\theta,T_{*},t+\tau}\left[\clT_{\tau,\eta}^{1}X_{2}\right]-\clA_{\theta,T_{*},t+\tau}\left[\clT_{\tau,\eta}^{1}X_{1}\right]\right]_{\eta=\tau=0}=-\frac{1}{2}X_{2}'(t-)^{2}+\frac{1}{2}X_{1}'(t)^{2},
\end{equation}
so~\eqref{eq:RHcondition} follows by implicit differentiation along with~\eqref{eq:Aequal}, \eqref{eq:etaderiv}, and~\eqref{eq:tauderiv}.
\end{proof}

\begin{rema}
Generalizing the computations in the last proof, we can compute, for
\emph{any} $\sff\colon(-\eps,\eps)\to\RR$, that
\begin{multline}\label{eq:generalfderiv}
\frac{\dif}{\dif\tau}  \left[\clA_{\theta,T_{*},t+\tau}\left[\clT_{\tau,\sff(\tau)}^{1}X_{2}\right]-\clA_{\theta,T_{*},t+\tau}\left[\clT_{\tau,\sff(\tau)}^{1}X_{1}\right]\right]_{\tau=0}\\
\begin{aligned}
& =-\frac{1}{2}X_{2}'(t-)^{2}+\frac{1}{2}X_{1}'(t-)^{2}+\sff'(0)\left(X_{2}'(t-)-X_{1}'(t-)\right)\\
& =\left(\sff'(0)-\frac{X_{2}'(t)+X_{1}'(t)}{2}\right)\left(X_{2}'(t)-X_{1}'(t)\right).
\end{aligned}
\end{multline}
\end{rema}

Now we will seek to identify the actual minimizers that occur after a perturbation. First we see what happens immediately after time $t$. The following proposition justifies the picture in Figure~\ref{fig:global_shock_minimizers_picture-1-2}.

\begin{prop}\label{prop:advanceint}
Suppose that $\mu\in\widetilde{\Omega}_{1}$. Suppose that $\theta\in\RR$ and $(t,x)\in\ssS_{\theta}\setminus\ssN$. There is an $\eps>0$ and a function $\sfr\colon[0,\eps)\to\RR$ such that $\sfr(0)=0$,
\begin{equation}\label{eq:advanceint-rderiv}
\sfr'(0+)=\frac{1}{2}\left(X_{\theta,t,x,\Left}'(t-)+X_{\theta,t,x,\Right}'(t-)\right),
\end{equation}
and, for all $\tau\in(0,\eps)$, we have
\begin{multline}\label{eq:usesameminimum}
\scM_{t+\tau,x+\sfr(\tau)}^{\theta}\\
=\left\{\clT_{\tau,\sfr(\tau)}^{1}X\st X\in\scM_{t,x}^{\theta}\quad\text{and}\quad X'(t-)\in\{X_{\theta,t,x,\Square}'(t-)\st\Square\in\LR\}\right\}.
\end{multline}
In particular, $(t+\tau,x+\sfr(\tau))\!\in\!\ssS_{\theta}$. Moreover, if $\eta\!\in\!(-\eps,\eps)\setminus\{\sfr(\tau)\}$, then $(t+\tau,x+\eta)\linebreak\notin\ssS_{\theta}$.
\end{prop}

\begin{proof}
We choose $\eps>0$, $T_{*}=T_{*}(t-\eps)$, and $\sfr\coloneqq\sfr_{X_{\theta,t,x,\Left},X_{\theta,t,x,\Right}}$ as defined in Lemma~\ref{lem:differentiate-A}, so~\eqref{eq:advanceint-rderiv} holds by~\eqref{eq:RHcondition}, and we have
\begin{equation}\label{eq:Aequal-apply}
\clA_{\theta,T_{*},t+\tau}\left[\clT_{\tau,\sfr(\tau)}^{1}X_{\theta,t,x,\Left}\right]\overset{\eqref{eq:Aequal}}{=}\clA_{\theta,T_{*},t+\tau}\left[\clT_{\tau,\sfr(\tau)}^{1}X_{\theta,t,x,\Right}\right]\quad\text{for all }\tau\in(-\eps,\eps).
\end{equation}
If $X\in\scM_{t,x}^{\theta}$ is such that
\begin{equation}\label{eq:intermediate}
X_{\theta,t,x,\Left}'(t-)>X'(t-)>X_{\theta,t,x,\Right}'(t-),
\end{equation}
then we have
\begin{multline}\label{eq:derivA}
\frac{\dif}{\dif\tau}  \left[\clA_{\theta,T_{*},t+\tau}\left[\clT_{\tau,\sfr(\tau)}^{1}X\right]-\clA_{\theta,T_{*},t+\tau}\left[\clT_{\tau,\sfr(\tau)}^{1}X_{\theta,t,x,\Square}\right]\right]_{\tau=0}\\
\overset{\eqref{eq:generalfderiv}} = \left(\sfr'(0)-\frac{X'(t-)+X_{\theta,t,x,\Square}'(t)}{2}\right)\left(X'(t)-X_{\theta,t,x,\Square}'(t)\right).
\end{multline}
Now, since $(t,x)\in\ssS_{\theta}$, we have $X_{\theta,t,x,\Left}'(t-)>X_{\theta,t,x,\Right}'(t-)$, and hence must either have
\[
\sfr'(0)<\frac{X'(t-)+X_{\theta,t,x,\Left}'(t-)}{2}\quad\text{or}\quad\sfr'(0)>\frac{X'(t-)+X_{\theta,t,x,\Right}'(t-)}{2}.
\]
In the first case, we see (also using~\eqref{eq:intermediate}) that the right side of~\eqref{eq:derivA} is strictly positive when $\Square=\Left$, and in the second case the right side of~\eqref{eq:derivA} is strictly positive when $\Square=\Right$. From this and Proposition~\ref{prop:main-continuity}, we see that the ``$\subseteq$'' direction of~\eqref{eq:usesameminimum} holds.

On the other hand, if $X'(t)\in\{X_{\theta,t,x,\Square}'(t-)\st\Square\in\LR\}$, then it is clear from the definitions and~\eqref{eq:Aequal-apply} that
\begin{multline}\label{eq:sidessame}
\clA_{\theta,T_{*},t+\tau}\left[\clT_{\tau,\sfr(\tau)}^{1}X\right]
=\clA_{\theta,T_{*},t+\tau}\left[\clT_{\tau,\sfr(\tau)}^{1}X_{\theta,t,x,\Square}\right]\\
\text{for all }\tau\in(-\eps,\eps)\text{ and }\Square\in\LR.
\end{multline}
Using this observation along with Proposition~\ref{prop:main-continuity}, we conclude the equality in~\eqref{eq:usesameminimum}. The last assertion of the proposition then follows from the uniqueness assertion in Lemma~\ref{lem:differentiate-A}.
\end{proof}

We also consider the case when $(t,x)\in\ssN$.
\begin{prop}\label{prop:advanceint-N}
Suppose that $\mu\in\widetilde{\Omega}_{1}$. Suppose that $\theta\in\RR$ and $(t,x)\in\ssN$. There is an $\eps>0$ and two continuous functions $\sfr_{\Left},\sfr_{\Right}\colon[0,\eps)\to\RR$ such that $\sfr_{\Square}(0)=0$ and for all $\tau\in(0,\eps)$, we have $\sfr_{\Left}(\tau)<\sfr_{\Right}(\tau)$ and
\begin{multline}\label{eq:usesameminimums-split}
\scM_{t+\tau,x+\sfr_{\Square}(\tau)}^{\theta}  =\left\{ \clT_{\tau,\sfr_{\Square}(\tau)}^{1}X\st X\in\scM_{t,x}^{\theta}\quad\text{and}\quad X'(t-)=X_{\theta,t,x,\Square}'(t-)\right\} \\
\cup\left\{ \clT_{\tau,\sfr_{\Square}(\tau)}^{2}X\st X\in\scM_{t,x}^{\theta}\right\}.
\end{multline}
In particular, this implies that if $(t,x)\in\GS_{\theta}\cap\ssN$, then $(t+\tau,x+\sfr_{\Square}(\tau))\in\GS_{\theta}$ for all $\tau\in[0,\eps)$ and $\Square\in\LR$.
\end{prop}

\begin{proof}
Let $X_{\Square}\coloneqq X_{\theta,t,x,\Square}$, $m_{\Square}\coloneqq X_{\Square}'(t-)-\theta$, and $s_{\Square}\coloneqq t-t_{-;1}[Y_{\Square}]$. We have, for $\tau>0$ and $\eta\in\RR$, that
\begin{multline*}
\clA_{\theta,T_{*}(t),t+\tau} \left[\clT_{\tau,\eta}^{2}X_{\Square}\right]-\clA_{\theta,T_{*}(t),t+\tau}\left[\clT_{\tau,\eta}^{1}X_{\Square}\right]\\
\begin{aligned}
& =\frac{1}{2}\tau\left(\frac{\eta}{\tau}-\theta\right)^{2}-\mu(\{(t,x)\})+\clA_{\theta,T_{*}(t),t}[Y_{\Square}]-\clA_{\theta,T_{*}(t),t+\tau}\left[\clT_{\tau,\eta}^{1}X_{\Square}\right]\\
\overset{\zcref{eq:AX}} & =\frac{(\eta-\tau\theta)^{2}}{2\tau}-\mu(\{(t,x)\})+\frac{1}{2}s_{\Square}m_{\Square}^{2}-\frac{1}{2}(s_{\Square}+\tau)\left(m_{\Square}+\frac{\eta-\tau(m_{\Square}+\theta)}{s_{\Square}+\tau}\right)^{2}\\
& =\frac{s_{\Square}}{2\tau(s_{\Square}+\tau)}\left((\eta-\tau\theta)^{2}+2\tau m_{\Square}(\eta-\tau\theta)+\tau^{2}m_{\Square}^{2}\right)-\mu(\{(t,x)\})\\
& =\frac{s_{\Square}}{2\tau(s_{\Square}+\tau)}\left(\eta-\tau\theta+\tau m_{\Square}\right)^{2}-\mu(\{(t,x)\}).
\end{aligned}
\end{multline*}
Therefore, if we define
\[
\sfr_{\Square}(\tau)\coloneqq\tau(\theta-m_{\Square})\pm_{\Square}\sqrt{2\tau(1+\tau/s_{\Square})\mu(\{(t,x)\})}
\]
(with $\pm_{\Square}$ as in~\eqref{eq:pmbookkeeping}), then the fact that $\sfr_{\Left}(\tau)<\sfr_{\Right}(\tau)$ is clear, and we moreover have for sufficiently small $\tau>0$ that
\begin{equation}\label{eq:setequal}
\clA_{\theta,T_{*}(t),t+\tau}\left[\clT_{\tau,\sfr_{\Square}(\tau)}^{2}X_{\Square}\right]=\clA_{\theta,T_{*}(t),t+\tau}\left[\clT_{\tau,\sfr_{\Square}(\tau)}^{1}X_{\Square}\right]\quad\text{for }\Square\in\LR.
\end{equation}

We note that
\begin{equation}\label{eq:independentof}
\clA_{\theta,T_{*}(t),t+\tau}\left[\clT_{\tau,\sfr_{\Square}(\tau)}^{2}X\right]\text{ is independent of }\Square\in\LR.
\end{equation}
Also, if $X\in\scM_{t,x}^{\theta}$ and $X'(t-)>X'_{\theta,t,x,\Right}(t-)$, then~\eqref{eq:generalfderiv} implies that, for sufficiently small $\tau>0$, we have
\begin{equation}\label{eq:biggerthanright}
\clA_{\theta,T_{*}(t),t+\tau}\left[\clT_{\tau,\sfr_{\Square}(\tau)}^{1}X\right]>\clA_{\theta,T_{*}(t),t+\tau}\left[\clT_{\tau,\sfr_{\Square}(\tau)}^{1}X_{\theta,t,x,\Right}\right],
\end{equation}
and similarly if $X'(t-)<X'_{\theta,t,x,\Left}(t-)$, then for sufficiently small $\tau>0$ we have
\begin{equation}\label{eq:biggerthanleft}
\clA_{\theta,T_{*}(t),t+\tau}\left[\clT_{\tau,\sfr_{\Square}(\tau)}^{1}X\right]>\clA_{\theta,T_{*}(t),t+\tau}\left[\clT_{\tau,\sfr_{\Square}(\tau)}^{1}X_{\theta,t,x,\Left}\right].
\end{equation}
Now using Proposition~\ref{prop:main-continuity} along with~\eqref{eq:setequal},
\eqref{eq:independentof}, \eqref{eq:biggerthanright}, and \eqref{eq:biggerthanleft}, we conclude \eqref{eq:usesameminimums-split}.

The last claim of the proposition statement is thus clear, since in this case, at least if $\tau$ is sufficiently small, then $\clT_{\tau,\sfr_{\Left}(\tau)}^{1}X_{\theta,t,x,\Left}$ and $\clT_{\tau,\sfr_{\Left}(\tau)}^{2}X_{\theta,t,x,\Right}$ will satisfy the conditions of Definition~\ref{def:globalshocks}, as will $\clT_{\tau,\sfr_{\Right}(\tau)}^{1}X_{\theta,t,x,\Right}$ and $\clT_{\tau,\sfr_{\Right}(\tau)}^{2}X_{\theta,t,x,\Left}$.
\end{proof}

Now we look at what happens just before time $t$. The following proposition justifies the picture in Figure~\ref{fig:shockmoveintime_picture1}.
\begin{prop}\label{prop:backwards}
Suppose that $\mu\in\widetilde{\Omega}_{1}$ and that $\theta\in\RR$, $(t,x)\in\ssS_{\theta}$, and $X_{1},X_{2}\in\scM_{t,x}^{\theta}$ are such that $X_{1}'(t-)>X_{2}'(t-)$ and, for any $X\in\scM_{t,x}^{\theta}$, we have $X'(t-)\not\in(X_{2}'(t-),X_{1}'(t-))$. Then there is an $\eps=\eps(\theta,t,x)>0$ and a function $\sfr_{X_{1},X_{2}}\colon(-\eps,0]\to\TT$ such that $\sfr(0)=0$,
\begin{equation}\label{eq:rderiv-LR-1}
\sfr_{X_{1},X_{2}}'(0-)=\frac{X_{1}'(t-)+X_{2}'(t-)}{2},
\end{equation}
and, for all $\tau\in(-\eps,0)$,
\begin{multline}\label{eq:useminimizers-before}
\scM_{t+\tau,x+\sfr_{X_{1},X_{2}}(\tau)}^{\theta}\\
=\left\{ \clT_{\tau,\sfr_{X_{1},X_{2}}(\tau)}^{1}X\st X\in\scM_{t,x}^{\theta}\quad\text{and}\quad X'(t-)\in\{X_{1}'(t-),X_{2}'(t-)\}\right\}.
\end{multline}
\end{prop}

\begin{proof}
We choose $\eps>0$, $T_{*}=T_{*}(t-\eps)$, and $\sfr\coloneqq\sfr_{X_{1},X_{2}}$ as in Lemma~\ref{lem:differentiate-A}, so~\eqref{eq:rderiv-LR-1} is simply~\eqref{eq:RHcondition}. If $X\in\scM_{t,x}^{\theta}$ is such that $X'(t-)\not\in[X_{2}'(t-),X_{1}'(t-)]$, then we can use~\eqref{eq:generalfderiv} and~\eqref{eq:rderiv-LR-1} to obtain
\begin{multline*}
\frac{\dif}{\dif\tau}\left[\clA_{\theta,T_{*},t+\tau}\left[\clT_{\tau,\sfr(\tau)}^{1}X\right]-\clA_{\theta,T_{*},t+\tau}\left[\clT_{\tau,\sfr(\tau)}^{1}X_{1}\right]\right]_{\tau=0}
\\ =\frac{1}{2}\left(X_{2}'(t-)-X'(t-)\right)\left(X'(t-)-X_{1}'(t-)\right) <0.
\end{multline*}
On the other hand, if $X'(t-)\in\{X_{1}'(t-),X_{2}'(t-)\}$, then it is clear from the definitions and~\eqref{eq:Aequal} that
\[
\clA_{\theta,T_{*},t+\tau}\left[\clT_{\tau,\sfr(\tau)}^{1}X\right]=\clA_{\theta,T_{*},t+\tau}\left[\clT_{\tau,\sfr(\tau)}^{1}X_{1}\right]=\clA_{\theta,T_{*},t+\tau}\left[\clT_{\tau,\sfr(\tau)}^{1}X_{2}\right]\quad\text{for all }\tau\in(-\eps,\eps).
\]
Thus we can take $\eps$ smaller if necessary and apply Proposition~\ref{prop:main-continuity} along with the last two displays to see that~\eqref{eq:useminimizers-before} holds for all $\tau\in(-\eps,0)$.
\end{proof}

\subsection{Global shock movement}
\label{subsec:movement-global}
We now use the results of the previous subsection to describe the movement in time of $\sfs_{\Left}(\theta,\cdot)$ and $\sfs_{\Right}(\theta,\cdot)$. The first part of the following proposition is Theorem~\ref{thm:shockstructure}$\MK$\eqref{enu:shockcontinuity}. At the end of this section, we also prove Theorem~\ref{thm:shockstructure}$\MK$\eqref{enu:splitatforcing}.

\begin{prop}\label{prop:smovement}
Suppose that $\mu\in\widetilde{\Omega}$. Fix $\theta\in\RR$ and $\Square\in\LR$.
\begin{enumerate}
\item \label{enu:scts}
The function $t\mapsto\sfs_{\Square}(\theta,t)$ is continuous in $t$.
\item \label{enu:sderiv}
The function $t\mapsto\sfs_{\Square}(\theta,t)$ is right-differentiable at every $t\in\RR$ such that $(t,\sfs_{\Square}(\theta,t))\not\in\ssN$, and for each such $t$ we have
\begin{equation}\label{eq:sderivintime}
\partial_{t}\sfs_{\Square}(\theta,t+)=\frac{1}{2}\sum_{\Diamond\in\LR}u_{\theta,\Diamond}(t,\sfs_{\Square}(\theta,t)).
\end{equation}
\end{enumerate}
\end{prop}

\begin{proof}
Let $t\in\RR$. We consider two cases.
\newcaseset
\begin{case}
First we consider the case when $\sfs_{\Left}(\theta,t)=\sfs_{\Right}(\theta,t)$. Let $\Square\in\LR$. Suppose for the sake of contradiction that there is a sequence $t_{k}\to t$ such that $\sfs_{\Square}(\theta,t_{k})$ does not converge to $\sfs_{\Square}(\theta,t)$ as $k\to\infty$. By the compactness of $\TT$, we can pass to a subsequence to assume that $x\coloneqq\lim_{k\to\infty}\sfs_{\Square}(\theta,t_{k})\ne\sfs_{\Square}(\theta,t)$ exists. But then Proposition~\ref{prop:GSclosed} implies that $x\in\GS_{\theta,t}$, whereas the assumption that $\sfs_{\Left}(\theta,t)=\sfs_{\Right}(\theta,t)$ implies that $\GS_{\theta,t}=\{\sfs_{\Square}(\theta,t)\}$, a contradiction. Therefore, we must have
\[
\lim_{t'\to t}\sfs_{\Square}(\theta,t')=\sfs_{\Square}(\theta,t),
\]
and hence $\sfs_{\Square}(\theta,\cdot)$ is continuous at $t$. If we now moreover assume that $(t,\sfs_{\Square}(\theta,t))\not\in\ssN$, then Proposition~\ref{prop:advanceint} implies that, if we take $\eps$ and $\sfr$ as in that proposition, we have $\sfs_{\Square}(\theta,t+\tau)=\sfs_{\Square}(\theta,t)+\sfr(\tau)$ for $\tau\in[0,\eps)$, and then~\eqref{eq:advanceint-rderiv} implies~\eqref{eq:sderivintime}.
\end{case}
\begin{case}
Now we consider the case when $\sfs_{\Right}(\theta,t)\ne\sfs_{\Left}(\theta,t)$. Note that this implies in particular that $\theta\in\Theta_{\otimes}$. We proceed in two steps.
\begin{step}
First we address continuity from above. For $\Square\in\LR$, define $\eps_{\Square},\sfr_{\Square}$ as in Proposition~\ref{prop:advanceint} with $x=\sfs_{\Square}(\theta,t)$, and put $\eps=\eps_{\Left}\wedge\eps_{\Right}$, so we know from that proposition that, for all $\tau\in(0,\eps)$, we have
\begin{equation}\label{eq:itsatranslation}
X_{\theta,t+\tau,\sfs_{\Square}(\theta,t)+\sfr_{\Square}(\tau),\Diamond}=\clT_{\tau,\sfr_{\Square}(\tau)}^{1}X_{\theta,t,\sfs_{\Square}(\theta,t),\Diamond}\quad\text{for }\Square,\Diamond\in\LR.
\end{equation}
From this we see that, for $\Square\in\LR$, we have
\[
\int_{T_{\vee}\left(\theta,t+\tau,\sfs_{\Square}(\theta,t)+\sfr_{\Square}(\tau)\right)}^{t+\tau}\left(X_{\theta,t+\tau,\sfs_{\Square}(\theta,t)+\sfr_{\Square}(\tau),\Left}'-X_{\theta,t+\tau,\sfs_{\Square}(\theta,t)+\sfr_{\Square}(\tau),\Right}'\right)(s)\,\dif s\ne0,
\]
so $(t+\tau,\sfs_{\Square}(\theta,t)+\sfr_{\Square}(\tau))\in\GS_{\theta}$. Also, using \eqref{eq:itsatranslation} and Proposition~\ref{prop:ifnotequalthenwherecollide}, we have
\[
X_{\theta,t+\tau,\sfs_{\Right}(\theta,t)+\sfr_{\Right}(\tau),\Left}(s_{\otimes}(\theta))=X_{\theta,t,\sfs_{\Right}(\theta,t),\Left}(s_{\otimes}(\theta))=y_{\otimes}(\theta)
\]
and
\[
X_{\theta,t+\tau,\sfs_{\Left}(\theta,t)+\sfr_{\Left}(\tau),\Right}(s_{\otimes}(\theta))=X_{\theta,t,\sfs_{\Left}(\theta,t),\Right}(s_{\otimes}(\theta))=y_{\otimes}(\theta),
\]
which means that $\sfs_{\Square}(\theta,t+\tau)=\sfs_{\Square}(\theta,t)+\sfr_{\Square}(\tau)$ for each $\Square\in\LR$. Then the continuity from above follows from the continuity of $\sfr_{\Square}(\tau)$. Moreover, if we assume that $(t,\sfs_{\Square}(\theta,t))\not\in\ssN$ (which is in fact guaranteed in this case since $s_{\otimes}(\theta)<t$ by Proposition~\ref{prop:findsplitpoint}), then~\eqref{eq:sderivintime} follows from~\eqref{eq:rderiv-LR-1}.
\end{step}
\begin{step}
Now we address continuity from below. We will prove this for $\sfs_{\Right}$; the proof for $\sfs_{\Left}$ is symmetrical. Continuity from below is somewhat more delicate than continuity from above because multiple shocks may be merging at time $t$, and so we need to figure out which of the merging shocks is the global shock to be followed backward in time. (At most one can be a global shock since we have assumed that $\sfs_{\Left}(\theta,t)\ne\sfs_{\Right}(\theta,t)$, so we cannot be at a point where global shocks merge.)

Let $\tilde{Z}_{\Left}$ be the rightmost element $Z$ of $\scM_{t,\sfs_{\Right}(\theta,t)}^{\theta}$ such that
\[
Z(s_{\otimes}(\theta))=X_{\theta,t,\sfs_{\Right}(\theta,t),\Left}(s_{\otimes}(\theta))=y_{\otimes}(\theta),
\]
and define the concatenated paths
\begin{align}
Z_{\Left}&\coloneqq X_{\theta,s_{\otimes}(\theta),y_{\otimes}(\theta),\Left}\odot_{s_{\otimes}(\theta)}\tilde{Z}_{\Left}\label{eq:ZLdef}
\\
\intertext{and}
Y&\coloneqq X_{\theta,s_{\otimes}(\theta),y_{\otimes}(\theta),\Right}\odot_{s_{\otimes}(\theta)}\tilde{Z}_{\Left},\label{eq:Ydef}
\end{align}
recalling the definition~\eqref{eq:odotdef} of $\odot$. Now let
\[
m_{0}=\max\left\{X'(t-)\st X\in\scM_{t,\sfs_{\Right}(\theta,t)}^{\theta}\quad \text{and}\quad X'(t-)<Z_{\Left}'(t-)\right\},
\]
and let $Z_{\Right}$ be the leftmost element $Z$ of $\scM_{t,\sfs_{\Right}(\theta,t)}^{\theta}$ such that $Z'(t)=m_{0}$. (See Figure~\ref{fig:find_Zs_picture} for an illustration.)

\begin{figure}[!htbp]
\begin{centering}
\input{figures/find_Zs_picture.pgf}
\end{centering}
\caption[Identifying the global shock out of several merging shocks]{Three shocks, including the right global shock $\protect\sfs_{\protect\Right}(\theta,\cdot)$, are merging at time $t$. In backward time, the right global shock continues between $Z_{\protect\Left}$ and $Z_{\protect\Right}$ (drawn with thicker black dashed lines).}
\label{fig:find_Zs_picture}
\end{figure}

From the definitions, we have $Z_{\Left}'(t-)>Z_{\Right}'(t-)$ and, if $X\in\scM_{t,\sfs_{\Right}(\theta,t)}^{\theta}$, then $X'(t-)\not\in(Z_{\Right}'(t-),Z_{\Left}'(t-))$. Thus, Proposition~\ref{prop:backwards} applies, so taking $\eps>0$ and $\sfr=\sfr_{Z_{\Left},Z_{\Right}}$ as in that proposition, we see that for any $\tau\in(-\eps,0)$, we have
\begin{multline}\label{eq:applybackward}
\scM_{t+\tau,\sfs_{\Right}(\theta,t)+\sfr(\tau)}^{\theta}\\
=\left\{ \clT_{\tau,\sfr(\tau)}^{1}X\st X\in\scM_{t,\sfs_{\Right}(\theta,t)}^{\theta}\quad \text{and}\quad X'(t-)\in
\left\{Z_{\Left}'(t-),Z_{\Right}'(t-)\right\}\right\}.
\end{multline}

Now we observe that
\begin{equation}\label{eq:ZRRightmost}
X\preccurlyeq Z_{\Right}\quad\text{for all }X\in\scM_{t,\sfs_{\Right}(\theta,t)}^{\theta}\text{ such that }X'(t-)=Z_{\Right}'(t-),
\end{equation}
which is clear from the definition of $Z_{\Right}$, and also that
\begin{equation}\label{eq:ZLleftmost}
X\succcurlyeq Z_{\Left}\quad\text{for all }X\in\scM_{t,\sfs_{\Right}(\theta,t)}^{\theta}\text{ such that }X'(t-)=Z_{\Left}'(t-).
\end{equation}
To see~\eqref{eq:ZLleftmost}, we note that if there is some $r\le t$ such that $X(r)=Z_{\Left}(r)$ but $X'(r-)\ne Z_{\Left}'(r-)$, then by Lemma~\ref{lem:whencanyousplit} we must have $r=s_{\otimes}(\theta)$ and $X(r)=Z_{\Left}(r)=y_{\otimes}(\theta)$. This means that $X|_{[s_{\otimes}(\theta),t]}=Z_{\Left}|_{[s_{\otimes}(\theta),t]}$. Then the fact that $X\succcurlyeq Z_{\Left}$ follows from the fact that $Z_{\Left}|_{(-\infty,s_{\otimes}(\theta))}$ is the leftmost minimizer from $(s_{\otimes}(\theta),y_{\otimes}(\theta))$ by the definition~\eqref{eq:ZLdef}. Using~\eqref{eq:ZRRightmost} and~\eqref{eq:ZLleftmost} in~\eqref{eq:applybackward}, we see that for any $\tau\in(-\eps,0)$, we have
\begin{equation}\label{eq:leftrightisZ}
X_{\theta,t+\tau,\sfs_{\Right}(\theta,t)+\sfr(\tau),\Square}=\clT_{\tau,\sfr(\tau)}^{1}Z_{\Square}\quad\text{for }\Square\in\LR.
\end{equation}

We also note that
\begin{equation}\label{eq:Zordering}
\int_{r}^{s_{\otimes}(\theta)}Z_{\Right}'(s)\,\dif s\le\int_{r}^{s_{\otimes}(\theta)}Y'(s)\,\dif s\le\int_{r}^{s_{\otimes}(\theta)}Z_{\Left}'(s)\,\dif s\quad\text{for all }r\in[T_{*}(t),t].
\end{equation}
The second inequality is clear from the definitions, and the first is a restatement of \eqref{eq:ZLleftmost}. Using this, we see that if we define
\[
T\coloneqq\sup\{s<t\st Z_{\Left}(s)=Z_{\Right}(s)\},
\]
then
\begin{equation}\label{eq:Zsintegral}
\int_{T}^{t}\left(Z_{\Left}'-Z_{\Right}'\right)(s)\,\dif s=1.
\end{equation}
Indeed, the fact that this integral is positive follows from~\eqref{eq:Zordering}, the definitions \eqref{eq:ZLdef} and~\eqref{eq:Ydef}, and the fact that $(s_{\otimes}(\theta),y_{\otimes}(\theta))\in\GS_{\theta}$, and then the definition of $T$ implies that it is an integer and at most $1$. From~\eqref{eq:Zsintegral} we conclude that, for any $\tau\in(-\eps,0)$, we have
\[
\int_{T}^{t}\left(\left(\clT_{\tau,\sfr(\tau)}^{1}Z_{\Left}\right)'-\left(\clT_{\tau,\sfr(\tau)}^{1}Z_{\Right}\right)'\right)(s)\,\dif s=1
\]
and hence that $(t+\tau,\sfs_{\Right}(\theta,t)+\sfr(\tau))\in\GS_{\theta}$ in light of~\eqref{eq:leftrightisZ}. Moreover, since
\[
X_{\theta,t+\tau,\sfs(\theta,t)+\sfr(\tau),\Left}(s_{\otimes}(\theta))=Z_{\Left}(s_{\otimes}(\theta))=y_{\otimes}(\theta)
\]
by~\eqref{eq:leftrightisZ} and the definition of $Z_{\Left}$, we in fact have $\sfs_{\Right}(t+\tau)=\sfs_{\Right}(\theta,t)+\sfr(\tau)$ for all $\tau\in(-\eps,0)$. Thus the continuity from below of $\sfs_{\Right}$ at $t$ follows from the continuity of~$\sfr$.\qedhere
\end{step}
\end{case}\let\qed\relax
\end{proof}

Now we can prove Theorem~\ref{thm:shockstructure}$\MK$\eqref{enu:splitatforcing}.
\begin{proof}[Proof of Theorem~\ref{thm:shockstructure}$\MK$\eqref{enu:splitatforcing}]
If $\theta\in\RR\setminus\Theta_{\otimes}$, then $\sfs_{\Left}(\theta,t)=\sfs_{\Right}(\theta,t)$ for all $t\in\RR$ by Proposition~\ref{prop:ifnotequalthenwherecollide}.

Thus we now assume that $\theta\in\Theta_{\otimes}$. Let $\sfT\coloneqq\{t\in\RR\st\sfs_{\Left}(\theta,t)\ne\sfs_{\Right}(\theta,t)\}$. By Proposition~\ref{prop:smovement}$\MK$\eqref{enu:scts}, $\sfT$ is an open subset of $\RR$. Also, Proposition~\ref{prop:findsplitpoint} tells us that $\sfT\subseteq(s_{\otimes}(\theta),\infty)$.

We claim that if $t_{k}\in\sfT$ and $t_{k}\downarrow t\not\in\sfT$ as $k\to\infty$, then in fact $t=s_{\otimes}(\theta)$. The continuity of $\sfs_{\Left}(\theta,\cdot)$ and $\sfs_{\Right}(\theta,\cdot)$ implies that $\lim_{k\to\infty}\sfs_{\Left}(\theta,t_{k})=\lim_{k\to\infty}\sfs_{\Right}(\theta,t_{k})\eqqcolon y$. But this is implies that $(t,y)\in\ssN$ by the last statement in
Proposition~\ref{prop:advanceint}, and since the continuity of $\sfs_{\Square}(\theta,\cdot)$ implies that $(t,y)\in\GS_{\theta}$ as well, we have $t=s_{\otimes}(\theta)$ by the assumption that $\mu\in\widetilde{\Omega}_{2}$.

Thus we can conclude that $\sfT=(s_{\otimes}(\theta),s_{\wedge}(\theta))$ for some $s_{\wedge}(\theta)\in[s_{\otimes}(\theta),\infty]$. But we know that $s_{\wedge}(\theta)<\infty$ by Proposition~\ref{prop:allshocksmerge}, and that $s_{\wedge}(\theta)>s_{\otimes}(\theta)$ by Proposition~\ref{prop:advanceint-N}. This completes the proof.
\end{proof}

\section{The dependence of \texorpdfstring{$u$}{u} on \texorpdfstring{$\theta$}{θ}}
\label{sec:udependenceontheta}
Now we begin our study of how the structure of the Burgers flow changes as $\theta$ changes, and ultimately prove
Theorem~\ref{thm:mainthm}. The key observation is the following.

\begin{lemm}
Suppose that $\mu\in\Omega$. If $s<t$, $X\colon[s,t]\to\TT$ and $\theta_{1},\theta_{2}\in\RR$, then
\begin{align}\label{eq:diffAtheta1theta2}
\clA_{\theta_{2},s,t}[X]-\clA_{\theta_{1},s,t}[X] & =(\theta_{1}-\theta_{2})\int_{s}^{t}X'(r)\,\dif r+\frac{1}{2}(t-s)\left(\theta_{2}^{2}-\theta_{1}^{2}\right).
\end{align}
Also, for $\theta\in\RR$, we have
\begin{equation}\label{eq:difAXderiv}
\frac{\dif}{\dif\theta}\clA_{\theta,s,t}[X]=-\int_{s}^{t}X'(r)\,\dif r+\theta(t-s).
\end{equation}
\end{lemm}

\begin{proof}
This follows immediately from the definition~\eqref{eq:Adef}.
\end{proof}

The remainder of this section is divided into two parts. In Section~\ref{subsec:only-change-at-global-shock}, we show that, as $\theta$ changes, $u_{\theta,\Square}(t,x)$ is constant except when $(\theta,t,x)\in\GS$, and at such values of $\theta$ the size of the jump is $u_{\theta,\Left,\Square}(t,x)-u_{\theta,\Right,\Square}(t,x)$. See Proposition~\ref{prop:differenceintermsoftopology} below. In Section~\ref{subsec:Movement-of-shocks}, we relate $u_{\theta,\Left,\Square}(t,x)-u_{\theta,\Right,\Square}(t,x)$ to $\partial_{\theta}\sfs_{\square}(\theta,\cdot)$ (see Proposition~\ref{prop:relatetooneoverderiv} below), and then complete the remaining proofs.


\subsection{Jumps of \texorpdfstring{$u$}{u} occur at global shocks}
\label{subsec:only-change-at-global-shock}

Recall the definition~\eqref{eq:udef} of $u_{\theta,\Square}$, and the definitions~\eqref{eq:Mstartx} of $\scM_{*\mid t,x}^{\theta}$ and~\eqref{eq:usquarediamonddef} of $u_{\theta,\Square,\Diamond}(t,x)$.
\begin{prop}\label{prop:GS'discrete}
Suppose that $\mu\in\widetilde{\Omega}_{1}$. For each fixed $t\in\RR$ and $x\in\TT$, the set
\begin{equation}\label{eq:GS'def}
\GS'_{t,x}\coloneqq\{\theta\in\RR\st(\theta,t,x)\in\GS\}
\end{equation}
is discrete.
\end{prop}

It is important in the statement of Proposition~\ref{prop:GS'discrete} that we do \emph{not} assume that $\mu\in\widetilde{\Omega}_{2}$, because we use Proposition~\ref{prop:GS'discrete} in the proof of Proposition~\ref{prop:omega21} below.
\begin{proof}
To have $(\theta,t,x)\in\GS$, we must have paths $X,Y\colon[T_{*}(t),t]\to\TT$, each consisting of straight line segments connecting points of $\ssN$, such that
\begin{equation}\label{eq:Xprimebigger}
\int_{T_{*}(t)}^{t}X'(s)\,\dif s>\int_{T_{*}(t)}^{t}Y'(s)\,\dif s
\end{equation}
and
\begin{equation}\label{eq:Asequal}
\clA_{\theta,T_{*}(t),t}[X]=\clA_{\theta,T_{*}(t),t}[Y].
\end{equation}
By~\eqref{eq:difAXderiv}, we have
\[
\frac{\dif}{\dif\theta}\left[\clA_{\theta,T_{*}(t),t}[X]-\clA_{\theta,T_{*}(t),t}[Y]\right]=\int_{T_{*}(t)}^{t}(Y'-X')(s)\,\dif s\overset{\eqref{eq:Xprimebigger}}{<}0,
\]
which means that for a given $X$ and $Y$, there is only at most a single value of $\theta$ for which~\eqref{eq:Asequal} can hold. Since, for $\theta$ in any bounded set, there are only finitely many $X$ and $Y$ that can achieve~\eqref{eq:Asequal}, this means that the set $\{\theta\in\RR\st(\theta,t,x)\in\GS\}$ must be discrete.
\end{proof}

\begin{prop}\label{prop:locallyconstant}
Suppose that $\mu\in\widetilde{\Omega}_{1}$. For each $\theta,t\in\RR$, $x\in\TT$, and $\Square\in\LR$, there is an $\eps=\eps(\theta,t,x,\Square)>0$ such that, whenever $\theta-\eps<\theta_{-}<\theta<\theta_{+}<\theta+\eps$, we have
\begin{equation}\label{eq:howdominimizerschange}
\scM_{*\mid t,x}^{\theta'}=
\begin{cases}
\scM_{*\mid t,x,\Right}^{\theta} & \text{if }\theta'\in(\theta-\eps,\theta);\\
\scM_{*\mid t,x,\Left}^{\theta} & \text{if }\theta'\in(\theta,\theta+\eps),
\end{cases}
\end{equation}
and hence that
\begin{equation}\label{eq:jumpsize-mo}
u_{\theta_{+},\Square}(t,x)=u_{\theta,\Left,\Square}(t,x)\quad\text{and}\quad u_{\theta_{-},\Square}(t,x)=u_{\theta,\Right,\Square}(t,x).
\end{equation}
\end{prop}

\begin{proof}
It is a consequence of Proposition~\ref{prop:main-continuity} that there is an $\eps>0$ such that if $|\theta'-\theta|<\eps$, then
\[
\scM_{*\mid t,x}^{\theta'}=\left\{ X\in\scM_{*\mid t,x}^{\theta}\st\clA_{\theta',T_{*}(t),t}[X]\le\clA_{\theta',T_{*}(t),t}[Y]\quad\text{for all }Y\in\scM_{*\mid t,x}^{\theta}\right\}.
\]
We have by~\eqref{eq:diffAtheta1theta2} that, for $X,Y\in\scM_{*\mid t,x}^{\theta'}$,
\begin{multline*}
\frac{\clA_{\theta',T_{*}(t),t}[X]-\clA_{\theta',T_{*}(t),t}[Y]}{\theta'-\theta}\\
\begin{aligned}
& =\frac{\clA_{\theta',T_{*}(t),t}[X]-\clA_{\theta,T_{*}(t),t}[X]-\left(\clA_{\theta',T_{*}(t),t}[Y]-\clA_{\theta,T_{*}(t),t}[Y]\right)}{\theta'-\theta}\\
& =-\int_{T_{*}(t)}^{t}(X'-Y')(s)\,\dif s\mathrel{
\begin{cases}
=0 & \text{if and only if }X\sim Y;\\
\ge0 & \text{if }X\in\scM_{*\mid t,x,\Right}^{\theta};\\
\le0 & \text{if }X\in\scM_{*\mid t,x,\Left}^{\theta},
\end{cases}}
\end{aligned}
\end{multline*}
where we recalled Definition~\ref{def:equiv-relation} of $\sim$ and $\scM_{*\mid t,x,\Square}^{\theta}$. Combining the last two displays, we conclude that~\eqref{eq:howdominimizerschange} holds, and then~\eqref{eq:jumpsize-mo} follows from~\eqref{eq:howdominimizerschange} and the definitions.
\end{proof}

The next statement gives us our first formula for $u_{\theta_{2},\Square}(t,x)-u_{\theta_{1},\Square}(t,x)$, which is the subject of Theorem\ref{thm:mainthm}. Recall that $\llbracket\theta_{1},\theta_{2}\rrbracket_{\Square}$ was defined in~\eqref{eq:pmbookkeeping}.

\begin{prop}\label{prop:differenceintermsoftopology}
Suppose that $\mu\in\widetilde{\Omega}_{1}$. Fix $t\in\RR$ and $x\in\TT$. We have, for $\Square\in\LR$, that
\begin{equation}\label{eq:differenceintermsoftopology}
u_{\theta_{2},\Square}(t,x)-u_{\theta_{1},\Square}(t,x)=\sum_{\theta\in\llbracket\theta_{1},\theta_{2}\rrbracket_{\Square}\cap\GS'_{t,x}}\bigl(u_{\theta,\Left,\Square}(t,x)-u_{\theta,\Right,\Square}(t,x)\bigr).
\end{equation}
\end{prop}

\begin{proof}
First we note that if $\theta\in\RR\setminus\GS'_{t,x}$, then $u_{\theta,\Left,\Diamond}(t,x)=u_{\theta,\Right,\Diamond}(t,x)$ for $\Diamond\in\LR$ by Proposition~\ref{prop:globalshocknojump}, which means by Proposition~\ref{prop:locallyconstant} that there is some $\eps>0$ such that if $|\theta'-\theta|<\eps$, then
\[
u_{\theta',\Square}(t,x)=u_{\theta,\Left,\Square}(t,x)=u_{\theta,\Right,\Square}(t,x)=u_{\theta,\Square}(t,x)
\]
(with the last identity by~\eqref{eq:leftleftrightright}). This, for each $\Square\in\LR$, the map $\theta\mapsto u_{\theta,\Square}(t,x)$ is constant on each connected component of $ \RR\setminus\GS'_{t,x}$. Moreover, for each $\theta\in\GS'_{t,x}$, we have again by Proposition~\ref{prop:locallyconstant} that
\[
\lim_{\theta'\downarrow\theta}u_{\theta',\Square}(t,x)-\lim_{\theta'\uparrow\theta}u_{\theta',\Square}(t,x)=u_{\theta,\Left,\Square}(t,x)-u_{\theta,\Right,\Square}(t,x).
\]
So, for any $\Square\in\LR$, we have
\begin{multline*}
u_{\theta_{2},\Square}(t,x)-u_{\theta_{1},\Square}(t,x) =u_{\theta_{2},\Square}(t,x)-u_{\theta_{2},\Right,\Square}(t,x)\\
\begin{aligned}
&+\sum_{\theta\in(\theta_{1},\theta_{2})\cap\GS'_{t,x}}\bigl(u_{\theta,\Left,\Square}(t,x)-u_{\theta,\Right,\Square}(t,x)\bigr)\\
&+u_{\theta_{1},\Left,\Square}(t,x)-u_{\theta_{1},\Square}(t,x).
\end{aligned}
\end{multline*}
We observe that the sum comprises only finitely many terms by Proposition~\ref{prop:GS'discrete}. For $\Square=\Right$, we note that $u_{\theta,\Right}(t,x)=u_{\theta,\Right,\Right}(t,x)$ by~\eqref{eq:leftleftrightright}, so we obtain
\[
u_{\theta_{2},\Right}(t,x)-u_{\theta_{1},\Right}(t,x)=\sum_{\theta\in[\theta_{1},\theta_{2})\cap\GS'_{t,x}}\bigl(u_{\theta,\Left,\Square}(t,x)-u_{\theta,\Right,\Square}(t,x)\bigr).
\]
Similarly, for $\Square=\Left$, we have $u_{\theta,\Left}(t,x)=u_{\theta,\Left,\Left}(t,x)$ by~\eqref{eq:leftleftrightright}, so we obtain
\[
u_{\theta_{2},\Left}(t,x)-u_{\theta_{1},\Left}(t,x)=\sum_{\theta\in(\theta_{1},\theta_{2}]\cap\GS'_{t,x}}\bigl(u_{\theta,\Left,\Left}(t,x)-u_{\theta,\Right,\Left}(t,x)\bigr).
\]
The conclusion~\eqref{eq:differenceintermsoftopology} is a summary of the last two displays using the notation \eqref{eq:pmbookkeeping}.
\end{proof}


\begin{prop}\label{prop:shocksdiffdiscrete}
Suppose that $\mu\in\widetilde{\Omega}$. For any fixed $t\in\RR$, the set $\{\theta\in\RR\st\sfs_{\Left}(\theta,t)\ne\sfs_{\Right}(\theta,t)\}$ is discrete.
\end{prop}

\begin{proof}
Let $\ssN_{*}(t)\coloneqq([T_{*}(t),t]\times\TT)\cap\ssN$. We know from Proposition~\ref{prop:findsplitpoint} and~\eqref{eq:Tveedef} that
\begin{equation}\label{eq:implication}
\text{if }\sfs_{\Left}(\theta,t)\ne\sfs_{\Right}(\theta,t),\text{ then }\ssN_{*}(t)\cap\GS_{\theta}\ne\varnothing.
\end{equation}
By Proposition~\ref{prop:GS'discrete}, we see that for each $(s,y)\in\ssN_{*}(t)$, the set $\{\theta\in\RR\st(s,y)\in\GS_{\theta}\}$ is discrete. But since $\ssN$ is discrete, the set $\ssN_{*}(t)$ is finite, and so
\[
\{\theta\in\RR\st\ssN_{*}(t)\cap\GS_{\theta}\ne\varnothing\}=\bigcup_{(s,y)\in\ssN(t)}\{\theta\in\RR\st(s,y)\in\GS_{\theta}\}
\]
is also discrete. Hence, using~\eqref{eq:implication}, we conclude that $\{\theta\in\RR\st\sfs_{\Left}(\theta,t)\ne\sfs_{\Right}(\theta,t)\}$ is discrete, as claimed.
\end{proof}

We can also complete the proof of Theorem~\ref{thm:mainthm}$\MK$\eqref{enu:switching}.
\begin{proof}[{Proof of Theorem~\ref{thm:mainthm}$\MK$\eqref{enu:switching}}]
We will just prove~\eqref{eq:rightstillshock}, as the proof of~\eqref{eq:leftstillshock} is symmetrical. By Proposition~\ref{prop:locallyconstant}, there is an $\eps>0$ such that, if $\theta'\in(\theta-\eps,\theta]$, then $\scM_{*\mid t,\sfs_{\Right}(\theta,t)}^{\theta'}=\scM_{*\mid t,\sfs_{\Right}(\theta,t),\Right}^{\theta}$, and so in particular
\begin{equation}\label{eq:theta'equalstheta}
u_{\theta',\Square}(t,\sfs_{\Right}(\theta,t))=u_{\theta,\Right,\Square}(t,\sfs_{\Right}(\theta,t))\quad\text{for }\Square\in\LR.
\end{equation}
By Proposition~\ref{prop:findshocktostay}, we have
\[
u_{\theta,\Right,\Left}(t,\sfs_{\Right}(\theta,t))>u_{\theta,\Right,\Right}(t,\sfs_{\Right}(\theta,t)),
\]
so~\eqref{eq:theta'equalstheta} implies that
\[
u_{\theta',\Left}(t,\sfs_{\Right}(\theta,t))>u_{\theta',\Right}(t,\sfs_{\Right}(\theta,t)),
\]
so $(\theta',t,\sfs_{\Right}(\theta,t))\in\ssS$ by the definition~\eqref{eq:Sdef}.
\end{proof}

\subsection{Movement of shocks as \texorpdfstring{$\theta$}{θ} is varied}
\label{subsec:Movement-of-shocks}

In this subsection, we will prove the following proposition, which we will combine with Proposition~\ref{prop:differenceintermsoftopology} to complete the proof of Theorem~\ref{thm:mainthm}$\MK$\eqref{enu:derivintheta}.

\begin{prop}\label{prop:relatetooneoverderiv}
Suppose that $\mu\in\widetilde{\Omega}$ and fix $\theta,t\in\RR$. We have
\begin{equation}\label{eq:rightside}
u_{\theta,\Left,\Right}(t,\sfs_{\Right}(\theta,t))-u_{\theta,\Right,\Right}(t,\sfs_{\Right}(\theta,t))=\frac{1}{\partial_{\theta}\sfs_{\Right}(\theta+,t)}
\end{equation}
and
\begin{equation}\label{eq:leftside}
u_{\theta,\Left,\Left}(t,\sfs_{\Left}(\theta,t))-u_{\theta,\Right,\Left}(t,\sfs_{\Left}(\theta,t))=\frac{1}{\partial_{\theta}\sfs_{\Left}(\theta-,t)}.
\end{equation}
In particular, the one-sided derivatives on the right sides of~\eqref{eq:rightside} and~\eqref{eq:leftside} both exist.
\end{prop}

In this section we will only prove~\eqref{eq:rightside}, as the proof of~\eqref{eq:leftside} is symmetrical. To simplify notation, we make the abbreviation
\[
\clT_{\eta}\coloneqq\clT_{0,\eta}^{1},
\]
with $\clT_{0,\eta}^{1}$ defined as in~\eqref{eq:Tidef}, and we also recall the definition~\eqref{eq:tminus1} of $t_{-;1}[X]$.

\begin{lemm}
Suppose that $\mu\in\widetilde{\Omega}$. Let $\theta,t\in\RR$ and $x\in\TT$. There is an $\eps=\eps(\theta,t,x)>0$ such that, if $0<|\zeta|,|\beta\zeta|<\eps$ and $X\in\scM_{*\mid t,x}^{\theta}$, then
\begin{multline}\label{eq:Adifflem}
\clA_{\theta+\zeta,T_{*}(t),t}  \left[\clT_{\beta\zeta}X\right]-\clA_{\theta,T_{*}(t),t}[X]-\frac{1}{2}\zeta^{2}(t-T_{*}(t)-2\beta)+\mu(\{(t,x)\})\\
 =\zeta\left[\beta\left(X'(t-)-\theta\right)-\int_{T_{*}(t)}^{t}\left(X'(s)-\theta\right)\,\dif s\right]+\frac{\beta^{2}\zeta^{2}}{2(t-t_{-;1}[X])}.
\end{multline}
\end{lemm}

An immediate consequence of~\eqref{eq:Adifflem} is that if $X_{1},X_{2}\in\scM_{*\mid t,x}^{\theta}$, then
\begin{multline}\label{eq:Adiff-double}
\clA_{\theta+\zeta,T_{*}(t),t}  \left[\clT_{\beta\zeta}X_{2}\right]-\clA_{\theta+\zeta,T_{*}(t),t}[\clT_{\beta\zeta}X_{1}]\\
=\zeta\left[\beta\left(X_{2}'(t-)-X'_{1}(t-)\right)-\int_{T_{*}(t)}^{t}\left(X_{2}'(s)-X_{1}'(s)\right)\,\dif s\right]\\
+\frac{1}{2}\zeta^{2}\beta^{2}\left(\frac{1}{t-t_{-;1}[X_{2}]}-\frac{1}{t-t_{-;1}[X_{1}]}\right).
\end{multline}

\begin{prop}\label{prop:findYs}
Suppose that $\mu\in\widetilde{\Omega}$. Let $\theta,t\in\RR$ and $x\in\TT$, and let $Y_{1}\coloneqq X_{\theta,t,x,\Left,\Right}$ and $\tilde{Y}_{2}\coloneqq X_{\theta,t,x,\Right,\Right}$. Let $s_{0}\coloneqq\sup\{s<t\st Y(s)=\tilde{Y}_{2}(s)\}$ and $Y_{2}\coloneqq Y_{1}\odot_{s_{0}}\tilde{Y}_{2}$. Then we have
\begin{equation}\label{eq:Y1Y2diff}
\int_{T_{*}(t)}^{t}\left(Y_{1}'-Y_{2}'\right)(s)\,\dif s=\int_{s_{0}}^{t}\left(Y_{1}'-Y_{2}'\right)(s)\,\dif s\in\{0,1\}
\end{equation}
and, for any $Z\in\scM_{*\mid t,x}^{\theta}$, there is an $i\in\{1,2\}$ such that
\begin{equation}\label{eq:ourconstruct}
Z'(t-)-Y_{i}'(t-)\ge0\quad\text{and}\quad\int_{T_{*}(t)}^{t}\left(Z'-Y_{i}'\right)(s)\,\dif s\le0.
\end{equation}
\end{prop}

\begin{proof}
It is clear from the definitions that the first two expressions in~\eqref{eq:Y1Y2diff} are equal, and moreover that they are in $\{-1,0,1\}$. Moreover, if they were equal to $-1$, then we would have
\[
-1=\int_{T_{*}(t)}^{t}\left(Y_{1}'-Y_{2}'\right)(s)\,\dif s=\int_{T_{*}(t)}^{t}\left(X_{\theta,t,x,\Left}'-X'_{\theta,t,x,\Right}\right)(s)\,\dif s\ge0,
\]
a contradiction. Thus we conclude that~\eqref{eq:Y1Y2diff} holds.

If $Z\in\scM_{*\mid t,x,\Left}^{\theta}$, then $\int_{T_{*}(t)}^{t}(Z'-Y_{1}')(s)\,\dif s=0$ and $Z'(t-)\ge Y_{1}'(t-)$, so~\eqref{eq:ourconstruct} holds with $i=1$ in this case. Thus we now assume that $Z\not\in\scM_{*\mid t,x,\Left}^{\theta}$. This implies in particular that
\[
\int_{T_{*}(t)}^{t}\left(Z'-Y_{1}'\right)(s)\,\dif s\le-1.
\]
Combining this with~\eqref{eq:Y1Y2diff}, we see that
\[
\int_{T_{*}(t)}^{t}\left(Z'-Y_{2}'\right)(s)\,\dif s\le0.
\]
Since it is also clear that $Z'(t-)\ge\tilde{Y}_{2}'(t-)=Y_{2}'(t-)$, we see that~\eqref{eq:ourconstruct} holds with $i=2$.
\end{proof}

\begin{prop}
Suppose that $\mu\in\widetilde{\Omega}_{1}$. Suppose that $(\theta,t,x)\in\GS$ and
\begin{equation}\label{eq:canmoveshock}
u_{\theta,\Left,\Right}(t,x)>u_{\theta,\Right,\Right}(t,x).
\end{equation}
Then there is an $\eps=\eps(\theta,t,x)>0$ and a function $\sfr\colon[0,\eps)\to\TT$ such that $\sfr(0)=x$, $(\theta+\zeta,t,\sfr(\zeta))\in\GS$ for each $\zeta\in[0,\eps)$, and
\begin{equation}\label{eq:rderiv}
\sfr'(0+)=\frac{1}{u_{\theta,\Left,\Right}(t,x)-u_{\theta,\Right,\Right}(t,x)}.
\end{equation}
\end{prop}

\goodbreak
\begin{proof}
The proof proceeds in several steps.
\newstepset
\begin{step}
By Proposition~\ref{prop:main-continuity}, there is an $\eps>0$ such that, if $|\beta\zeta|<\eps$ and $|\theta'-\theta|<\eps$, then $x+\beta\zeta\in\GS_{\theta',t}$ if there exist $Y_{1},Y_{2}\in\scM_{*\mid t,x}^{\theta}$ such that
\begin{equation}\label{eq:globalshock}
(\clT_{\beta\zeta}Y_{2})'(t-)>(\clT_{\beta\zeta}Y_{1})'(t-),\quad\int_{T_{*}(t)}^{t}\left[(\clT_{\beta\zeta}Y_{2})'-(\clT_{\beta\zeta}Y_{1})'\right](s)\,\dif s>0,
\end{equation}
and, for all $Z\in\scM_{*\mid t,x}^{\theta}$, we have
\begin{equation}\label{eq:globalshockm}
\clA_{\theta',T_{*}(t),t}\left[\clT_{\beta\zeta}Z\right]\ge\clA_{\theta',T_{*}(t),t}\left[\clT_{\beta\zeta}Y_{1}\right]=\clA_{\theta',T_{*}(t),t}\left[\clT_{\beta\zeta}Y_{2}\right].
\end{equation}

We also assume that $\eps$ is chosen small enough that, if $X\in\scM_{*\mid t,x}^{\theta}$ and $|\beta\zeta|<\eps$, then there are no points of $\ssN$ on $\clT_{\beta\zeta}X$ that are not also on $X$; this is possible by the discreteness of $\ssN$.
\end{step}
\begin{step}
We select $Y_{1},\tilde{Y}_{2},Y_{2}$ as in the statement of Proposition~\ref{prop:findYs}. We note that
\begin{equation}\label{eq:Yderivorder}
Y_{2}'(t-)<Y_{1}'(t-)
\end{equation}
by the assumption~\eqref{eq:canmoveshock}. We claim that
\begin{equation}\label{eq:diffisL}
\int_{T_{*}(t)}^{t}(Y_{1}'-Y_{2}')(s)\,\dif s=1.
\end{equation}
In light of \zcref{eq:Y1Y2diff}, it suffices to show that
\[
\int_{T_{*}(t)}^{t}(Y_{1}'-Y_{2}')(s)\,\dif s\ne0,
\]
but this is true because otherwise we would have $Y_{2}\in\scM_{*\mid t,x,\Left}^{\theta}$ and then~\eqref{eq:Yderivorder} would contradict the definition of $Y_{1}$ as the rightmost element of $\scM_{*\mid t,x,\Left}^{\theta}$. Now~\eqref{eq:Yderivorder} and~\eqref{eq:diffisL} imply that~\eqref{eq:globalshock} holds.
\end{step}
\begin{step}
By~\eqref{eq:Adiff-double} and~\eqref{eq:diffisL}, we have
\begin{multline}\label{eq:Adiff-double-1}
\clA_{\theta+\zeta,T_{*}(t),t}\left[\clT_{\beta\zeta}Y_{2}\right]-\clA_{\theta+\zeta,T_{*}(t),t}\left[\clT_{\beta\zeta}Y_{1}\right]\\
=\zeta\left(\beta\left(Y_{2}'(t-)-Y'_{1}(t-)\right)-1+\beta^{2}Q\zeta/2\right),
\end{multline}
where we have defined
\[
Q\coloneqq\frac{1}{t-t_{-;1}[Y_{2}]}-\frac{1}{t-t_{-;1}[Y_{1}]}.
\]
Now define, as long as $\zeta$ is sufficiently small,
\begin{equation}\label{eq:betazetadef}
\beta_{\zeta}\coloneqq
\begin{cases}
\left(Y_{2}'(t-)-Y_{1}'(t-)\right)\cdot\frac{\left(1+2Q\zeta\left(Y_{2}'(t-)-Y_{1}'(t-)\right)^{-2}\right)^{1/2}-1}{Q\zeta}, & \text{if }Q\zeta\ne0;\\
\left(Y_{2}'(t-)-Y_{1}'(t-)\right)^{-1}, & \text{if }Q\zeta=0.
\end{cases}
\end{equation}
Using~\eqref{eq:betazetadef} in~\eqref{eq:Adiff-double-1}, we see (for $\zeta$ small enough that~\eqref{eq:betazetadef} is well-defined) that
\[
\clA_{\theta+\zeta,T_{*}(t),t}\left[\clT_{\beta_{\zeta}\zeta}Y_{2}\right]=\clA_{\theta+\zeta,T_{*}(t),t}\left[\clT_{\beta_{\zeta}\zeta}Y_{1}\right],
\]
which verifies the identity in~\eqref{eq:globalshockm}. We also observe that, for small $\zeta$, we have the Taylor expansion
\begin{equation}\label{eq:betazetaTaylor}
\beta_{\zeta}=\frac{1}{Y_{2}'(t-)-Y_{1}'(t-)}-\frac{Q\zeta}{2\left(Y_{2}'(t-)-Y_{1}'(t-)\right)^{3}}+O(\zeta^{2}).
\end{equation}
\end{step}
\begin{step}
Now we want to verify the inequality in~\eqref{eq:globalshockm}. Let $Z\in\scM_{*\mid t,x}^{\theta}$. By Proposition~\ref{prop:findYs}, we can find an $i\in\{1,2\}$ such that~\eqref{eq:ourconstruct} holds. Using~\eqref{eq:Adiff-double}, we can compute, for $\zeta\ge0$, that
\begin{multline}\label{eq:derivZY}
\clA_{\theta+\zeta,T_{*}(t),t}\left[\clT_{\beta_{\eta}\eta}Z\right]-\clA_{\theta+\zeta,T_{*}(t),t}\left[\clT_{\beta_{\eta}\eta}Y_{i}\right]\\
\begin{aligned}
&=\zeta\beta_{\zeta}\left(Z'(t-)-Y_{i}'(t-)\right)-\zeta\int_{T_{*}(t)}^{t}\left(Z'-Y_{i}'\right)(s)\,\dif s\\
&\quad+\frac{1}{2}\zeta^{2}\beta_{\zeta}^{2}\left(\frac{1}{t-t_{-;1}[Z]}-\frac{1}{t-t_{-;1}[Y_{i}]}\right)\\
&\ge\zeta\beta_{\zeta}\left(Z'(t-)-Y_{i}'(t-)\right)+\frac{1}{2}\zeta^{2}\beta_{\zeta}^{2}\left(\frac{1}{t-t_{-;1}[Z]}-\frac{1}{t-t_{-;1}[Y_{i}]}\right),
\end{aligned}
\end{multline}
where the inequality is by the second inequality in~\eqref{eq:ourconstruct}. Now if the first inequality in~\eqref{eq:ourconstruct} is strict, then the right side of~\eqref{eq:derivZY} is (strictly) positive for sufficiently small $\zeta>0$ (and is zero for $\zeta=0$). On the other hand, if $Z'(t-)=Y_{i}'(t-)$, then $t_{-;1}[Z]=t_{-;1}[Y_{i}]$ as well (recalling the definition~\eqref{eq:tminus1}), and so in this case the right side of~\eqref{eq:derivZY} is zero. Therefore, the inequality in~\eqref{eq:globalshockm} is verified.
\end{step}
\begin{step}
We now define $\sfr(\zeta)=x+\beta_{\zeta}\zeta$. For $\zeta\ge0$, the conditions~\eqref{eq:globalshock} and \eqref{eq:globalshockm} have been verified and so we have $(\theta+\zeta,t,\sfr(\zeta))\in\GS$. The derivative~\eqref{eq:rderiv} follows from~\eqref{eq:betazetaTaylor}, and so the proof is complete.\qedhere
\end{step}\let\qed\relax
\end{proof}

We can now complete the proof of Proposition~\ref{prop:relatetooneoverderiv} and in fact simultaneously prove~\eqref{eq:sLsRareleftrightlimits}.

\begin{proof}[Proof of Proposition~\ref{prop:relatetooneoverderiv} and~\eqref{eq:sLsRareleftrightlimits}]
As noted above, we will only prove~\eqref{eq:rightside} of Proposition~\ref{prop:relatetooneoverderiv}, since the proof of~\eqref{eq:leftside} is symmetrical. Similarly, for the proof of~\eqref{eq:sLsRareleftrightlimits}, we will only prove that
\begin{equation}\label{eq:sctsgoal}
\sfs_{\Right}(\theta,t)=\lim_{\theta'\downarrow\theta}\sfs_{\Diamond}(\theta,t),\quad\Diamond\in\LR,
\end{equation}
as the proof of the other limit is again symmetrical. By Proposition~\ref{prop:findleftrightshocks}, we have
\[
u_{\theta,\Left,\Right}(t,\sfs_{\Right}(\theta,t))\ne u_{\theta,\Right,\Right}(t,\sfs_{\Right}(\theta,t)),
\]
which means that~\eqref{eq:canmoveshock} is satisfied and Proposition~\ref{prop:findYs} applies with $x=\sfs_{\Right}(\theta,t)$. Therefore, we have an $\eps>0$ and a function $\sfr\colon[0,\eps)\to\TT$ such that
\begin{equation}\label{eq:r0issR}
\sfr(0)=\sfs_{\Right}(\theta,t),
\end{equation}
\begin{equation}\label{eq:inGS}
(\theta+\zeta,t,\sfr(\zeta))\in\GS\quad\text{for each }\zeta\in[0,\eps),
\end{equation}
and
\begin{align}\label{eq:rprime0}
\sfr'(0+) & =\frac{1}{u_{\theta,\Left,\Right}(t,\sfs_{\Right}(\theta,t))-u_{\theta,\Right,\Right}(t,\sfs_{\Right}(\theta,t))}.
\end{align}
Now Proposition~\ref{prop:shocksdiffdiscrete} and~\eqref{eq:inGS} tell us that, by reducing $\eps$ if necessary, we can assume that $\sfr(\zeta)=\sfs_{\Left}(\theta+\zeta,t)=\sfs_{\Right}(\theta+\zeta,t)$ for all $\zeta\in(0,\eps)$. Combining this with~\eqref{eq:r0issR}, we see that in fact $\sfr(\zeta)=\sfs_{\Right}(\theta+\zeta,t)$ for all $\zeta\in[0,\eps)$. The continuity of $\sfr$ then implies~\eqref{eq:sctsgoal}, and the conclusion~\eqref{eq:rightside} of Proposition~\ref{prop:relatetooneoverderiv} is now simply~\eqref{eq:rprime0}.
\end{proof}

Finally, we can prove Theorem~\ref{thm:mainthm}$\MK$\eqref{enu:derivintheta}.


\begin{proof}[{Proof of Theorem~\ref{thm:mainthm}$\MK$\eqref{enu:derivintheta}}]
We assume that $\Square=\Right$; the proof in the case $\Square=\Left$ is analogous. We have by Proposition~\ref{prop:differenceintermsoftopology} that, for any $\theta_{1}<\theta_{2}$, $t\in\RR$, and $x\in\TT$, we have
\begin{equation}\label{eq:applyuthetadiff}
u_{\theta_{2},\Right}(t,x)-u_{\theta_{1},\Right}(t,x)=\sum_{\substack{\theta\in[\theta_{1},\theta_{2})\\
(\theta,t,x)\in\GS } }\bigl(u_{\theta,\Left,\Right}(t,x)-u_{\theta,\Right,\Right}(t,x)\bigr).
\end{equation}
Now if $(\theta,t,x)\in\GS$, then $x\in\{s_{\Square}(\theta,t)\st\Square\in\LR\}$. But if $x=\sfs_{\Left}(\theta,t)\ne\sfs_{\Right}(\theta,t)$, then
Proposition~\ref{prop:converse} tells us that $u_{\theta,\Left,\Right}(t,x)=u_{\theta,\Right,\Right}(t,x)$, so the contribution to the right side of~\eqref{eq:applyuthetadiff} is zero. Thus, we obtain
\[
u_{\theta_{2},\Right}(t,x)-u_{\theta_{1},\Right}(t,x)=\sum_{\substack{\theta\in[\theta_{1},\theta_{2})\\
\sfs_{\Right}(\theta,t)=x } }\bigl(u_{\theta,\Left,\Right}(t,x)-u_{\theta,\Right,\Right}(t,x)\bigr)=\sum_{\substack{\theta\in[\theta_{1},\theta_{2})\\
\sfs_{\Right}(\theta,t)=x } }\frac{1}{\partial_{\theta}\sfs_{\Right}(\theta+,t)},
\]
with the last identity by Proposition~\ref{prop:relatetooneoverderiv}, and thus we obtain~\eqref{eq:utheta2theta1}.
\end{proof}

\section{Verifying the hypotheses for compound Poisson forcing}
\label{sec:poisson}

In this section we show that compound Poisson processes are in $\widetilde{\Omega}$ (defined in Definition~\ref{def:Omegatilde2}) with probability $1$, proving Theorem~\ref{thm:poisson-ok}. In this section, we let $\bbP$ be a probability measure under which $\mu$ is the measure associated to a nonnegative homogeneous compound Poisson process on $ \RR\times\TT$. Theorem~\ref{thm:poisson-ok} is a combination of the three propositions in this section.

First we address $\widetilde{\Omega}_{1}$, defined in Definition~\ref{def:smallnoisezones}.
\begin{prop}\label{prop:POmega11}
We have $\bbP(\widetilde{\Omega}_{1})=1$.
\end{prop}

\begin{proof}
Let $\bbU$ be the weight distribution of $\bbP$, so $\bbU$ is a probability measure on $(0,\infty)$. This means that
\begin{multline}\label{eq:Udef}
\bbU(B)=\bbP\bigl(\mu(A)\in B\mid\#(\ssN\cap A)=1\bigr)\\
\text{for any Borel }A\subseteq\RR\times\TT\text{ and }B\subseteq(0,\infty).
\end{multline}
Choose $M>0$ large enough that $\bbU\left(\left(\frac{1}{4M},\infty\right)\right)>0$. Defining $A_{r,N}\coloneqq[r-N,r+N]\times\TT$, we see that
\begin{multline*}
\bbP \left(\#(\ssN\cap A_{r,2M})=\#(\ssN\cap A_{r,M})=1\quad\text{and}\quad\mu(\ssN\cap A_{r,M})>\frac{1}{4M}\right)\\
\begin{aligned}
& =\bbP\left(\#(\ssN\cap(A_{r,2M}\setminus A_{r,M}))=0\right)\cdot\bbP\left(\#(\ssN\cap A_{r,M})=1\right)\\
&\qquad \cdot\bbP\left(\mu(\ssN\cap A_{r,M})>\frac{1}{4M}\ \middle|\ \#(\ssN\cap A_{r,M})=1\right)\\
& =\bbP\left(\#(\ssN\cap(A_{r,2M}\setminus A_{r,M}))=0\right)\cdot\bbP\left(\#(\ssN\cap A_{r,M})=1\right)\cdot\bbU\left(\left(\frac{1}{4M},\infty\right)\right)>0.
\end{aligned}
\end{multline*}
By the independence and spatial homogeneity of the Poisson process, this means that there are constants $\rho<1$ and $C<\infty$, independent of $t$, such that for any $k\ge3$, we have
\begin{align*}
\bbP & \left(
\begin{multlined}
\exists r\in[t-kM,t-2M]\text{ s.t. }\#(\ssN\cap A_{r,2M})\\=\#(\ssN\cap A_{r,M})=1\quad\text{and}\quad\mu(\ssN\cap A_{r,M})>\frac{1}{4M}
\end{multlined}\right)
\ge1-C\rho^{k}.
\end{align*}
The fact that $\bbP(\widetilde{\Omega}_{1})=1$ then follows from the Borel--Cantelli lemma.
\end{proof}

Now we address $\widetilde{\Omega}_{2}$, defined in Definition~\ref{def:Omegatilde2}. We note that $\widetilde{\Omega}_{2}=\widetilde{\Omega}_{2;1}\cap\widetilde{\Omega}_{2;2}$, with $\widetilde{\Omega}_{2;1}$ and $\widetilde{\Omega}_{2;2}$ defined in Proposition~\ref{prop:omega21} and Proposition~\ref{prop:omega22} below, respectively.

\begin{prop}\label{prop:omega21}
Let $\widetilde{\Omega}_{2;1}$ be the set of all $\mu\in\widetilde{\Omega}$ such that $\#(\GS_{\theta}\cap\ssN)\le1$ for all $\theta\in\RR$. We have $\bbP(\widetilde{\Omega}_{2;1})=1$.
\end{prop}

\begin{proof}
If $\mu\in\widetilde{\Omega}_{1}$ and $t\in\RR$, then we define
\[
\clR_{t}\coloneqq\{\theta\in\RR\st\GS_{\theta}\cap\ssN\cap((-\infty,t]\times\TT)\ne\varnothing\}=\bigcup_{(s,x)\in\ssN\cap((-\infty,t]\times\TT)}\GS'_{s,x},
\]
with $\GS'_{s,x}$ defined in \eqref{eq:GS'def}. By Proposition~\ref{prop:GS'discrete} and the countability of $\ssN$, we see that $\clR_{t}$ is countable for each $t$, being the countable union of countable sets. Also, for $t\le t'$, we define
\[
\clE_{t,t'}\coloneqq\bigl\{ \theta\in\RR\st\exists(s,y)\in\ssS_{\theta}\text{ s.t. }[(t,t')\times\TT]\cap\ssN=\{(s,y)\}\bigr\}.
\]
We note that, if $\mu\not\in\widetilde{\Omega}_{2;1}$, then there must be some $t,t'\in\bbQ$ such that $\clR_{t}\cap\clE_{t,t'}\ne\varnothing$, which means that
\begin{equation}\label{eq:ctblunion}
\bbP\left(\widetilde{\Omega}_{1}\setminus\widetilde{\Omega}_{2;1}\right)\le\sum_{t,t'\in\bbQ}\bbP\left(\mu\in\widetilde{\Omega}_{1}\quad\text{and}\quad\clR_{t}\cap\clE_{t,t'}\ne\varnothing\right).
\end{equation}

Let $\clF_{t,t'}$ be the $\sigma$-algebra generated by $\mu|_{\RR\setminus(t,t')}$. From Definition~\ref{def:minimizers}, it is not difficult to check that $\widetilde{\Omega}_{1}$ is measurable with respect to $\clF_{t,t'}$ for any $-\infty<t\le t'<+\infty$. (That is, it is a tail event in an appropriate sense, although we will not need any zero-one theorems here.) It is also not difficult to see from Proposition~\ref{prop:basicproperties}$\MK$\eqref{enu:straightlines} that, for any fixed $\theta\in\RR$ and $t\le t'$, we have, $\bfun_{\widetilde{\Omega}_{1}(\mu)}\bbP(\theta\in\clE_{t,t'}\mid\clF_{t,t'})=0$. Thus we can compute
\begin{multline*}
\bbP\left(\mu\in\widetilde{\Omega}_{1}\text{ and }\clR_{t}\cap\clE_{t,t'}\ne\varnothing\right) \le\bbE\left[\#(\clR_{t}\cap\clE_{t,t'});\mu\in\widetilde{\Omega}_{1}\right ]\\
=\bbE\left[\bbE[\#(\clR_{t}\cap\clE_{t,t'})\mid\clF_{t,t'}];\mu\in\widetilde{\Omega}_{1}\right] =\bbE\left[\sum_{\theta\in\clR_{t}}\bbP\left(\theta\in\clE_{t,t'}\,\middle|\,\clF_{t,t'}\right);\mu\in\widetilde{\Omega}_{1}\right]=0.
\end{multline*}
Using this in~\eqref{eq:ctblunion}, we see that $\bbP(\widetilde{\Omega}_{1}\setminus\widetilde{\Omega}_{2;1})=0$ and hence that $\bbP(\widetilde{\Omega}_{2;1})=1$ by Proposition~\ref{prop:POmega11}.
\end{proof}

\begin{prop}\label{prop:omega22}
Let $\widetilde{\Omega}_{2;2}$ be the set of all $\mu\in\widetilde{\Omega}$ such that $\ssS_{\theta}\cap\ssN\setminus\GS_{\theta}=\varnothing$ for all $\theta\in\RR$. We have $\bbP(\widetilde{\Omega}_{2;2})=1$.
\end{prop}

\begin{proof}
Suppose that $\mu\in\widetilde{\Omega}_{1}$. If $(t,x)\in\ssS_{\theta}\setminus\GS_{\theta}$, then we must have $X_{\theta,t,x,\Left}'(t-)\ne X_{\theta,t,x,\Right}'(t-)$ but
\begin{equation}\label{eq:notglobal}
\int_{T_{\vee}(\theta,t,x)}^{t}X_{\theta,t,x,\Left}'(s)\,\dif s=\int_{T_{\vee}(\theta,t,x)}^{t}X_{\theta,t,x,\Right}'(s)\,\dif s,
\end{equation}
as well as
\begin{align*}
0 & =\clA_{\theta,T_{\vee}(\theta,t,x),t}[X_{\theta,t,x,\Right}]-\clA_{\theta,T_{\vee}(\theta,t,x),t}[X_{\theta,t,x,\Left}]\\
\oversetx{\eqref{eq:Adef}} & =\frac{1}{2}\int_{T_{\vee}(\theta,t,x)}^{t}\left(X_{\theta,t,x,\Right}'(s)^{2}-X_{\theta,t,x,\Left}'(s)^{2}\right)\,\dif s
-\!\theta\!\int_{T_{\vee}(\theta,t,x)}^{t}\left(X_{\theta,t,x,\Right}'(s)-X_{\theta,t,x,\Left}'(s)\right)\,\dif s\\
&\quad
-\!\mu\bigl(\{(s,X_{\theta,t,x,\Right}(s))\st s\in(T_{\vee}(\theta,t,x),t)\}\bigr)\\
&\qquad+\mu\bigl(\{(s,X_{\theta,t,x,\Left}(s))\st s\in(T_{\vee}(\theta,t,x),t)\}\bigr)\\
\oversetx{\eqref{eq:notglobal}} & =\left[\frac{1}{2}\int_{T_{\vee}(\theta,t,x)}^{t}X_{\theta,t,x,\Square}'(s)^{2}\,\dif s-\mu\bigl(
\bigl\{(s,X_{\theta,t,x,\Right}(s))\st s\in(T_{\vee}(\theta,t,x),t)\bigr\}\bigr)\right]_{\Square=\Left}^{\Square=\Right}.
\end{align*}
Fix some arbitrary $y\in\TT$ such that $(t_{1},y)\not\in\ssN$. If $\mu\not\in\widetilde{\Omega}_{2;2}$, then there must exist $t_{1},t_{2}\in\bbQ$ with $t_{1}<t_{2}$, paths $Y_{1},Y_{2}$ connecting $(y,t_{1})$ and elements of $\ssN$ by straight line segments (of which there are at most countably many), and $\tau,\eta\in\RR$ such that
\begin{multline}\label{eq:Ydoesnotdepend}
\frac{1}{2}\int_{T_{*}(t)}^{t}\left(\clT_{\tau,\eta}^{1}Y_{1}\right)'(s)^{2}\,\dif s-\mu\left(\left\{(s,\clT_{\tau,\eta}^{1}Y_{\Square}(s))\st s\in(T_{*}(t),t)\right\}\right)
\\
\text{does not depend on }i\in\{1,2\}.
\end{multline}
Now given such $t_{1},t_{2},Y_{1},Y_{2}$, the set $\clJ(t_{1},t_{2},Y_{1},Y_{2})$ of $(t,x)\in[t_{1},t_{2}]\times\TT$ such that~\eqref{eq:Ydoesnotdepend} can hold is at most countable. This means that
\[
\bbP\left(\widetilde{\Omega}_{1}\setminus\widetilde{\Omega}_{2;2}\right)\le\sum_{\substack{t_{1},t_{2}\in\bbQ\\
Y_{1},Y_{2} } }\; \sum_{(t,x)\in\clJ(t_{1},t_{2},Y_{1},Y_{2})}\bbP((t,x)\in\ssN)=0,
\]
and hence that $\bbP(\widetilde{\Omega}_{2;2})=1$ by Proposition~\ref{prop:POmega11}.
\end{proof}

\let\OLDthebibliography\thebibliography
\renewcommand\thebibliography[1]{
\OLDthebibliography{#1}
\setlength{\parskip}{0pt}
\setlength{\itemsep}{7pt plus 0.5ex}}

\bibliography{dunlap}
\end{document}
